% Options for packages loaded elsewhere
\PassOptionsToPackage{unicode}{hyperref}
\PassOptionsToPackage{hyphens}{url}
\PassOptionsToPackage{dvipsnames,svgnames,x11names}{xcolor}
%
\documentclass[
  10pt,
]{article}
\usepackage{amsmath,amssymb}
\usepackage{iftex}
\ifPDFTeX
  \usepackage[T1]{fontenc}
  \usepackage[utf8]{inputenc}
  \usepackage{textcomp} % provide euro and other symbols
\else % if luatex or xetex
  \usepackage{fontspec}
  \defaultfontfeatures{Scale=MatchLowercase}
  \defaultfontfeatures[\rmfamily]{Ligatures=TeX,Scale=1}
\fi
\usepackage{lmodern}
\ifPDFTeX\else
  % xetex/luatex font selection
    \setmainfont[]{STIX Two Text}
\fi
% Use upquote if available, for straight quotes in verbatim environments
\IfFileExists{upquote.sty}{\usepackage{upquote}}{}
\IfFileExists{microtype.sty}{% use microtype if available
  \usepackage[]{microtype}
  \UseMicrotypeSet[protrusion]{basicmath} % disable protrusion for tt fonts
}{}
\makeatletter
\@ifundefined{KOMAClassName}{% if non-KOMA class
  \IfFileExists{parskip.sty}{%
    \usepackage{parskip}
  }{% else
    \setlength{\parindent}{0pt}
    \setlength{\parskip}{6pt plus 2pt minus 1pt}}
}{% if KOMA class
  \KOMAoptions{parskip=half}}
\makeatother
\usepackage{xcolor}
\usepackage[margin=1in]{geometry}
\usepackage{longtable,booktabs,array}
\usepackage{calc} % for calculating minipage widths
% Correct order of tables after \paragraph or \subparagraph
\usepackage{etoolbox}
\makeatletter
\patchcmd\longtable{\par}{\if@noskipsec\mbox{}\fi\par}{}{}
\makeatother
% Allow footnotes in longtable head/foot
\IfFileExists{footnotehyper.sty}{\usepackage{footnotehyper}}{\usepackage{footnote}}
\makesavenoteenv{longtable}
\usepackage{graphicx}
\makeatletter
\newsavebox\pandoc@box
\newcommand*\pandocbounded[1]{% scales image to fit in text height/width
  \sbox\pandoc@box{#1}%
  \Gscale@div\@tempa{\textheight}{\dimexpr\ht\pandoc@box+\dp\pandoc@box\relax}%
  \Gscale@div\@tempb{\linewidth}{\wd\pandoc@box}%
  \ifdim\@tempb\p@<\@tempa\p@\let\@tempa\@tempb\fi% select the smaller of both
  \ifdim\@tempa\p@<\p@\scalebox{\@tempa}{\usebox\pandoc@box}%
  \else\usebox{\pandoc@box}%
  \fi%
}
% Set default figure placement to htbp
\def\fps@figure{htbp}
\makeatother
\setlength{\emergencystretch}{3em} % prevent overfull lines
\providecommand{\tightlist}{%
  \setlength{\itemsep}{0pt}\setlength{\parskip}{0pt}}
\setcounter{secnumdepth}{-\maxdimen} % remove section numbering
% Map the Unicode math operators used in prose to LaTeX math commands,
% so they render even though the text font lacks the glyphs.
\usepackage{amssymb,amsmath}
\usepackage{newunicodechar}
\newunicodechar{→}{\ensuremath{\rightarrow}}
\newunicodechar{∂}{\ensuremath{\partial}}
\newunicodechar{∈}{\ensuremath{\in}}
\newunicodechar{∎}{\ensuremath{\blacksquare}}
\newunicodechar{□}{\ensuremath{\Box}}
\newunicodechar{∝}{\ensuremath{\propto}}
\newunicodechar{∞}{\ensuremath{\infty}}
\newunicodechar{∫}{\ensuremath{\int}}
\newunicodechar{≈}{\ensuremath{\approx}}
\newunicodechar{≠}{\ensuremath{\neq}}
\newunicodechar{≡}{\ensuremath{\equiv}}
\newunicodechar{≤}{\ensuremath{\leq}}
\newunicodechar{≥}{\ensuremath{\geq}}
\newunicodechar{≪}{\ensuremath{\ll}}
\newunicodechar{≫}{\ensuremath{\gg}}
\newunicodechar{≲}{\ensuremath{\lesssim}}
\newunicodechar{⊙}{\ensuremath{\odot}}
\newunicodechar{⊗}{\ensuremath{\otimes}}
\newunicodechar{⟹}{\ensuremath{\Longrightarrow}}
\newunicodechar{⟺}{\ensuremath{\Longleftrightarrow}}
\newunicodechar{√}{\ensuremath{\surd}}
\newunicodechar{∼}{\ensuremath{\sim}}
\newunicodechar{∇}{\ensuremath{\nabla}}
\newunicodechar{↔}{\ensuremath{\leftrightarrow}}
\newunicodechar{⟨}{\ensuremath{\langle}}
\newunicodechar{⟩}{\ensuremath{\rangle}}
% keep figures from floating far from their callouts
\usepackage{float}
\let\origfigure\figure
\let\endorigfigure\endfigure
\usepackage{bookmark}
\IfFileExists{xurl.sty}{\usepackage{xurl}}{} % add URL line breaks if available
\urlstyle{same}
\hypersetup{
  pdftitle={Structural Entropy Dark Energy: a fixed-parameter, growth-gated holographic dark-energy model without a cosmological constant},
  pdfauthor={Stilian Pandev},
  colorlinks=true,
  linkcolor={RoyalBlue},
  filecolor={Maroon},
  citecolor={Blue},
  urlcolor={RoyalBlue},
  pdfcreator={LaTeX via pandoc}}

\title{Structural Entropy Dark Energy: a fixed-parameter, growth-gated
holographic dark-energy model without a cosmological constant}
\author{Stilian Pandev}
\date{Draft --- 29 June 2026}

\begin{document}
\maketitle
\begin{abstract}
We present Structural Entropy Dark Energy (SEDE), a model without a
cosmological constant in which dark energy is identified with the
thermodynamic conjugate of the cosmic apparent-horizon entropy,
\(\rho _{\mathrm{DE}} = T_{\mathrm{AH}}\) \(s_{\mathrm{grav}}\)
\(f_{\mathrm{sat}}(z)\), gated by a structure-growth factor
\(f_{\mathrm{sat}}\) whose redshift dependence follows a
halo-binding-entropy prescription. A single Friedmann fixed-point
equation replaces Λ by a structure-gated, horizon-coupled term that adds
\textbf{no continuously-sampled dark-energy parameter} relative to ΛCDM:
the dark sector is fixed instead by a few stated, discrete modelling
choices --- the H-coupling λ = 1/2 from a volume-law horizon entropy,
the equation of state from spatial flatness, the coupling γ ≈ 1.5 from a
halo-binding prescription, the sound speed \(c_{\mathrm{s}}^{2} = 1\)
--- one of which (the entropy weight p) was revised once from a halo to
a horizon reservoir. In a marginalised, CAMB-in-the-loop joint analysis
(DESI DR2 BAO, Pantheon+/DES-5YR/Union3 supernovae, cosmic chronometers,
fσ8, and compressed Planck data), SEDE is \textbf{mildly preferred} over
ΛCDM (ΔDIC ≈ −2.9). The preference survives folding the full primary CMB
likelihood into the joint (Δχ² = −3.17) and the full non-lite Planck
plik with all foregrounds sampled (Δχ² = −0.33) --- so it is not an
artifact of CMB compression; a nested-sampling Bayesian evidence gives
ln B = +1.89 (positive-but-weak, favouring SEDE). A per-probe
decomposition shows this preference is carried by the CMB-distance and
H₀ (SH0ES) channels, \textbf{not} by the DESI BAO data --- which SEDE
fits comparably to ΛCDM --- so ``mildly preferred'' should not be read
as a DESI-BAO preference. Dropping the contested SH0ES prior entirely,
the preference weakens but survives and shifts onto the harder channel
(Δχ² = −1.75, ln B = +0.95, \textasciitilde1.9σ; the CMB-distance
carries it with H₀ no longer pulled high). A false-preference
calibration indicates the preference is \textbf{not obviously
attributable to model flexibility} (\textasciitilde2--3σ with SH0ES,
\textasciitilde1.9σ without, depending on the CMB-distance treatment),
and a pre-DESI → DESI split shows a SEDE realisation trained on pre-DESI
data predicts the DESI DR2 BAO trend better than ΛCDM. The equation of
state crosses −1 with \(w_{\mathrm{a}} < 0\), in DESI's
evolving-dark-energy quadrant, and the vacuum-energy-scale problem is
recast via the Cohen--Kaplan--Nelson bound rather than inherited. The
construction rests on \textbf{one explicit quantum-gravitational
postulate} --- a volume-law (thermalised) horizon entropy --- testable
by a measurement of the deformation Δ; under the stated Fisher
assumptions, DESI DR3 and Euclid reach σ(Δ) ≈ 0.09, sharply
distinguishing the volume-law and area-law branches. SEDE is
\textbf{suggestive but not established}; the next generation of surveys
will decide.
\end{abstract}

{
\hypersetup{linkcolor=}
\setcounter{tocdepth}{2}
\tableofcontents
}
\section{1. Introduction}\label{introduction}

The standard cosmological model, ΛCDM, fits an extraordinary range of
data with a cosmological constant Λ whose value confronts two
long-standing puzzles. The \textbf{cosmological-constant problem}: if Λ
is the energy of the quantum vacuum, naïve effective field theory
overshoots the observed value by \textasciitilde10¹²⁰. The
\textbf{coincidence problem}: the dark-energy and matter densities are
comparable today, \(\rho _{\mathrm{DE}}\) \textasciitilde{}
\(\rho _{\mathrm{m}}\), despite scaling differently with expansion ---
so we appear to live at a special epoch. Recent DESI DR2 results,
moreover, give a mild preference for \emph{evolving} dark energy (w
crossing −1) over a constant Λ, reopening the question of whether Λ is
fundamental at all.

We approach this through an accounting question. As cosmic structure
forms, collapsing halos release gravitational binding energy, and dark
energy activates over the \emph{same} epoch that structure formation
approaches saturation. Is that a coincidence, or is the dark sector tied
to structure formation by an energy ledger? It is tempting to guess that
the binding energy released by structure \emph{is} the dark energy ---
but a one-line budget settles it: a virialised system has
\textbar{}\(E_{\mathrm{bind}}\)\textbar/Mc² \textasciitilde{} σ²/c², and
even clusters have σ²/c² \textasciitilde{} 10⁻⁵ (galaxies
\textasciitilde10⁻⁷), so the binding-energy density is
\textasciitilde10⁶ times too small to \emph{be} the observed
\(\rho _{\mathrm{DE}}\). The link, if there is one, cannot be an energy
supply.

The resolution we propose is thermodynamic. We identify dark energy with
the \textbf{conjugate thermal energy density of the cosmic horizon} ---
\(\rho _{\mathrm{DE}} = T_{\mathrm{AH}}\cdot s_{\mathrm{grav}}\), the
conjugate energy of the apparent-horizon entropy (motivated by the
horizon thermodynamics of Jacobson 1995; Cai--Kim 2005) --- whose
\emph{scale} is the horizon's own, \(T_{\mathrm{AH}}\) s
\textasciitilde{} \(M_{\mathrm{P}}^{2}H^{2}\) \textasciitilde{}
\(\rho _{\mathrm{crit}}\), rather than the QFT vacuum (so the
vacuum-energy-scale problem is recast, not inherited) and not structure.
What structure supplies is not energy but \textbf{entropy bookkeeping}:
by the horizon first law δQ = T dS, the binding-entropy ledger
associated with halo virialisation fixes the \emph{redshift dependence}
of the effective activation fraction
\(f_{\mathrm{sat}}(z) \in  [0\),1{]}, while its amplitude and the
\(f_{\mathrm{sat}}(0) = 1\) normalisation are fixed separately by the
horizon-scale normalisation. Energy is conserved throughout --- total
covariant conservation is automatic, and the dark-energy fluid is
self-conserved --- with the binding energy playing a purely entropic,
shape-setting role (§2). The coincidence is then a consequence of the
mechanism, not a tuning: dark energy turns on with a gate whose redshift
dependence tracks structure formation as it approaches saturation,
i.e.~near ``now.''

The idea that dark energy is tied to the entropy/information produced as
structure forms is not new --- it originates with Gough's information
dark energy {[}Gough 2008, 2011, 2025{]} --- and we credit that priority
explicitly. What is specific to SEDE is the \emph{mechanism} (a
gravitational horizon-entropy realisation,
\(\rho _{\mathrm{DE}} = T_{\mathrm{AH}}\cdot s_{\mathrm{grav}}\), driven
by the gravitational growth factor rather than the baryonic
star-formation history) and a parameter-free, first-principles dark
sector; the relationship is quantified in §7.6.

The model is, to a useful approximation, \textbf{ΛCDM with a
cosmological constant that fades in as structure forms}: schematically
\(\Lambda _{\mathrm{eff}}(z) \propto  f_{\mathrm{sat}}(z)\) (the exact
running term is \(3\Omega _{\mathrm{DE0}}\) H H₀ \(f_{\mathrm{sat}}\),
§2), with \(f_{\mathrm{sat}}\) running from ≈0 in the early universe to
1 today. By \textbf{Λ-free} we mean specifically that the late-time
acceleration is \emph{not} a constant stress-energy term (p = −ρ =
const.) but a self-conserved horizon fluid whose present amplitude is
fixed by flatness
\((\Omega _{\mathrm{DE0}} = 1 - \Omega _{\mathrm{m}} - \Omega _{\mathrm{r}})\)
--- exactly as flat ΛCDM fixes \(\Omega _\Lambda 0\) --- and whose
redshift dependence is set by the horizon/growth prescription rather
than held constant; it does \emph{not} mean the absence of a dark-energy
normalisation. Its distinctive feature is that the dark sector adds
\textbf{no fitted parameter} --- every input is fixed by a stated
theoretical choice (§3), so SEDE has the same cosmological parameter
count as the ΛCDM baseline used in this analysis (the five of the
compressed late-time pipeline of §4). The price is a single postulate
about the quantum-gravitational structure of the horizon (§8), which we
make explicit, plus the accompanying fixed choices (the halo-entropy
prescription for γ, the smooth-DE closure \(c_{\mathrm{s}}^{2}=1\)),
which we flag as such rather than as derivations.

This paper is organised around the energy ledger. §2 states the model
and establishes energy conservation as its backbone, fixing the
(entropic, not energetic) role of binding energy; §3 fixes the
dark-sector inputs --- the structure coupling γ from a halo-entropy
prescription, then w(z), λ, and the sound speed \(c_{\mathrm{s}}^{2}\)
--- separating the horizon-set \emph{magnitude} from the structure-set
\emph{shape}. §4 describes the data and the CAMB-in-the-loop
methodology; §5 presents the results, including a false-preference
calibration and out-of-sample tests; §6 gives the falsifiable
predictions; §7 the relation to dark matter and other sectors. §8 then
supplies the quantum-gravity underpinning --- the magnitude scale and
the maximal deformation --- and isolates the one irreducible postulate;
§9 discusses the open problems and §10 concludes. Throughout we maintain
that SEDE is a \emph{falsifiable proposal}, preferred at a modest level,
not an established result.

\begin{center}\rule{0.5\linewidth}{0.5pt}\end{center}

\section{2. The model and energy
conservation}\label{the-model-and-energy-conservation}

SEDE is general relativity with a single extra fluid: a horizon-fluid
dark-energy density identified with the conjugate thermal energy density
of the cosmic apparent horizon. The whole construction is governed by
one bookkeeping rule --- that energy is conserved when this fluid's
redshift dependence is gated by structure growth. We first state the
model (the conjugate horizon-fluid ansatz, the structure gate, the
Friedmann equation, and how Λ is replaced), and then make energy
conservation the backbone: it is what fixes the role of structure's
gravitational binding energy as \emph{entropic}, not energetic, and what
shows the w = −1 crossing to be sourcing rather than a phantom field.

\textbf{The conjugate horizon-fluid ansatz.} At the cosmological
apparent horizon \(R_{\mathrm{AH}} = 1/H\) (flat FRW), the dynamical
Cai--Kim temperature is
\(T_{\mathrm{AH}} = (H/2\pi )(1 - \epsilon /2)\), ε ≡ −Ḣ/H². For the
dark-energy sector, which tends to a de Sitter attractor (ε → 0), we
adopt the \textbf{Gibbons--Hawking temperature
\(T_{\mathrm{AH}} = H/2\pi\)} --- the ε → 0 limit, i.e.~the appropriate
\emph{leading} temperature for a slowly-evolving horizon fluid. This is
an approximation, made deliberately: the dynamical (1 − ε/2) correction
is largest in the matter era, exactly where \(f_{\mathrm{sat}}\) → 0 and
dark energy is negligible, and including it in full on the volume-law
fluid is in fact \emph{disfavoured} --- it yields an unstable,
phantom-everywhere background with no w = −1 crossing (§6,
\texttt{run\_dynT\_crossing.py}). So \(T_{\mathrm{AH}} = H/2\pi\) is the
physically-appropriate choice for this sector and it is what enters the
calibrated E(z). We \emph{identify} the dark-energy density with the
conjugate horizon-fluid energy density,
\[\rho_{\rm DE} = T_{AH}\, s_{\rm grav}\, f_{\rm sat}(z),\] where
\(s_{\mathrm{grav}}\) is the gravitational entropy \emph{density} and
\(f_{\mathrm{sat}}\) is an effective horizon-entropy activation
\textbf{gate}, normalised to \(f_{\mathrm{sat}}(0) = 1\) --- i.e.~a
value \emph{relative to today}, not a literal fraction bounded by 1 (it
reaches ≈ 1.23 in the future, below); its redshift dependence is set by
the growth-weighted halo prescription of §3.1, not derived from the
Bousso bound itself. \emph{Range and future behaviour (defining the
convention explicitly):} for z ≥ 0 --- the regime of every data fit here
--- the linear growth obeys 0 ≤ D(z) ≤ 1, so
\(0 \leq  f_{\mathrm{sat}} \leq  1\). \(f_{\mathrm{sat}}\) is
\textbf{not clipped}; in the future (z \textless{} 0) linear growth
\emph{saturates} (D freezes at \(D_\infty  \approx  1.44\) as dark
energy suppresses growth --- it does not run away), so
\(f_{\mathrm{sat}}\) saturates to a \textbf{finite} de Sitter value
\(f_{\mathrm{sat}}(D_\infty ) \approx  1.23\) (modestly above 1, never
reaching the formal x → ∞ ceiling \(1/(1-e^{-\gamma }) \approx  1.29\)).
The background tends to a de Sitter attractor (H → \(H_\infty\),
\(\rho _{\mathrm{DE}}\) → const, w → −1; A.7, Result 12) --- bounded and
well-defined for all z, with the ≤ 1 statement specific to z ≥ 0. (The
gate is normalised to its present value; it is \emph{not} the literal
fraction of horizon entropy deposited by collapsed structures, which is
far smaller --- ∼10⁻⁷; App. F.) Clausius horizon thermodynamics
motivates the relation and supplies the conservation logic, but the
absolute identification is the defining SEDE \textbf{ansatz}, not an
identity (App. A.1). We take \(s_{\mathrm{grav}}\) to be the
coarse-grained \textbf{entanglement entropy in the apparent-horizon
volume},
\(S_{\mathrm{grav}} \propto  V_{\mathrm{AH}} \propto  R_{\mathrm{AH}}^{3}\)
--- equivalently a \emph{constant} entropy density
\(s_{\mathrm{grav}} = s_{0}\) (in contrast to the Bekenstein--Hawking
area-law density \(s_{\mathrm{AH}} \propto  H\)). With
\(T_{\mathrm{AH}} \propto  H\),
\[\rho_{\rm DE} = T_{AH}\,s_0\,f_{\rm sat} \;\propto\; H\,f_{\rm sat}, \qquad
s_0 = \frac{\Omega_{\rm DE0}\,\rho_{\rm crit,0}}{T_{AH,0}}\ \text{(fixed by flatness)},\]
giving \textbf{\(\rho _{\mathrm{DE}} \propto  H\) \(f_{\mathrm{sat}}\)}
with no free parameter. This volume law is exactly the maximal Barrow
deformation \((S \propto  A^{1+\Delta /2} = A^{3/2} = R^{3}\) at Δ=1, λ
= 1 − Δ/2 = 1/2); we keep Δ only as a \emph{test} parameter quantifying
deviations from the volume law (§8, and the data measurement of Δ), not
as a model input.

\textbf{The structure gate.} Structure formation supplies the
\emph{redshift dependence} of the activation gate at a rate set by the
collapsed-halo abundance (Result 3, App. A), yielding a closed form for
the normalised activation fraction,
\[f_{\rm sat}(z) = \frac{1 - e^{-\gamma D^2(z)}}{1 - e^{-\gamma}}, \qquad D(z)\ \text{the linear growth factor},\]
where D(z) is the model's \textbf{self-consistent} growth factor
(computed with SEDE's own E(z), which in turn depends on
\(f_{\mathrm{sat}}\) --- the background is solved as the joint fixed
point; §4.4, App. B). with \(f_{\mathrm{sat}}(0) = 1\) imposed as the
present-day normalisation convention and \(f_{\mathrm{sat}}\) → 0 at
high z. The single shape parameter γ is fixed by a halo-entropy
prescription (§3.1), not fitted.

\textbf{The Friedmann equation.} Inserting \(\rho _{\mathrm{DE}}\) into
the (unmodified) Friedmann constraint gives a single \textbf{fixed-point
equation} for the dimensionless expansion rate E = H/H₀:
\[\boxed{\;E^2(z) = \Omega_m(1+z)^3 + \Omega_r(1+z)^4 + \Omega_{\rm DE0}\, f_{\rm sat}(z)\, E(z)\;}\tag{1}\]
(for λ = 1/2, so \(E^{2\lambda } = E\)). Equation (1) is the entire
background model. It is ΛCDM with the constant Λ replaced by the
structure-gated, horizon-coupled term \(\Omega _{\mathrm{DE0}}\)
\(f_{\mathrm{sat}}(z)\) E(z); at z = 0 with \(f_{\mathrm{sat}} = 1\) it
reproduces the observed
\(\Omega _{\mathrm{DE0}} = 1 - \Omega _{\mathrm{m}}\).

\textbf{Replacement of Λ in the field equations.} SEDE is
\emph{holographic dark energy}, not modified gravity: the Einstein
tensor and Newton's constant are untouched, and the deformation enters
only a smooth \((c_{\mathrm{s}}^{2} = 1\); §3.4) perfect fluid on the
right-hand side,
\[G_{\mu\nu} = 8\pi G\,(T^{\rm (matter)}_{\mu\nu} + T^{\rm (DE)}_{\mu\nu}),\qquad
\rho_{\rm DE} = \tfrac{3}{8\pi G}\,\Omega_{\rm DE0}\, H H_0\, f_{\rm sat},\quad p_{\rm DE} = w(z)\rho_{\rm DE}.\]
Equivalently, Λ becomes a \emph{running} scalar
\(\Lambda _{\mathrm{eff}}(z) = 3\) \(\Omega _{\mathrm{DE0}}\) H H₀
\(f_{\mathrm{sat}}\) that equals the observed Λ today and is negligible
in the early universe. The choice of right-hand side (fluid) over
left-hand side (geometry) is not cosmetic: a
Barrow-\emph{modified-gravity} Friedmann equation would rescale the
early-universe expansion by \((M_{\mathrm{P}}/H)^{\Delta }\)
\textasciitilde{} 10⁶¹, a Big-Bang-nucleosynthesis catastrophe (§8,
Result 11). Keeping the deformation in the dark-energy fluid leaves the
early universe exactly standard.

\textbf{Energy conservation.} Because SEDE is general relativity with a
fluid source, the Bianchi identity \(\nabla^\mu G_{\mu\nu}=0\) forces
the \textbf{total} stress--energy to be covariantly conserved,
\(\nabla^\mu(T^{(\rm matter)}_{\mu\nu} + T^{(\rm DE)}_{\mu\nu})=0\),
identically --- there is no violation. At the background level matter
dilutes standardly, \(\rho _{\mathrm{m}} \propto  a^{-3}\), and dark
energy is a \textbf{self-conserved fluid} whose density evolves entirely
through its own pressure work,
\[\dot\rho_{\rm DE} + 3H\,(1+w_{\rm DE})\,\rho_{\rm DE} = 0, \qquad
w_{\rm DE}(a) = -1 - \tfrac13\,\frac{d\ln\rho_{\rm DE}}{d\ln a}\] (the
w(z) of §3.2 and Fig. 4). There is no literal energy transfer from
matter to dark energy: such a transfer of order
\(\rho _{\mathrm{crit}}\) over a Hubble time would spoil
\(\rho _{\mathrm{m}} \propto  a^{-3}\) and with it BBN and the CMB. This
is the crucial accounting point of the model. The \textbf{magnitude} of
\(\rho _{\mathrm{DE}}\) is anchored to the horizon gravitational scale ∼
\(M_{\mathrm{P}}^{2}\) H² ∼ \(\rho _{\mathrm{crit}}\) through the CKN
bound and the flatness normalisation (§8.2), not to energy borrowed from
structure: the gravitational binding energy released as halos virialise
\((E_{\mathrm{bind}} \propto  M^{5/3})\) is ∼ σ²/c² ∼ 10⁻⁶ of the
collapsed rest-mass energy, six orders of magnitude too small to
\emph{be} the dark energy. This makes the genuinely new content precise:
with the area law (Δ=0) the ratio
\(\rho _{\mathrm{DE}}/\rho _{\mathrm{crit}} \propto  H^{-\Delta }\) is a
fixed fraction that merely tracks --- GR's own horizon energy repackaged
(Jacobson) --- whereas the Barrow deformation Δ\textgreater0 makes it
grow as H falls, turning a tracking term into dark energy that comes to
dominate; the new physics is \emph{entirely} the deformation (``is the
energy new?'' ⟺ ``is Δ≠0?''). What structure formation supplies is
therefore not energy but \textbf{entropy bookkeeping}: by the
apparent-horizon first law (A.1), the deposited binding entropy
\(\Delta S_{\mathrm{AH}} = E_{\mathrm{bind}}/T_{\mathrm{AH}}\) fixes the
\emph{weight} p = 5/3 that sets the shape γ of the gate
\(f_{\mathrm{sat}}\) (§3.1) --- it determines \emph{when} and \emph{how
sharply} dark energy turns on, while the overall scale, and the
normalisation \(f_{\mathrm{sat}}(0) = 1\), come from the horizon (§3,
§8). Read this way, SEDE's crossing of w = −1 is \textbf{not a phantom
field} --- there is no negative kinetic term or gradient instability ---
but the pressure work of a horizon-sourced fluid whose density rises as
\(f_{\mathrm{sat}}\) grows (Result 13: 1 + w is the horizon
entropy-production rate). (A single canonical or k-essence field cannot
cross w = −1 without its kinetic term
\(\rho  + p = 2X\cdot P_{\mathrm{X}}\) passing through zero --- the
Vikman no-go --- so a fundamental-field description would require a
ghostly quintom; SEDE avoids this because its phantom is
\emph{effective}:
\(\rho _{\mathrm{DE}} = T_{\mathrm{AH}}\cdot s_{\mathrm{grav}}\) is
built from covariant horizon scalars, not a fundamental scalar.)
Finally, since \(f_{\mathrm{sat}}\) is a functional of the
\emph{background} growth \((c_{\mathrm{s}}^{2} = 1\), §3.4), the
coupling lives at the background level only and introduces no local
energy non-conservation in the perturbations.

\pandocbounded{\includegraphics[keepaspectratio]{../../output/paper/fig1_mechanism}}
\textbf{Figure 1.} The mechanism. \emph{Left:} the effective activation
fraction (structure-growth gate) \(f_{\mathrm{sat}}(z)\) rises toward 1
today as structure forms, activating the dark-energy gate --- the
effective running term
\(\Lambda _{\mathrm{eff}} = 3\Omega _{\mathrm{DE0}}\) H H₀
\(f_{\mathrm{sat}}\) \emph{fades in} (computed from the model's
self-consistent linear growth factor D(z), §2). \emph{Right:} the cosmic
density inventory; SEDE's \(\rho _{\mathrm{DE}}\) is orders of magnitude
below matter and radiation at high z and falls below ΛCDM's constant
\(\rho _\Lambda\), so the early universe (BBN, recombination) is
standard --- the deformation acts only on the late-time dark-energy
fluid.

\begin{center}\rule{0.5\linewidth}{0.5pt}\end{center}

\section{3. The fixed inputs --- why the dark sector adds no fitted
parameter}\label{the-fixed-inputs-why-the-dark-sector-adds-no-fitted-parameter}

In flat ΛCDM the dark-energy amplitude
\(\Omega _\Lambda 0 = 1 - \Omega _{\mathrm{m}} - \Omega _{\mathrm{r}}\)
is itself fixed by the flatness closure, not an independent fit; SEDE
keeps exactly this --- no independent dark-energy amplitude beyond
flatness --- and differs only in that the \emph{redshift dependence} is
set by the horizon/growth prescription rather than held constant. What
that prescription fixes is the triplet (γ, w(z)-shape, λ) plus the sound
speed \(c_{\mathrm{s}}^{2}\) --- and each is \textbf{fixed by a stated
theoretical choice} (not fitted to cosmological data), so SEDE adds no
free dark-energy parameter. We are explicit about the status of each
choice below (some follow from flatness or the model construction;
others, notably γ and \(c_{\mathrm{s}}^{2}\), are fixed prescriptions).
The choices separate along the two scales of §2. The \textbf{magnitude}
of dark energy is the horizon scale, ∼ \(M_{\mathrm{P}}^{2}H^{2}\) ∼
\(\rho _{\mathrm{crit}}\) (the CKN bound normalised by flatness, §8.2)
--- not a fitted dark-sector parameter. What the choices control is the
\textbf{shape}: \emph{when} and \emph{how sharply} that energy turns on.
We therefore lead with the structure coupling γ --- the shape, set by
the binding energy --- and then give w₀, λ, and \(c_{\mathrm{s}}^{2}\).

\subsection{3.1 The structure coupling γ ≈ 1.5 from a
halo-binding-energy
prescription}\label{the-structure-coupling-ux3b3-1.5-from-a-halo-binding-energy-prescription}

γ sets the sharpness of the structure gate, γ = d ln
\(\Sigma _{\mathrm{S}} / d\) ln σ8² with \(\Sigma _{\mathrm{S}}\) the
entropy-weighted collapsed mass, \(\Sigma _{\mathrm{S}}\)=∫ \(M^{p}\)
(dn/d ln M) d ln M --- the derivative is with respect to the
\textbf{variance} σ8² because the gate's argument is x = D² =
σ8²(z)/σ8²(0) (§3.4, A.2); equivalently γ = ½ d ln
\(\Sigma _{\mathrm{S}}/d\) ln σ8. The weight exponent p is fixed by
\emph{where} the structure-formation entropy is deposited --- and the
reservoir is fixed \emph{before any data enter}, by the model's own §2
identity. Because SEDE identifies
\(\rho _{\mathrm{DE}} = T_{\mathrm{AH}}\cdot s_{\mathrm{grav}}\) (§2),
\(f_{\mathrm{sat}}\) is \textbf{by construction} the
cosmic-\emph{horizon} entropy activation fraction (A.1), so the binding
energy released by virialisation is accounted in the \textbf{horizon}
ledger at its temperature \(T_{\mathrm{AH}} \propto  H\) --- \emph{not}
the halo's internal virial temperature
\(T_{\mathrm{vir}} \propto  M^{2/3}\). With \(T_{\mathrm{AH}}\)
mass-independent at a fixed epoch, the Clausius relation gives
\(\Delta S_{\mathrm{AH}} = E_{\mathrm{bind}}/T_{\mathrm{AH}} \propto  E_{\mathrm{bind}}\),
and the virial binding energy \(E_{\mathrm{bind}} \propto  M^{5/3}\)
then fixes \textbf{p = 5/3}. This reservoir choice is forced by the §2
identity; we flag explicitly that an earlier version used
\(T_{\mathrm{vir}}\) (giving p = 1), which was a \textbf{reservoir error
corrected for consistency with §2}, not a tuning --- though we
acknowledge it is a discrete modelling revision, and that the corrected
value happens to be the one favoured in diagnostic fits.

γ has a transparent closed form. A self-similar reduction gives **γ = (p
− 1)⟨1/α⟩\_Σ** (α ≡ −d ln σ/d ln M the effective spectral slope, ⟨·⟩\_Σ
the entropy-weighted average), which reproduces the full mass-function
integral to \textbf{0.1\%} (A.3) --- so the dark-sector input is
\emph{only} p = 5/3, with ⟨1/α⟩ ordinary halo physics. The
\emph{numerical} γ then follows from a standard, \textbf{pre-specified}
halo model and inherits its systematics: a Sheth--Tormen mass function
over 10¹⁰--10¹⁶ M⊙ with an EH98 transfer gives \textbf{γ = 1.50} (Result
4C), and varying the mass function (Press--Schechter, Tinker08,
Despali16, Watson13), mass range, and transfer function spans \textbf{γ
∈ {[}1.14, 1.60{]}} (median 1.50; \texttt{run\_gamma\_systematics.py}),
dominated by the mass-range low-cutoff and robust (1.48--1.57) to the
mass-function and transfer choice. We therefore claim precisely: the
reservoir --- hence the weight p = 5/3 --- is fixed by the model's own
construction; the numerical γ ≈ 1.5 carries standard halo-model
systematics of order ±0.2, \textbf{none tuned to cosmological data}.
That it lands near the diagnostic-fit value is then a \emph{consistency
check}, not the motivation. (We do \textbf{not} claim a demonstrated
microphysical transport channel depositing the binding heat onto the
horizon; the prescription is the accounting, and the smooth-DE closure
of §3.4 ensures it induces no local dark-energy perturbation.) Binding
energy thus sets the gate \emph{shape}, not the magnitude (§2, §10
ledger).

\subsection{3.2 The equation of state w(z) from the fixed background,
conservation, and
flatness}\label{the-equation-of-state-wz-from-the-fixed-background-conservation-and-flatness}

The model has a single physical equation of state --- that of the
canonical fixed-point background (eq. 1) --- obtained from energy
conservation,
\[w(z) = -1 - \tfrac13\,\frac{d\ln\rho_{\rm DE}}{d\ln a},\] with
\(\rho _{\mathrm{DE}} = \Omega _{\mathrm{DE0}}\)
\(\rho _{\mathrm{crit}}\),0 E \(f_{\mathrm{sat}}\) and the present-day
normalisation fixed by spatial flatness
\((\Omega _{\mathrm{DE0}} = 1 - \Omega _{\mathrm{m}} - \Omega _{\mathrm{r}}\),
\(f_{\mathrm{sat}}(0) = 1\)). Evaluated on the self-consistent
background this gives \textbf{w₀ ≈ −1.0}, \textbf{crossing −1} at low
redshift (z ≈ 0.2) and dipping to ≈ −1.08 by z ≈ 2 --- DESI DR2's
thawing/crossing quadrant. So w(z) is fixed by \(\Omega _{\mathrm{m}}\)
(through the background) with no free EOS parameter; this is the w(z)
used throughout (§5, Fig. 4). Three w₀ values appear in this paper for
\emph{different} objects; to avoid confusion:

\begin{longtable}[]{@{}
  >{\raggedright\arraybackslash}p{(\linewidth - 6\tabcolsep) * \real{0.2500}}
  >{\raggedright\arraybackslash}p{(\linewidth - 6\tabcolsep) * \real{0.2500}}
  >{\raggedright\arraybackslash}p{(\linewidth - 6\tabcolsep) * \real{0.2500}}
  >{\raggedright\arraybackslash}p{(\linewidth - 6\tabcolsep) * \real{0.2500}}@{}}
\toprule\noalign{}
\begin{minipage}[b]{\linewidth}\raggedright
object
\end{minipage} & \begin{minipage}[b]{\linewidth}\raggedright
scaling
\end{minipage} & \begin{minipage}[b]{\linewidth}\raggedright
role
\end{minipage} & \begin{minipage}[b]{\linewidth}\raggedright
w₀
\end{minipage} \\
\midrule\noalign{}
\endhead
\bottomrule\noalign{}
\endlastfoot
\textbf{canonical SEDE} & \(\rho _{\mathrm{DE}} \propto  H\)
\(f_{\mathrm{sat}}\) & the actual model & ≈ −1.0, crosses −1 \\
λ=1 ``temperature-factor'' cousin & --- & closed-form structural
estimate (side note) & ≈ −0.85 \\
area-law branch & \(\rho _{\mathrm{DE}} \propto  H^{2}\)
\(f_{\mathrm{sat}}\) & disfavoured diagnostic alternative (§5.6, §8.1) &
≈ −0.49 \\
\end{longtable}

Only the first row is the model's prediction; the −0.85 is a sanity
check on the \(\Omega _{\mathrm{m}}\)-dependence (closed form
\(w_{0} = (4\Omega _{\mathrm{m}}/3 - 1)/(1 - \Omega _{\mathrm{m}})\))
and the −0.49 is the Δ=0 alternative the data disfavour.

\subsection{3.3 The H-coupling λ = 0.5 from the horizon's volume
entropy}\label{the-h-coupling-ux3bb-0.5-from-the-horizons-volume-entropy}

λ fixes the \emph{magnitude}'s scaling with H, with no free parameter:
the gravitational entropy is taken to be the coarse-grained entanglement
entropy in the apparent-horizon \textbf{volume},
\(S_{\mathrm{grav}} \propto  V_{\mathrm{AH}} \propto  R_{\mathrm{AH}}^{3}\),
i.e.~a \emph{constant} horizon entropy density. The conjugate
horizon-fluid ansatz
\((\rho _{\mathrm{DE}} = T_{\mathrm{AH}}\cdot s_{\mathrm{grav}}\),
\(T_{\mathrm{AH}} \propto  H\)) then gives
\(\rho _{\mathrm{DE}} \propto  H \equiv  H^{2\lambda }\) with \textbf{λ
= 1/2} --- no deformation parameter. Equivalently, in the Barrow
language \(S = (A/A_{0})^{1+\Delta /2}\) this is the maximal deformation
Δ = 1 \((s_{\mathrm{AH}} \propto  H^{1-\Delta }\), λ = 1−Δ/2): the
volume law \emph{is} Δ=1 \((A^{3/2} = R^{3} = V)\). We keep Δ only as a
\textbf{falsifiable test} of the volume law (its data measurement,
§5/§6), not as a model input; §8 motivates the volume/space-filling
state (capping Δ≤1 and postulating saturation). (Verified Δ-free ≡ Δ=1
bit-for-bit; \texttt{friedmann.E\_SEDE\_volume},
\texttt{run\_volume\_equiv.py}.)

\subsection{\texorpdfstring{3.4 The sound speed
\(c_{\mathrm{s}}^{2} = 1\) from a background
functional}{3.4 The sound speed c\_\{\textbackslash mathrm\{s\}\}\^{}\{2\} = 1 from a background functional}}\label{the-sound-speed-c_mathrms2-1-from-a-background-functional}

SEDE's dark energy is \emph{background-growth-gated} (its
\(f_{\mathrm{sat}}\) depends on the growth history), which might suggest
it clusters with structure. We adopt the \textbf{minimal
smooth-dark-energy closure}: \(f_{\mathrm{sat}}\) is a functional of the
\textbf{background} growth factor D(z) --- a function of cosmic time
only --- so the structure coupling is \textbf{background-level /
nonlocal / coarse-grained, not local}. Hence \(\rho _{\mathrm{DE}}\) is
homogeneous by construction, ∂\(f_{\mathrm{sat}}/\partial (local\)
\(\delta _{\mathrm{m}}) = 0\), and the dark energy has the standard
metric-driven response with \textbf{\(c_{\mathrm{s}}^{2} = 1\)}. We are
explicit about the wavelength scope: this is a statement about the
\textbf{sub-horizon} response --- there is no small-scale dark-energy
clustering --- and it is what the \(c_{\mathrm{s}}^{2} = 1\) validation
below tests. A distinct \emph{long-wavelength, separate-universe}
modulation of the background gate by a super-survey mode is a separate,
super-horizon effect (the optional near-critical layer, deferred to a
companion paper, §6); it is not sub-horizon clustering and is not probed
by the validation here. Operationally SEDE is implemented as an
\textbf{effective smooth dark-energy fluid using a PPF crossing
prescription} (which handles the w = −1 crossing; §2's Vikman-no-go
argument is why a fundamental single perfect fluid would \emph{not}
cross −1 without extra structure, and why SEDE's phantom is
\emph{effective}). This is a modelling closure, not a theorem; a full
CLASS Boltzmann treatment confirms the resulting fσ8(z) matches this
smooth-dark-energy growth to 0.1\% (§4).

\begin{longtable}[]{@{}
  >{\raggedright\arraybackslash}p{(\linewidth - 4\tabcolsep) * \real{0.3333}}
  >{\raggedright\arraybackslash}p{(\linewidth - 4\tabcolsep) * \real{0.3333}}
  >{\raggedright\arraybackslash}p{(\linewidth - 4\tabcolsep) * \real{0.3333}}@{}}
\toprule\noalign{}
\begin{minipage}[b]{\linewidth}\raggedright
input
\end{minipage} & \begin{minipage}[b]{\linewidth}\raggedright
value
\end{minipage} & \begin{minipage}[b]{\linewidth}\raggedright
status --- fixed by
\end{minipage} \\
\midrule\noalign{}
\endhead
\bottomrule\noalign{}
\endlastfoot
λ (H-coupling) & 0.5 & \textbf{postulate} --- volume-law horizon entropy
\(S_{\mathrm{grav}} \propto  V_{\mathrm{AH}}\) (≡ Barrow Δ=1), §3.3,
§8 \\
w(z) (EOS) & w₀ ≈ −1.0, crosses −1 & \textbf{derived} --- spatial
flatness + eq. 1, §3.2 (no free EOS parameter) \\
γ (structure gate) & ≈ 1.50 (range {[}1.14,1.60{]}) & \textbf{fixed
prescription} --- halo-binding-entropy weighting (fixed, not fitted),
§3.1 \\
\(c_{\mathrm{s}}^{2}\) (sound speed) & 1 & \textbf{closure} --- minimal
smooth-DE treatment (validated numerically), §3.4 \\
\end{longtable}

\pandocbounded{\includegraphics[keepaspectratio]{../../output/paper/fig2_inputs}}
\textbf{Figure 2.} The two fixed dark-sector couplings. \emph{Left:} the
H-coupling λ = 1 − Δ/2 as a function of the Barrow deformation; the
volume-law endpoint Δ = 1 (the postulate of §8) fixes λ = 0.5.
\emph{Right:} the structure coupling γ = d ln \(\Sigma _{\mathrm{S}}/d\)
ln σ8² (variance convention, §3.1) as a function of the entropy-weight
exponent p; the binding-energy-into-the-horizon argument (Result 4C /
Prescription 4C, A.3) fixes p = 5/3, giving γ = 1.50 --- close to the
value favoured in diagnostic fits. Neither is a fitted parameter.

\begin{center}\rule{0.5\linewidth}{0.5pt}\end{center}

\section{4. Data \& methodology}\label{data-methodology}

\subsection{4.1 Data}\label{data}

We use a standard late-time-plus-compressed-CMB compilation, with both
models fit to the identical data vector: - \textbf{BAO:} DESI DR2
(galaxy, QSO, and Lyα tracers, z = 0.3--2.33), as
\(D_{\mathrm{M}}/r_{\mathrm{d}}\) and \(D_{\mathrm{H}}/r_{\mathrm{d}}\).
- \textbf{Supernovae:} Pantheon+ as the baseline, with DES-5YR and
Union3 used independently for the robustness cross-check of §5.4. -
\textbf{Cosmic chronometers:} the Moresco et al.~H(z) compilation
(model-independent expansion rates). - \textbf{Redshift-space
distortions:} the Gold-2018 fσ8(z) growth compilation (a standard,
vetted set chosen for comparability with prior growth analyses;
combining it with DESI DR2 BAO follows common late-time practice, and
the audit below corrects a few mis-entered points), with the data audit
of App. C (a small number of mis-entered points corrected to their
published values). - \textbf{CMB:} the compressed Planck distance priors
(shift parameter R, acoustic scale \(l_{\mathrm{A}}\),
\(\omega _{\mathrm{b}}\)), with the full \(plik_{\mathrm{lite}}\)
likelihood used as a cross-check. - \textbf{Local distance ladder:} the
SH0ES H₀ prior. - \textbf{Out-of-sample (§5.3):} the pre-DESI eBOSS DR16
full-shape BAO, held entirely out of the main fit.

\subsection{4.2 CAMB-in-the-loop: a physical sound
horizon}\label{camb-in-the-loop-a-physical-sound-horizon}

The single most important methodological point is that the sound horizon
is \textbf{not} a free parameter. Earlier exploratory fits that floated
\(r_{\mathrm{d}}\) found large apparent preferences for SEDE; those did
not survive a physical sound horizon. Here, at every point in the chain,
\(r_{\mathrm{drag}}\), \(r_{\mathrm{star}}\), and θ* are computed
self-consistently from the background by CAMB (``CAMB-in-the-loop''), so
the BAO and CMB distances are tied to the same early-universe physics
for both models. This is what makes the ΔDIC of §5 a like-for-like
comparison rather than an artifact of a free standard ruler. The CMB
enters as the compressed (R, \(l_{\mathrm{A}}\)) distances computed on
the SEDE background through the same CAMB call. \emph{Scope of the
compressed likelihood, and its full-likelihood validation:} the
compressed priors capture the background distance to last scattering
(and hence the early-dark-energy fraction that drives the Δ leverage of
§5.6), but not the full acoustic-peak/polarization or lensing/ISW
information. We therefore \textbf{validate the headline against the full
primary CMB likelihood}: a CAMB-in-the-loop fit on the full Planck
\(plik_{\mathrm{lite}}\) TTTEEE + lowℓ TT/EE, with the five cosmological
parameters and the Planck nuisance \emph{all free}, gives
\textbf{Δχ²(SEDE−ΛCDM) = −0.43} --- the preference sitting in the high-ℓ
peak structure (−0.37), \emph{negative} (mildly favouring SEDE),
\textbf{not} the large positive Δχ² that the holographic-DE early--late
tension would produce (§5.1, §9; \texttt{run\_full\_cmb\_mcmc.py}). This
directly tests the acoustic-peak channel the compression omits, and we
have confirmed it against the \textbf{full non-lite plik with all 15
foregrounds sampled} (Δχ² = −0.33, consistent with the lite −0.43; §9).
\emph{Stated scope of what remains:} only the ACT DR6 \emph{primary}
multifrequency likelihood is not run (ACT \emph{lensing}, the
DE-sensitive channel, is, evaluated at the primary best-fit; §5.1). It
is \textbf{formally deferred} under a pre-registered trigger
(PREREGISTRATION.md §E): the \texttt{act\_dr6\_mflike} likelihood and
its multi-GB data products are not installed, and with three concordant
CMB datasets --- \(plik_{\mathrm{lite}}\), full non-lite plik, and ACT
DR6 lensing --- all at \textbar Δχ²\textbar{} \textless{} 1, a fourth
primary likelihood is not expected to overturn them; the committed
(falsifiable) expectation is \textbar Δχ²(SEDE−ΛCDM)\textbar{}
\textless{} 1, a large positive value being a genuine failure. The
\%-level metric effects (§4.4) are reported but not used in the headline
fit.

\subsection{4.3 Inference and model
comparison}\label{inference-and-model-comparison}

Parameters are sampled by MCMC, marginalising over
\(\Omega _{\mathrm{m}}\), H₀ (with the SH0ES prior),
\(\omega _{\mathrm{b}}\), the amplitude/σ8, and
\(r_{\mathrm{d}}\)-determining physics --- the \textbf{same five} for
SEDE and ΛCDM, since SEDE's dark sector is fixed (§3). Model comparison
uses the deviance information criterion (DIC), which penalises the
effective number of parameters \(p_{\mathrm{D}}\); because
\(p_{\mathrm{D}}\) is nearly equal (numerically close) for the two
models, ΔDIC reduces essentially to the χ² difference. Since DIC alone
is distrusted for non-Gaussian posteriors, we report the \textbf{AIC and
BIC as cross-checks} (\texttt{run\_info\_criteria.py}): with k = 5
fitted parameters for \emph{both} models (Δk = 0), AIC = χ² + 2k and BIC
= χ² + k ln N (N = 1628) give \textbf{ΔAIC = ΔBIC = ΔDIC = Δχ² = −2.95}
(−3.17 for the full-CMB joint) --- the criteria coincide because the
comparison is \textbf{penalty-free}, the preference is not bought with
complexity. We also compute a genuine \textbf{Bayesian evidence} by
nested sampling (dynesty, \texttt{run\_bayesian\_evidence.py}),
integrating the joint likelihood over the shared prior: ln Z = −737.6
(ΛCDM) and −735.7 (SEDE), giving \textbf{ln B(SEDE−ΛCDM) = +1.89 ± 0.33}
--- \emph{positive but weak} on the Jeffreys/Kass--Raftery scale,
favouring SEDE, with the Occam factor explicit (and ≈ 0, since the prior
volume is identical). The evidence slightly \emph{exceeds}
½\(\Delta \chi ^{2}_{\mathrm{min}}\) (1.48), i.e.~SEDE is well-supported
across the prior, not only at the best fit. With the contested SH0ES
prior dropped the evidence falls to \textbf{ln B = +0.95 ± 0.41} ---
still favouring SEDE but formally \emph{inconclusive} (\textbar ln
B\textbar{} \textless{} 1) --- so about half the evidence is carried by
SH0ES and the rest by the CMB distance (§5.4 ii′). Significance is then
assessed not by any single criterion but by the false-preference
calibration of §5.2.

\subsection{4.4 Perturbations: smooth dark energy, validated by
CLASS}\label{perturbations-smooth-dark-energy-validated-by-class}

SEDE's dark energy has \(c_{\mathrm{s}}^{2} = 1\) and does not cluster
(§3.4), so its growth is captured analytically by a smooth-dark-energy
treatment. We validated this against a full Boltzmann computation:
implemented in CLASS as a PPF fluid (which handles the w = −1 crossing),
the full-perturbation fσ8(z) matches the analytic background-only growth
to \textbf{0.1\%} over z = 0.3--1.5. The smooth-DE closure is therefore
numerically validated for the observables used here, and the
metric-level effects (CMB lensing \(C_\ell ^{\phi \phi }\), low-ℓ ISW,
S8) are all at the \%-level relative to ΛCDM.

\section{5. Results}\label{results}

\subsection{5.1 The joint fit and the equation of
state}\label{the-joint-fit-and-the-equation-of-state}

In the marginalised joint analysis, SEDE is \textbf{mildly preferred}
over ΛCDM, with a deviance information criterion difference ΔDIC ≈ −2.9
in SEDE's favour. Because the dark sector adds no fitted parameter (§3),
the two models share the same five parameters and the same effective
complexity,
\(p_{\mathrm{D}}(SEDE) \approx  p_{\mathrm{D}}(\Lambda CDM)\): the
preference is in fit, not bought with parameters. \textbf{The headline
survives the full primary CMB likelihood}, not only the compressed
priors: a marginalised CAMB-in-the-loop fit on the full Planck
\(plik_{\mathrm{lite}}\) TTTEEE + lowℓ TT/EE (all five cosmological
parameters and the Planck nuisance free) gives Δχ²(SEDE−ΛCDM) = −0.43,
the preference carried by the high-ℓ acoustic peaks (−0.37) with the
low-ℓ ISW-sensitive channel flat (≈ −0.03) --- i.e.~\emph{mildly
favouring} SEDE rather than the large positive Δχ² the standard
holographic-DE early--late tension would produce (§4.2, §9). More
decisively, \textbf{folding the full primary CMB \emph{into} the joint}
(replacing the compressed R, ℓ\_A priors with \(plik_{\mathrm{lite}}\)
TTTEEE + lowℓ, profiled over all parameters;
\texttt{run\_joint\_fullcmb.py}) leaves the preference intact and
marginally stronger: joint Δχ²(SEDE−ΛCDM) = \textbf{−3.17} (ΛCDM 2443.99
→ SEDE 2440.83), versus ≈ −2.9 with compressed CMB --- so the headline
is \textbf{not an artifact of the CMB compression}. The physical reason
the CMB poses no obstacle is that SEDE carries \emph{less} early dark
energy than ΛCDM and deforms only the late-time expansion (Figure 5):
there is no early--late wall of the kind that disfavours standard
holographic DE. The recovered background parameters are physical and
close to ΛCDM \((\Omega _{\mathrm{m}} \approx  0.30\), H₀ ≈ 68.8, σ8 ≈
0.76). The absolute χ²/dof ≈ 3 (§5.5) reflects the heterogeneous,
partly-compressed compilation (and the standard-but-not- goodness-of-fit
dof definition); it should \emph{not} be read as a goodness-of-fit
probability --- we use only matched-model differences (ΔDIC, Δχ²) on the
identical data vector for inference. The canonical (λ = 1/2) equation of
state \textbf{crosses w = −1}, with a present-day value near −1 and a
dip to ≈ −1.08 by z ≈ 2, i.e.~(w₀, \(w_{\mathrm{a}}) \approx  (-0.98\),
−0.11) --- squarely in DESI DR2's evolving-dark-energy
(thawing/crossing) quadrant, \textasciitilde2.7σ from the ΛCDM point.
The crossing is not imposed; it is where the expansion term and the
horizon-entropy-production term in 1 + w balance (Result 13, §6).

\pandocbounded{\includegraphics[keepaspectratio]{../../output/paper/fig3_datafit}}
\textbf{Figure 3.} Matched residual comparison (relative to ΛCDM).
\emph{Left:} DESI DR2 BAO residuals relative to the ΛCDM best fit, in
units of σ --- the data (slate) with the SEDE prediction (red); SEDE is
consistent with the measurements and sits between ΛCDM and the points
that pull away from it. \emph{Right:} Pantheon+ supernova residuals
(binned) versus ΛCDM, with the SEDE shape overplotted. SEDE and ΛCDM are
close at present precision; the preference is a multi-probe aggregate
(§5.4), not one decisive measurement.

\pandocbounded{\includegraphics[keepaspectratio]{../../output/paper/fig4_eos}}
\textbf{Figure 4.} The equation of state. \emph{Left:} the canonical
SEDE-H fluid w(z) (red) crosses −1 near z ≈ 0 and dips into the phantom
regime (≈ −1.08 by z ≈ 2), the same sense as DESI DR2's CPL fit (green)
and distinct from ΛCDM (dashed). \emph{Right:} the (w₀,
\(w_{\mathrm{a}}\)) plane --- SEDE sits between ΛCDM and the DESI DR2
w0waCDM contour, in the evolving-dark-energy quadrant. The crossing is
not imposed; it is where the expansion and horizon-entropy-production
terms in 1 + w balance (Result 13).

\pandocbounded{\includegraphics[keepaspectratio]{../../output/fig_cmb_earlyde}}
\textbf{Figure 5.} Why SEDE survives the CMB where standard holographic
dark energy does not --- the physical content of the full-CMB result
(full non-lite plik Δχ² = −0.33, \(plik_{\mathrm{lite}} -0.43\), ACT DR6
lensing −0.31; all \textbar Δχ²\textbar{} \textless{} 1). \emph{Left:}
the early-dark-energy fraction \(\Omega _{\mathrm{DE}}(z)\). SEDE's
structure gate \((f_{\mathrm{sat}}\) → 0 at high z) holds it
\emph{below} ΛCDM \((\Omega _{\mathrm{DE}}(z=3) = 0.029 < 0.035\) at
\(\Omega _{\mathrm{m}} = 0.30\)), whereas un-gated standard Barrow-HDE
carries substantially more early DE --- the runaway that drives the
early--late tension of refs. {[}Wu et al. 2509.02945{]}. \emph{Right:}
the fractional expansion deviation \(H/H_\Lambda CDM - 1\). SEDE's
deviation is sub-percent and confined to z ≲ 3, decaying to zero through
BBN and recombination, so the sound horizon, the acoustic-peak
positions, and the distance to last scattering are all essentially
unchanged; standard Barrow-HDE deviates over a far wider range. SEDE has
\emph{less} early dark energy than ΛCDM, not more --- hence no
early--late wall. (\texttt{run\_cmb\_earlyde.py}.)

\subsection{5.2 Is the preference real, or model
flexibility?}\label{is-the-preference-real-or-model-flexibility}

A ΔDIC ≈ −3 is suggestive but not, on its own, decisive --- and a
structure-gated model could in principle be winning by flexibility
rather than by capturing a real signal. We test this directly with a
\textbf{false-preference calibration} (a parametric bootstrap). We
generate mock datasets from a ΛCDM truth --- re-drawing \emph{every}
probe self-consistently, crucially including the compressed CMB
distances (R, \(l_{\mathrm{A}}\)) and the SH0ES prior, not just the
late-time data --- and refit both models to each mock. Under ΛCDM truth
a flexible model would still tend to ``win''; SEDE does not: the null
distribution of Δχ²(SEDE − ΛCDM) has mean \textbf{+0.42} (SEDE if
anything fits ΛCDM-truth mocks slightly \emph{worse}, consistent with
having no spare flexibility). The real data give Δχ² = −2.96, beyond
essentially all of the null draws: in a 2000-mock run (null mean +0.37,
s.d. 1.25), \textbf{8 mocks in 2000 are as SEDE-preferring as the real
data, p ≈ 0.004 (\textasciitilde2.6σ), 95\% upper limit 0.006}. (An
earlier 500-mock run gave p ≈ 0.002 but with only one mock at the
threshold --- resolution-limited; the 2000-mock value is the resolved,
slightly more conservative number.) Rerun with the contested SH0ES prior
removed from both the real fit and the mocks, the real value moves to
Δχ² = −1.75 and the calibration gives \textbf{p = 0.030
(\textasciitilde1.9σ, 95\% UL 0.050)} --- the preference degrades but
stays beyond the null, now carried by the CMB-distance channel alone
(§5.4 ii′).

Two caveats. First, an earlier version of this test held the CMB
distances \emph{fixed} across mocks; that biased the null mean to −2.87
and gave p ≈ 0.5 --- the artifact that made the preference look like
pure flexibility. Re-drawing the CMB self-consistently is the correct
procedure and is what yields p ≈ 0.004. Second, the significance is
\textbf{CMB-distance-treatment-dependent}: an independent sibling
pipeline, with a different \(l_{\mathrm{A}}/CAMB\) handling, obtains a
milder real Δχ² ≈ −1.98 and p ≈ 0.025 (\textasciitilde2σ). Both analyses
agree the preference is \textbf{not obviously attributable to model
flexibility}, and trace the 2σ--3σ spread to the CMB-distance engine
rather than to SH0ES. We therefore quote the preference as a
\textbf{suggestive \textasciitilde2--3σ}, pending independent
reproduction.

\pandocbounded{\includegraphics[keepaspectratio]{../../output/fig_F5_calibration}}
\textbf{Figure 6.} The false-preference calibration. Histogram: the null
distribution of Δχ²(SEDE − ΛCDM) from ΛCDM-truth mocks (all probes,
including the compressed CMB and SH0ES, re-drawn self-consistently),
centred near zero (mean +0.37) --- SEDE has no flexibility advantage
under ΛCDM truth. Red line: the real data, Δχ² = −2.96, in the far tail.
In the 2000-mock run only 8 mocks are as SEDE-preferring as the real
data, p ≈ 0.004 (\textasciitilde2.6σ, 95\% UL 0.006); the preference is
not obviously attributable to flexibility.

\subsection{5.3 Out-of-sample check}\label{out-of-sample-check}

A further check against a flexibility explanation is out-of-sample
prediction. We trained SEDE on \textbf{real pre-DESI data only} (eBOSS
DR16 full-shape BAO, LRG/QSO) and used it to predict the later, held-out
DESI DR2 BAO: in this split SEDE matches the held-out DESI DR2 trend
\textbf{better than ΛCDM by Δχ² = −8.2}; a low-z→high-z split of the
same kind gives −9.9. A model winning purely by flexibility would tend
to do \emph{worse} out of sample, so this is encouraging --- i.e.~a
pre-DESI-trained SEDE realisation \textbf{predicts the DESI DR2 trend}
in this split. \emph{Two clarifications so this number is read
correctly.} (a) \textbf{Not comparable to the in-sample ΔDIC.} The −8.2
is the Δχ² on the \textbf{held-out 13-point DESI BAO subset alone},
whereas the in-sample ΔDIC ≈ −2.95 is the net over the \textbf{full
1628-point compilation} (most probes of which do not discriminate;
§5.4); the two are different objects on different data, so ``OOS −8.2
exceeds in-sample −3'' is \emph{expected}, not anomalous (a single-probe
Δχ² is naturally larger than a diluted whole-compilation total) and is
\textbf{not} evidence of over-fitting (which would push OOS the other
way). (b) \textbf{Direction-dependent.} The signal lives in the low-z →
high-z direction (the DESI evolving-DE direction, Δχ² = −9.9); the
reverse split (high-z → low-z) is flat and marginally favours ΛCDM
(+0.3), so part of the OOS strength is that eBOSS-trained parameters
land where DESI's specific tracers sit --- partly favourable, to be
independently reproduced before strong weight is placed on it.

\pandocbounded{\includegraphics[keepaspectratio]{../../output/paper/fig5_oos}}
\textbf{Figure 7.} Out-of-sample test. BAO residuals relative to ΛCDM
(in σ) for the \textbf{training} set (pre-DESI eBOSS DR16, purple) and
the \textbf{held-out} DESI DR2 data (slate), with the SEDE prediction
(red) from parameters fit to eBOSS alone. SEDE, trained without ever
seeing DESI, predicts the held-out DESI BAO better than ΛCDM (Δχ² =
−8.2) in this split --- encouraging, and to be independently reproduced.

\subsection{5.4 Robustness}\label{robustness}

The preference is not a single-dataset artifact. (i)
\textbf{Supernova-set independence:} Δχ²(SEDE − ΛCDM) is −2.95 / −2.87 /
−2.97 with Pantheon+ / DES-5YR / Union3 respectively --- not a Pantheon+
effect. (ii) \textbf{Per-probe decomposition (where the preference comes
from).} Decomposing the joint Δχ²(SEDE − ΛCDM) = −2.95 into per-probe
contributions \emph{at the joint best-fit}
(\texttt{run\_probe\_decomposition.py}) is more informative than the
leave-one-out, and we report it plainly: the preference is carried by
the \textbf{geometry channels} --- the compressed CMB distance (R, ℓ\_A)
contributes \textbf{−1.73} and the SH0ES H₀ prior \textbf{−1.25}
(together ≈ 100\% of the total) --- while DESI BAO (−0.14), fσ8 (−0.03)
and \(\omega _{\mathrm{b}}\) (−0.07) are nearly flat and Pantheon+ SN
(+0.16) and cosmic chronometers (+0.11) marginally favour ΛCDM. So this
is \textbf{not a broad multi-probe consensus}: SEDE's evolving w(z) fits
the distance to last scattering better and accommodates a marginally
higher H₀ (68.9 vs 68.7) that eases SH0ES --- a \emph{paired}
(CMB-distance + H₀) effect. We show the decomposition precisely because
the \textbf{leave-one-probe-out is blind to this paired structure}
(dropping either member leaves the other, so every single-probe holdout
stays negative: range {[}−3.06, −2.16{]}, mean −2.64) --- the holdout
robustness is real but does not by itself establish a consensus. Two
honest qualifications follow. First, the CMB-distance contribution is
\emph{not} a compressed-prior artifact: folding the full primary CMB
into the joint \textbf{strengthens} it (Δχ² = −3.17, §5.1). Second, SEDE
does \textbf{not} resolve the Hubble tension --- its best-fit H₀ ≈ 69
still sits \textasciitilde4σ from SH0ES (§7, §9) --- but its geometry
does ease the SH0ES prior by Δχ² ≈ 1.2, so part of the preference is a
mild H₀ accommodation, which we flag rather than hide. (ii′)
\textbf{No-SH0ES robustness (the contested prior removed).} Because the
SH0ES H₀ prior is the field's most disputed input --- and DESI's own
baseline excludes it --- we recompute all three headline statistics with
SH0ES dropped (\texttt{run\_no\_shoes\_robustness.py}). The preference
\textbf{weakens but survives, and what remains sits in the harder
channel.} The joint Δχ²(SEDE − ΛCDM) = \textbf{−1.75} (from −2.95), now
carried by the compressed CMB distance (−1.37) and DESI BAO (−0.45),
with the best-fit H₀ falling to \textbf{68.2} once it is no longer
pulled toward 73 --- i.e.~the CMB-distance improvement is genuine, not
an H₀ accommodation. The nested-sampling evidence falls to \textbf{ln B
= +0.95 ± 0.41} --- still favouring SEDE but now \emph{inconclusive}
(\textbar ln B\textbar{} \textless{} 1) on the Jeffreys scale --- and
the false-preference calibration (§5.2) gives \textbf{p = 0.030} (6/200
ΛCDM-truth mocks beyond the real value; \textasciitilde1.9σ, 95\% UL
0.050) versus p ≈ 0.004 (\textasciitilde2.6σ) with SH0ES. So roughly
40\% of the χ² preference and about half the Bayesian evidence are
carried by SH0ES; the remainder is \textbf{modest but robust and lives
in the CMB-distance channel rather than the contested local prior.} We
read this as sharpening, not weakening, the paper's ``suggestive, not
established'' posture: with the one disputed prior removed, SEDE remains
\emph{mildly} preferred (\textasciitilde1.9σ) by hard geometric data,
and no part of that residual depends on resolving the Hubble tension.
(iii) \textbf{Code independence:} the SEDE expansion history reproduced
with CLASS matches CAMB to 5×10⁻⁵ in \(r_{\mathrm{drag}}\) (≪ 0.2\%) ---
not a CAMB artifact. (iv) \textbf{Growth:} the full-Boltzmann
\(\gamma _{\mathrm{growth}} \approx  0.55\) confirms standard-GR dark
energy, not modified gravity. (v) \textbf{Tracer-level robustness:}
recent analyses argue the DESI DR2 evolving-DE signal is driven mainly
by the LRG1 (z≈0.51) and LRG2 (z≈0.71) BAO bins {[}arXiv:2508.10514,
2504.15222{]}. Dropping these bins individually and together
(\texttt{run\_tracer\_loo.py}) leaves SEDE's preference intact ---
Δχ²(SEDE − ΛCDM) = −3.27 (−LRG1), −2.76 (−LRG2), and \textbf{−3.00 with
both removed} (versus −2.95 for the full set), negative in every tracer
holdout (range {[}−3.27, −2.76{]}). SEDE's mild preference is therefore
\emph{not} the LRG1--2 artefact the steep-CPL signal is attributed to:
it is a geometry-channel preference (CMB-distance + H₀, §5.4(ii)) that
does not rest on the contested bins. The dissociation is quantitative:
the ΛCDM-deviation measured by the \(ratio(\omega _{\mathrm{m}})\)
early-vs-late consistency metric drops from 2.6σ to 1.2σ when LRG1--2
are removed {[}arXiv:2407.04385{]}, whereas SEDE's Δχ² is essentially
unchanged (−2.95 → −3.00) --- SEDE's preference and the LRG-driven
steep-CPL deviation are \emph{different objects}. This is the natural
answer to the post-DR2 robustness critiques (§5.7).

\subsection{5.5 Evidential summary}\label{evidential-summary}

SEDE is mildly preferred over ΛCDM (ΔDIC ≈ −2.9; calibration
\textasciitilde2--3σ); the preference is suggestive --- not obviously
attributable to model flexibility, and supported by a pre-DESI → DESI
out-of-sample check --- and the model adds no continuously-sampled
dark-energy parameter. But it is modest: both models have χ²/dof ≈ 3, so
SEDE is ``less wrong,'' not ``right.'' This is stronger than
indistinguishable, short of established. DESI DR3 and Euclid provide the
sharp future test (§6).

\subsection{5.6 The deformation Δ from current data --- a conditional
diagnostic}\label{the-deformation-ux3b4-from-current-data-a-conditional-diagnostic}

The Barrow parameter Δ is a useful \emph{test} axis: the volume
postulate corresponds to Δ = 1, the area law to Δ = 0. SEDE's
\emph{equation of state} is CPL-degenerate in shape
(max\textbar{}\(w - w_{\mathrm{CPL}}\)\textbar{} = 0.027), so a
diagnostic must use channels other than the EOS curvature. We stress at
the outset that the numbers below are \textbf{conditional diagnostics}
(computed with the other cosmological and nuisance parameters held at or
near their best-fit values, not a full marginalised likelihood with Δ
varied); they indicate where current data sit, not a marginalised
detection. Four channels:

\begin{longtable}[]{@{}
  >{\raggedright\arraybackslash}p{(\linewidth - 4\tabcolsep) * \real{0.3333}}
  >{\raggedright\arraybackslash}p{(\linewidth - 4\tabcolsep) * \real{0.3333}}
  >{\raggedright\arraybackslash}p{(\linewidth - 4\tabcolsep) * \real{0.3333}}@{}}
\toprule\noalign{}
\begin{minipage}[b]{\linewidth}\raggedright
channel
\end{minipage} & \begin{minipage}[b]{\linewidth}\raggedright
what it uses
\end{minipage} & \begin{minipage}[b]{\linewidth}\raggedright
Δ̂ (conditional profile --- \textbf{not a posterior})
\end{minipage} \\
\midrule\noalign{}
\endhead
\bottomrule\noalign{}
\endlastfoot
CMB shift parameter R (Planck) & early-DE fraction,
\(\Omega _{\mathrm{DE}}(z_{\mathrm{rec}}) \propto  H^{-\Delta }\) & 0.98
± 0.11 \\
lever arm: DESI BAO + Lyα(z=2.33) + R &
\(\rho _{\mathrm{DE}}/\rho _{\mathrm{crit}} \propto  H^{-\Delta }\)
across the H-range & 1.03 ± 0.11 \\
DESI DR2 (w₀,\(w_{\mathrm{a}}\)) quadrant (official chains) & EOS values
+ crossing & 0.93 / 0.83 / 0.90 (PantheonPlus/DESY5/Union3) \\
growth fσ8 (Gold-2018) --- \textbf{orthogonal} & growth amplitude, not
distance & 0.0 ± 0.61 \\
\end{longtable}

\begin{quote}
\textbf{These are not posterior constraints on Δ; they are
profile/conditional diagnostics} --- each Δ̂ is obtained with the other
cosmological and nuisance parameters held at (or near) their best-fit
values, \emph{not} marginalised. The quoted ± is the conditional
curvature width, \textbf{not} a marginalised error bar; read each row as
``the data lean to Δ ≈ 1 in this channel,'' not as a measurement. The
robust, marginalised Δ measurement awaits DESI DR3 + Euclid (§6).
\end{quote}

These diagnostics \textbf{prefer the volume endpoint (Δ ≈ 1) and
disfavour the area-law branch (Δ = 0)}, consistent with the EOS argument
of §8.1 (the area-law branch has w₀ = −0.49 and an outsized early-DE
fraction). The growth channel is weak but \textbf{orthogonal} (growth
amplitude, not expansion) and is consistent with the geometric value at
1.6σ --- no growth--geometry tension.

The high apparent significance of the area-law disfavouring is a
\textbf{conditional} statement and must not be read as a marginalised
detection: the CMB and lever-arm channels share the high-z expansion, so
their combination is an upper bound on the information, and it is
conditional on \(\Omega _{\mathrm{m}}\), H₀, \(r_{\mathrm{d}}\) ---
marginalising, the area-law signal is partly reabsorbed by shifts in
those parameters. This is exactly why the \emph{joint, marginalised}
model preference is only modest (ΔDIC ≈ −3, §5.1--§5.5), not
overwhelming; the two statements are not in tension once the
conditional-vs-marginal distinction is kept. The (w₀,\(w_{\mathrm{a}}\))
channel also pulls slightly low because SEDE's gentle \(w_{\mathrm{a}}\)
under-matches DESI's steeper preferred \(w_{\mathrm{a}}\) --- the same
mild \textasciitilde2.7σ EOS tension visible in the joint fit (Fig. 4).
\textbf{Honest summary: current late-time data prefer the volume
endpoint and disfavour the area-law branch, but a robust
\emph{marginalised} measurement of Δ awaits DESI DR3 and Euclid (§6).}

\subsection{5.7 Relation to the post-DESI DR2
literature}\label{relation-to-the-post-desi-dr2-literature}

The DESI DR2 BAO release {[}arXiv:2503.14738{]} sharpened interest in
evolving dark energy, with a preference for the (w₀ \textgreater{} −1,
\(w_{\mathrm{a}} < 0\)) quadrant; the subsequent literature is actively
debating whether that signal is physical or an artefact of dataset
combination, tracer selection, or supernova calibration {[}reviewed in
arXiv:2602.05368{]}. Two threads bear directly on SEDE. First,
\textbf{robustness critiques} argue the steep signal reflects CMB/BAO/SN
combination tension {[}arXiv:2504.15222{]} and is driven mainly by the
LRG1--2 BAO bins {[}arXiv:2508.10514, 2407.04385{]}. This \emph{helps}
SEDE rather than threatening it: SEDE does not chase the steep
\(w_{\mathrm{a}}\) (it sits \textasciitilde2.7σ from the DESI CPL point,
§5.6), and its preference survives a \textbf{tracer-level} leave-one-out
that removes exactly those LRG1--2 bins (§5.4, Δχ² = −3.00 with both
dropped) --- so SEDE realises the \emph{conservative} reading the
critiques call for, rather than the fragile steep-CPL one. Second,
\textbf{reconstruction and mechanism papers}: model-agnostic
Gaussian-process reconstructions place the w = −1 crossing at
\(z_{\mathrm{wt}} \approx  0.46\) (+0.24/−0.12) {[}arXiv:2511.02220{]},
versus SEDE's canonical z ≈ 0.19 --- a ∼2.3σ offset (SEDE crosses more
recently) that is a genuine future discriminator, not yet a tension
given the broad reconstruction errors. Among mechanism competitors,
data-driven analyses read the evolving signal as a \emph{coupled} dark
sector (∼2.2σ DE--DM interaction at low z) {[}arXiv:2504.00985{]};
SEDE's CHR ledger is exactly such a coupling but with the interaction
\emph{fixed} to the structure-growth rate, \(c_{\mathrm{s}}^{2} = 1\),
and instability-free (§7), rather than a free coupling. And the
\textbf{quintom} programme {[}arXiv:2505.24732, 2601.02284{]} rests on a
\emph{no-go theorem} --- no single canonical field or perfect fluid
crosses w = −1, so viable models require two fields, higher derivatives,
modified gravity, or an effective-field-theory description. SEDE is the
last of these: it crosses w = −1 \emph{effectively}, as a smooth horizon
fluid in the GW-safe EFT corner
\((\alpha _{\mathrm{T}} = \alpha _{\mathrm{M}} = 0\),
\(c_{\mathrm{s}}^{2} = 1\); §7), avoiding the ghost/two-field machinery,
and --- unlike the two-field quintom whose late attractor is
phantom-dominated (a Big-Rip tendency) {[}arXiv:2601.02284{]} --- SEDE's
future attractor is de Sitter (w → −1, with \(f_{\mathrm{sat}}\)
\emph{saturating} to a finite ≈1.23 as growth freezes, not diverging;
§2, App. A.7, Result 12), so it does not predict a Big Rip.

Two further consistency points sharpen the post-DR2 placement. (a)
\textbf{Turyshev's model filter.} The DR2 status review
{[}arXiv:2602.05368{]} screens evolving-DE models by
\emph{gravitational-wave propagation} and \emph{perturbation stability},
and stresses that the CPL preference is sensitive to redshift-dependent
SN calibration at the few×10⁻² mag level. SEDE passes the filter by
construction \((\alpha _{\mathrm{T}} = 0\) ⟹ \(c_{\mathrm{GW}} = c\);
\(c_{\mathrm{s}}^{2} = 1\) ⟹ gradient/ghost-stable; §7), and its SN-set
independence (Δχ² = −2.95/−2.87/−2.97 across Pantheon+/DES-5YR/Union3;
§5.4) limits the calibration exposure. (b) \textbf{The
\(r_{\mathrm{d}}\)-independent shape test.} Turyshev's calibration-free
observable
\(F_{\mathrm{AP}}(z) \equiv  D_{\mathrm{M}}/D_{\mathrm{H}} = (D_{\mathrm{M}}/r_{\mathrm{d}})/(D_{\mathrm{H}}/r_{\mathrm{d}})\)
cancels the sound horizon and is also H₀-independent, isolating the pure
expansion shape E(z). On the six DESI DR2 tracer bins that report both
ratios (\texttt{run\_fap\_test.py}), SEDE's evolving-w shape gives χ²/n
= 0.96 --- consistent with the data and indistinguishable from ΛCDM (Δχ²
= −0.14) --- with \textbf{no} reliance on \(r_{\mathrm{d}}\), H₀, the
CMB sound horizon, or SN calibration, the very systematics the
robustness debate targets. (The LRG1 bin sits ∼1.8σ low for \emph{both}
SEDE and ΛCDM, i.e.~it is a data feature, not a model failure.) SEDE is
thus a \emph{physical} member of the post-DR2 evolving-DE family,
distinguished by its mechanism (the growth--expansion lock, §6), by
passing the stability/GW filter, and by surviving the robustness and
calibration-free tests the steep signal may not.

\section{6. Predictions \&
falsifiability}\label{predictions-falsifiability}

Because SEDE adds no fitted dark-energy parameter, it makes sharp
predictions rather than fits. We \textbf{pre-registered} these (a
sha256-locked statement deposited ahead of the future data; App. D) to
forbid post-hoc tuning:

\begin{longtable}[]{@{}
  >{\raggedright\arraybackslash}p{(\linewidth - 4\tabcolsep) * \real{0.3333}}
  >{\raggedright\arraybackslash}p{(\linewidth - 4\tabcolsep) * \real{0.3333}}
  >{\raggedright\arraybackslash}p{(\linewidth - 4\tabcolsep) * \real{0.3333}}@{}}
\toprule\noalign{}
\begin{minipage}[b]{\linewidth}\raggedright
quantity
\end{minipage} & \begin{minipage}[b]{\linewidth}\raggedright
SEDE prediction
\end{minipage} & \begin{minipage}[b]{\linewidth}\raggedright
discriminating from
\end{minipage} \\
\midrule\noalign{}
\endhead
\bottomrule\noalign{}
\endlastfoot
Barrow deformation Δ & \textbf{1.0 ± 0.09} (forecast) & a smooth horizon
Δ = 0 \\
equation of state (w₀, \(w_{\mathrm{a}}\)) & ≈ (−0.98, −0.11), crossing
−1 & the ΛCDM point (−1, 0) at \textasciitilde2.7σ \\
w = −1 crossing redshift (consistency check, \emph{not} a discriminator)
& z ≈ 0.19 (genuine volume-law value; ∼2.2σ mild tension) & GP
reconstructions \(z_{\mathrm{wt}} \approx  0.46\) (+0.24/−0.12)
{[}arXiv:2511.02220{]} \\
structure amplitude σ8 & ≈ 0.76 & (consistent with weak-lensing S8) \\
growth index \(\gamma _{\mathrm{growth}}\) & ≈ 0.55 & modified gravity
\((\gamma _{\mathrm{growth}} \neq  0.55)\) \\
\end{longtable}

\textbf{An honest near-term tension --- and why z ≈ 0.19 is the genuine
prediction, not a knob.} Canonical SEDE crosses w = −1 at
\textbf{\(z_{\mathrm{cross}} = 0.195\)} (w₀ = −0.996), using the
Gibbons--Hawking horizon temperature \(T_{\mathrm{AH}} = H/2\pi\). This
is \textbf{not} an arbitrary convention: T = H/2π is the de Sitter
temperature appropriate for a de Sitter-attractor dark-energy fluid, and
we \emph{tested} the alternative --- applying the full dynamical
Cai--Kim factor (1−ε/2) to the volume-law fluid
(\texttt{run\_dynT\_crossing.py}, the self-consistent
\(\rho _{\mathrm{DE}} \propto  H(1-\epsilon /2)f\) model). It
\textbf{fails}: it has no stable self-consistent background, and in the
stable perturbative limit it gives w₀ = −1.15 and stays phantom (w
\textless{} −1) at all z --- \textbf{no crossing at all}, and in the
\emph{wrong} DESI quadrant (w₀ \textless{} −1). The dynamical correction
is largest in the matter era, exactly where \(f_{\mathrm{sat}}\) → 0 and
dark energy is negligible; forcing it where it does not belong destroys
the late-time crossing. So T = H/2π is the physically-correct choice for
this sector, and z ≈ 0.19 is the genuine volume-law prediction. (The
unrelated z ≈ 0.59 is the \emph{area-law} Δ = 0 model ---
\(\rho _{\mathrm{DE}} \propto  H^{2}(1-\epsilon /2)f\) --- a different
model, not a convention of this one.) Post-DR2 Gaussian-process
reconstructions favour \(z_{\mathrm{wt}} \approx  0.46\) (+0.24/−0.12)
{[}arXiv:2511.02220, 2504.00985{]}; canonical SEDE sits \textbf{∼2.2σ
more recent} --- a \emph{real, mild} tension (the reconstruction's lower
error is tight), the one place current data pull against canonical SEDE.
We flag it as such, \textbf{not} as a discriminator: the crossing
redshift is CPL-degenerate (§6, below) and, read as a Δ-probe, would
point the \emph{wrong} way (toward Δ = 0), conflicting with the early-DE
diagnostics that favour Δ = 1 (§5.6). The model's actual test is the
deformation amplitude Δ, not the crossing; a DR3/Euclid measurement of
the crossing is a useful consistency check, not the decider.

\textbf{The sharp future test.} The discriminating observable is the
\emph{amplitude} of the deformation Δ, not the detailed shape of w(z)
--- indeed the SEDE w(z) is CPL-degenerate in shape to better than the
DR3 precision (max \textbar w − CPL\textbar{} ≈ 0.027), so a CPL fit
cannot distinguish the \emph{mechanism}; Δ can. A Fisher forecast in
\((\Omega _{\mathrm{m}}\), Δ, ln A \(r_{\mathrm{d}}\)) gives σ(Δ) ≈
0.115 from DESI DR3 alone, 0.133 from Euclid, and \textbf{σ(Δ) ≈ 0.087
combined} --- an \textasciitilde11σ separation of the volume-law horizon
(Δ = 1) from a smooth one (Δ = 0), under the Fisher assumptions. DESI
DR3 and Euclid could therefore move SEDE substantially toward
confirmation or exclusion under those assumptions.

\textbf{The mechanism test is also forecast to be decisive ---
separately from the amplitude.} The σ(Δ) above is the
\emph{geometry-side} deformation; the \emph{distinctive} SEDE claim
(dark energy sourced by structure) is tested by measuring Δ on the
\textbf{growth} side and checking it against the geometry side --- the
E1/P2 growth--expansion lock. A Fisher forecast for the growth-side
deformation (\texttt{run\_mechanism\_forecast.py}) gives
\(\sigma (\Delta )_{\mathrm{growth}} = 0.25\) from Euclid/LSST weak
lensing alone, tightening to \textbf{0.13 (∼4σ)} when combined with
CMB-lensing, cluster abundance, RSD, and kSZ. Reframed as the
parameter-free \textbf{null test}
\(r(z) \equiv  \rho _{\mathrm{DE}}^{geom}(z) - G[D^{growth}(z)] = 0\)
--- dark energy a fixed function of the linear growth factor --- the
test uses the full redshift shape with no dark-energy parameters; ΛCDM
\((\rho _{\mathrm{DE}} = const\), no D-dependence) and Gough's
information-DE \((\rho _{\mathrm{DE}} = G\)′{[}SMD{]}, baryonic) violate
it distinctly. So Euclid/LSST-era growth data are forecast to
\textbf{confirm or refute the structure-sourcing mechanism at ∼4σ}, not
merely the EOS shape --- closing the gap between ``a w(z) with a
thermodynamic story'' and a tested mechanism (§9). The exact data vector
(the geometry leg DESI DR3 \(D_{\mathrm{H}}/r_{\mathrm{d}}\)→H(z) vs the
growth leg Euclid/LSST 3×2pt + CMB-lensing + clusters + RSD) and the
\textbf{pre-registered decision rule} --- CONFIRM if
\textbar{}\(\Delta _{\mathrm{grow}} - \Delta _{\mathrm{geom}}\)\textbar{}
\textless{} 2σ with the null
\(r(z) = \rho _{\mathrm{DE}}^{geom} - G[D^{grow}] = 0\) at χ²/dof ≲ 1;
REFUTE if \(\Delta _{\mathrm{grow}}\) is consistent with 0 at
\textgreater3σ while \(\Delta _{\mathrm{geom}} = 1\) --- are locked in
the pre-registration (PREREGISTRATION.md §F; App. D), so the staged test
carries no post-hoc freedom.

\pandocbounded{\includegraphics[keepaspectratio]{../../output/paper/fig6_forecast}}
\textbf{Figure 8.} The sharp future test. A Fisher forecast of the
separation of the volume-law horizon (Δ = 1) from a smooth one (Δ = 0):
DESI DR3 reaches 8.7σ, Euclid 7.5σ, and the combination σ(Δ) ≈ 0.087, an
\textasciitilde11σ separation. The mild present preference would become
a strong test with the next generation of surveys.

\pandocbounded{\includegraphics[keepaspectratio]{../../output/fig_e1_mechanism}}
\textbf{Figure 9.} The structure-sourcing mechanism on \emph{current}
data (the companion to the forecast of Figure 8, and the paper's honest
open frontier). SEDE's founding claim is that the activation fraction
\(f_{\mathrm{sat}}\) read off the \textbf{expansion} (geometry leg: DESI
\(D_{\mathrm{H}}/r_{\mathrm{d}}\) → H(z), blue) coincides with the one
read off \textbf{structure growth} (RSD fσ8 → σ8(z) → D(z), red); the
black curve is the SEDE prediction both should follow. The two legs
\textbf{track at z \textless{} 1} but \textbf{diverge at z
\textgreater{} 1} (the structure leg flattens while the geometry/model
declines). The significance is error-budget-dependent --- χ²/dof = 0.40
in the RSD-error budget (consistent) versus 1.40 in the precise-geometry
budget (marginally disfavoured) --- so on present data the mechanism is
\textbf{inconclusive and RSD-limited}: SEDE's EOS is favoured, but its
distinctive mechanism is not yet confirmed. Euclid/LSST weak-lensing
tomography sharpens the structure leg and decides this at ∼4σ (Figure 8;
pre-registered decision rule, PREREGISTRATION.md §F).
(\texttt{run\_e1\_figure.py} / \texttt{run\_e1\_mechanism.py}.)

\textbf{The EOS as a meter of entropy production.} SEDE gives the
equation of state a thermodynamic meaning (Result 13): within the SEDE
background ansatz, an exact identity
\(1 + w_{\mathrm{DE}} = (1/3)[2\lambda \epsilon  - d\) ln
\(f_{\mathrm{sat}}/d\) ln a{]} splits 1 + w into an expansion term
(which pushes w \textgreater{} −1) and a horizon-entropy-production term
(which pushes w \textless{} −1). The phantom crossing at z ≈ 0.2 is
exactly where the two balance. Measuring w(z) thus probes the effective
activation rate of the horizon-fluid sector --- the
structure-growth-gated arrow of time made observable in the dark-energy
equation of state.

\textbf{Structure sourcing without an entropy-amplitude problem.} The
identity above gives w(z) a thermodynamic reading, and it invites the
natural objection that a tiny structure-sourced entropy cannot gate an
order-unity dark-energy term --- the deposited binding entropy is only
\(\epsilon _{\mathrm{dep}} \approx  \Omega _{\mathrm{m}}\)
\(f_{\mathrm{coll}}\) ⟨\(\sigma _{\mathrm{v}}^{2}\)⟩/c² ≈ 1.5 × 10⁻⁷ of
the horizon entropy (the same six orders by which the binding
\emph{energy} falls short of \(\rho _{\mathrm{DE}}\); §2). \textbf{The
objection is a category error, and the amplitude problem is
\emph{avoided}, not solved by amplification.} The activation coordinate
is \textbf{not} the deposited-entropy fraction
\(\Delta S_{\mathrm{bind}}/S_{\mathrm{AH}}\); it is the structure
variance x ≡ σ²(z)/σ²(0) = D²(z), an \textbf{order-unity} quantity,
which is the \emph{argument} of the gate \(f_{\mathrm{sat}}\) (the
Press--Schechter collapsed fraction; §3.1, App. A.2). The deposited
binding entropy enters in a \emph{different} place: it fixes the entropy
\emph{weight} γ ≈ 1.5 --- the \emph{shape} of the kernel --- through the
heat \(\Delta S = E_{\mathrm{bind}}/T_{\mathrm{AH}}\) delivered to the
horizon (§3.1, App. A.3). Hence
\(\Delta S_{\mathrm{bind}}/S_{\mathrm{AH}}\) ∼ 10⁻⁷ does \textbf{not}
imply \(f_{\mathrm{sat}}\) ∼ 10⁻⁷: they are different objects. With the
\emph{argument} (variance, O(1)) and the \emph{weight} (entropy, 10⁻⁷)
kept distinct, the gate activates simply because the structure variance
grows by order unity --- nothing is amplified and nothing is tuned. This
is a consequence of the fixed inputs of §3, \textbf{derived, not
assumed.}

\textbf{The exact growth--expansion lock.} Write the gate as
\(\phi  \equiv  f_{\mathrm{sat}}(x)\), x = D²(z), and define its derived
response to the activation coordinate,
\(\chi _{\mathrm{x}} \equiv  \partial \phi /\partial x = d\)
\(f_{\mathrm{sat}}/dx = \gamma\) \(e^(-\gamma x)/(1 - e^(-\gamma ))\),
which, once γ is fixed (§3.1), is an \textbf{order-unity} function with
no free parameter --- \(\chi _{\mathrm{x}}(x=0) \approx  1.93\),
\(\chi _{\mathrm{x}}(x=1) \approx  0.43\). The fluid equation of state
then follows exactly from continuity (the entropy-production identity
above), using d ln φ/d ln
\(a = (\chi _{\mathrm{x}}/\phi )\cdot x\cdot (d\) ln x/d ln a) and d ln
x/d ln a = 2g with g ≡ d ln D/d ln a: \(1 + w_{\mathrm{DE}} = (1/3)[\)
\(2\lambda \epsilon  - (\chi _{\mathrm{x}}/\phi )\cdot 2x\cdot g\) {]},
verified against the direct fluid EOS to 3 × 10⁻³
(\texttt{run\_chr\_experiments.py}; suite tests V52--V53). Because
\(\chi _{\mathrm{x}}\), x, and g are each derived or directly measured,
the phantom term is \textbf{growth data}, not a free function. This is
the \textbf{w(z) ↔ σ²(z) lock}: given the measured expansion history,
the dark-energy equation of state and the structure-growth history are
not independent; inverting the relation predicts g(z) from w(z) to 0.6\%
within the model. A kinematic w(z) --- including a CPL fit that
reproduces SEDE's w(z) to within the DR3 precision --- obeys no such
relation. \textbf{Expansion-only mimicry is therefore not enough: SEDE
requires growth--expansion phase-locking.} This is the test (P2 below)
that separates the structure-sourcing mechanism from a curve fit --- w
from BAO/SNe/CMB distances versus g, D, fσ8 from RSD/weak lensing. It is
parameter-free and uses nothing beyond §3; no criticality assumption
enters.

\textbf{The observable form of the lock (P2), and a GR consistency
null.} The lock derived above --- the phantom part of w(z) equals
\((1/3)(\chi _{\mathrm{x}}/\phi )\cdot 2x\cdot g\) reconstructed from
growth data, obeyed by \emph{no} kinematic dark energy --- is the
model's one parameter-free, degeneracy-breaking prediction. Its
established observational realisation is the \(E_{\mathrm{G}}\)
estimator (the lensing-to-velocity amplitude ratio): SEDE is GR with
smooth dark energy, so it predicts no gravitational slip,
\(E_{\mathrm{G}}(z) = \Omega _{\mathrm{m}}\),0/f(z), within \textless1\%
of ΛCDM and consistent with the DESI-DR1 + weak-lensing measurement to z
≈ 1 {[}arXiv:2507.16098{]} (\texttt{run\_eg\_test.py}).
\(E_{\mathrm{G}}\) is a \emph{shared} GR null (with the GW-speed
\(c_{\mathrm{T}} = c\) and the zero-slip nulls SEDE passes), not a
SEDE-vs-ΛCDM discriminator --- at 0.3\% it sits far below the ∼5--20\%
\(E_{\mathrm{G}}\) errors; the model's only discriminator remains the
deformation amplitude Δ (§5.6, §6). \textbf{Failure of the lock (P2)
would falsify the structure-sourcing mechanism itself} --- SEDE would
then be an ordinary kinematic w(z) --- so P2 is the test that separates
the mechanism from a curve fit.

\textbf{An optional near-critical layer (deferred to a companion
paper).} If, in addition, the horizon sits near the area↔volume spinodal
--- the flat-landscape limit, i.e.~the Δ → 1 endpoint itself --- the
response sector becomes \emph{near-critical} (the \textbf{Critical
Horizon Response} regime) and acquires fluctuation-level signatures: a
localised enhancement of large-scale activity at the transition z* ≈ 1.2
(P3); a small, \textbf{nonzero}, redshift-localised
\emph{separate-universe} (long-wavelength) modulation of growth/ISW
peaked at z* (P4 --- the kill test, distinct from the \emph{broad,
sub-horizon} signal that clustering dark energy with
\(c_{\mathrm{s}}^{2} < 1\) would give, and from ΛCDM's exact zero); and
a critical-slowing transient in \(\gamma _{\mathrm{growth}}\) (P5).
These are \textbf{conditional, falsifiable} predictions that add no
fitted parameter but depend on the proximity-to-criticality hypothesis.
Because they are not required for --- and do not affect --- the core
background result, we develop them, with the
separate-universe/super-sample-covariance formalism
{[}arXiv:1711.07467{]} and the supporting figures, in a
\textbf{companion paper (in preparation)}; the corresponding kill test
is within reach of Euclid/Rubin weak-lensing tomography.

\textbf{Cross-horizon universality (speculative; companion).} Whether Δ
= 1 extends to black-hole horizons --- giving
\(S_{\mathrm{BH}} \propto  A^{3/2} \propto  M^{3}\) and strongly altered
primordial-black-hole evaporation --- is a speculative extension
logically separate from the core model, which requires the deformation
only in the dark-energy sector. We do \textbf{not} adopt it: SEDE is
consistent with the standard area law for astrophysical black holes, and
the black-hole-vs-cosmic-horizon asymmetry is a genuine
\textbf{conceptual cost} (§8.5), not a neutral prediction. We pursue
this, and the Verlinde dark-matter connection (§7), in the companion
paper.

\section{7. Relation to other sectors}\label{relation-to-other-sectors}

\textbf{Dark matter.} SEDE \textbf{does not derive dark matter}; it
assumes the standard cold component. What it \emph{does} provide are
falsifiable \emph{relationships}: the equation of state is locked to
\(\Omega _{\mathrm{m}}\) through the fixed-point background --- the
canonical model giving w₀ ≈ −1 with a crossing set by the growth gate
(§3.2; the closed-form w₀ ≈ −0.85 enters only as a structural
diagnostic, not the relation) --- and the coincidence problem is
reframed: dark energy's activation gate \(f_{\mathrm{sat}}(z)\) tracks
the (dark-matter-dominated) growth of structure, so ``why now''
corresponds to \(f_{\mathrm{sat}}\) crossing ≈ 0.5 near z ≈ 1, shortly
before matter--DE equality. These are testable links, not a unification,
and we state them as such.

There is, however, a \emph{principled} route to more, which we flag as a
future direction and pursue in a companion paper, not here. The same
volume-law dark-sector entropy SEDE uses for dark \emph{energy} is the
object from which Verlinde's emergent gravity (2016) derives apparent
dark \emph{matter} (baryonic displacement of the de Sitter volume-law
entropy → the MOND scale a₀ ∼ cH₀); if SEDE's Δ = 1 horizon entropy and
Verlinde's volume-law de Sitter entropy are the same object, \textbf{one
volume-law postulate could in principle source both dark energy and
apparent dark matter}. We do \textbf{not} adopt this here: Verlinde's
proposal is contested (it struggles with clusters and lensing and lacks
a covariant CMB/structure completion --- the regime SEDE \emph{is}
tested in), and our own cross-check
(\texttt{run\_verlinde\_crosscheck.py}) finds only an
\emph{order-of-magnitude} consistency --- the CKN/flatness normalisation
sets a₀ within an O(1) factor of observed
\(((c\surd \Omega _{\mathrm{DE0}}\) H₀)/2π ≈ 0.73 a₀, ∼25\% low), not a
derivation. We record it as a motivated, partially-consistent target for
the companion; SEDE's data fits use standard cold dark matter
throughout.

\textbf{Quantum and cross-field.} SEDE's horizon is fractal/volume-law
while its \emph{bulk} is standard GR (holographic-DE scope), so it
predicts a deformed horizon \textbf{without} a Lorentz-violating, foamy
bulk. Two real datasets are consistent with this restricted-scope
picture (they constrain bulk effects SEDE does not predict, so they do
not contradict it): the Fermi GRB 090510 31 GeV photon excludes
Planck-scale linear Lorentz violation \((E_{\mathrm{QG}} > several\)
\(M_{\mathrm{P}}\)), and the Fermilab Holometer's null result excludes
bulk Planck-amplitude ``holographic noise.'' The CKN reading gives a
gravitating-vacuum cutoff \(\Lambda _{\mathrm{UV}} \approx  4\) meV at
the dark-energy scale (§8.2). We tested one cross-field extrapolation
and it \textbf{failed}: the cosmological δ = 3/2 does \emph{not}
transfer to galaxy-cluster velocity statistics --- a real SDSS/Coma fit
gives a Gaussian (Tsallis q = 1.01, not the q ≈ 1.5 a naïve transfer
would predict). We report this negative result plainly; the
galaxy-dynamics link is the weakest cross-field claim and the data do
not support it.

\textbf{A contingent tension: cosmic birefringence.} Minimal SEDE
predicts \textbf{β = 0} by symmetry (its parity-even thermodynamic
scalar has no φFF̃ coupling --- the same structural fact as
\(\alpha _{\mathrm{T}} = 0\)), whereas a nonzero isotropic rotation is
reported (ACT DR6 0.215°±0.074°, 2.9σ {[}arXiv:2509.13654{]}). We treat
this as a \emph{real-but-contingent} threat, not a present
falsification, and assess it in the §9 red-team (threat 2): the signal's
cosmic, dark-energy-sourced nature is not established (miscalibration-
degenerate α+β; the cleanest β×galaxy channel is null at SEDE's zero),
but a confirmed cosmic β would favour a single-axion rival that does
SEDE's crossing and more.

\textbf{Where SEDE sits among dark-energy models --- and what it
borrows.} SEDE is a thermodynamic model, but it is not outside the
standard effective-field-theory description of dark energy, and locating
it there removes its main field-theoretic disadvantage. In the EFT of
dark energy (Gleyzes--Gubitosi--Piazza-- Vernizzi), SEDE occupies the
\textbf{safe corner}: the bulk is general relativity (holographic
scope), so the Planck-mass run rate
\textbf{\(\alpha _{\mathrm{M}} = 0\)} (Ġ/G = 0, §6/V7) and the tensor
speed excess \textbf{\(\alpha _{\mathrm{T}} = 0\)} --- gravitational
waves travel at c, so SEDE is \textbf{GW170817-safe by construction}
(\textbar{}\(c_{\mathrm{T}}/c - 1\)\textbar{} = 0, against the ∼10⁻¹⁵
bound), having never switched on the higher Horndeski operators that
GW170817 excludes. The dark energy is smooth with
\textbf{\(c_{\mathrm{s}}^{2} = 1\)} (§3.4), so it is gradient-stable
\((c_{\mathrm{s}}^{2} > 0)\) and does not cluster. SEDE is thus the
\((\alpha _{\mathrm{T}} = \alpha _{\mathrm{M}} = 0\),
\(c_{\mathrm{s}}^{2} = 1\)) point --- the \emph{least}-constrained
region of EFT-of-DE. \textbf{Field-theory representation and stability
through the crossing} (\texttt{run\_eft\_lagrangian.py}): a
\emph{single} minimally-coupled scalar cannot reproduce the w = −1
crossing --- for any P(X,φ), ρ + p = 2X \(P_{\mathrm{X}}\), so the
crossing forces \(P_{\mathrm{X}}\) → 0 (verified: 1 + w → 0 at z ≈
0.18), at which point the perturbation kinetic term
\(\alpha _{\mathrm{K}} \propto  P_{\mathrm{X}} + 2X\)
\(P_{\mathrm{XX}}\) passes through zero (a ghost) and
\(c_{\mathrm{s}}^{2}\) is ill-defined (the Vikman no-go). A
\emph{fundamental} ghost-free crossing nonetheless exists with
\textbf{braiding} (Kinetic Gravity Braiding / Horndeski G₃(φ,X)□φ): the
no-ghost determinant
\(Q_{\mathrm{S}} = \alpha _{\mathrm{K}} + (3/2)\alpha _{\mathrm{B}}^{2}\)
stays positive even when \(\alpha _{\mathrm{K}}\) → 0, because a nonzero
\(\alpha _{\mathrm{B}}\) carries it (we verify
\(Q_{\mathrm{S}} = 1.15 > 0\) at the crossing for a representative
\(\alpha _{\mathrm{B}} = 1.5\) \(\Omega _{\mathrm{DE}}\) realisation,
with \(c_{\mathrm{s}}^{2} > 0\) throughout) {[}Deffayet et al.~2010;
Kimura \& Yamamoto 2010{]}. SEDE, however, needs neither:
\(\rho _{\mathrm{DE}} = T_{\mathrm{AH}}\cdot s_{\mathrm{grav}}\) is
built from covariant horizon scalars, not a fundamental field, and is
treated as an effective smooth fluid (PPF, \(c_{\mathrm{s}}^{2} = 1\))
--- gradient-stable \((c_{\mathrm{s}}^{2} = 1 > 0)\), with no
propagating scalar to ghost, so the no-go simply does not apply (the
thermodynamic content, not a Lagrangian, carries the prediction).

The structure-coupling --- SEDE's most distinctive and least-confirmed
ingredient (the structure-sourced- dark-energy idea itself traces to
Gough's information dark energy, §7.6; what is specific here is the
horizon-thermodynamic realisation) --- is also \textbf{free of the two
instability classes that afflict coupled dark energy}
(\texttt{run\_eft\_stability.py}, V55). There is no gradient or ghost
instability because \(c_{\mathrm{s}}^{2} = 1 > 0\); and there is no
early-time large-scale instability because the interaction fraction
\(\Omega _{\mathrm{DE}}(z)\cdot\)\textbar d ln φ/d ln a\textbar{} → 0 at
high z (it peaks at ≈ 0.38 near z ≈ 0.5 and falls to ∼ 10⁻¹⁰ by
recombination, since the dark-energy fraction itself vanishes as
\(f_{\mathrm{sat}}\) → 0). The coupling switches \emph{off} in precisely
the radiation/early-matter era where coupled-DE instabilities are
seeded; the full perturbation evolution is validated end-to-end by
CLASS-PPF (fσ8 to 0.1\%, §4.4). So two of the apparent disadvantages
relative to other programmes in fact \emph{invert}: against Horndeski
SEDE is less constrained (GW-safe), and against interacting dark energy
it is more stable.

\begin{longtable}[]{@{}
  >{\raggedright\arraybackslash}p{(\linewidth - 4\tabcolsep) * \real{0.3333}}
  >{\raggedright\arraybackslash}p{(\linewidth - 4\tabcolsep) * \real{0.3333}}
  >{\raggedright\arraybackslash}p{(\linewidth - 4\tabcolsep) * \real{0.3333}}@{}}
\toprule\noalign{}
\begin{minipage}[b]{\linewidth}\raggedright
framework
\end{minipage} & \begin{minipage}[b]{\linewidth}\raggedright
what SEDE shares or borrows
\end{minipage} & \begin{minipage}[b]{\linewidth}\raggedright
net position
\end{minipage} \\
\midrule\noalign{}
\endhead
\bottomrule\noalign{}
\endlastfoot
ΛCDM (Planck) & empirical benchmark & \textbf{honest limit:} far less
established --- only DR3/Euclid closes this \\
quintessence (Ratra--Peebles) & dynamical DE; tracker/attractor crossing
(§8.1, Result 12--13) & structure-locked w(z) with a non-tuned −1
crossing \\
CKN bound & the UV--IR energy bound fixes the DE scale (§8.2) & SEDE
turns it into a cosmology; \textbf{limit:} needs the volume-law form on
top \\
Li holographic DE & holographic DE with a horizon cutoff & uses the
\textbf{apparent} horizon (no future-horizon causality/circularity);
\textbf{limit:} volume-law is radical \\
standard Barrow-HDE (Saridakis 2020) & the Barrow entropy
\(S \propto  A^{1+\Delta /2}\) & SEDE's Δ enters through the H-coupling
λ = 1−Δ/2 \emph{and} the structure gate, not the HDE density directly
--- so neighbouring Barrow-HDE fits that prefer Δ ≪ 1
{[}arXiv:2506.03019{]} constrain a \emph{different} model and do not
exclude SEDE's Δ = 1 (§9); \textbf{limit:} the volume-law endpoint is
the radical input \\
Hu--Sawicki f(R) & viable late-time acceleration & keeps GR
\((\alpha _{\mathrm{M}} = \alpha _{\mathrm{T}} = 0)\): no screening/GW
tension; \textbf{limit:} less mature \\
CPL & the data language; w(z) is CPL-degenerate in shape (§6) & adds the
growth--expansion lock CPL lacks; \textbf{limit:} model-dependent, not a
neutral fit \\
early dark energy & --- & \textbf{honest limit:} does not address H₀
(out of scope, \(r_{\mathrm{d}}\)-pinned, §9) \\
Jacobson (1995) & Clausius → Friedmann;
\(\rho _{\mathrm{DE}} = T\cdot s\) (App. A.1) & applies horizon
thermodynamics to DE; entropy-law change \textbf{confined to the DE
fluid} (§8.3), not gravity \\
Horndeski / EFT-of-DE & the α-function language & sits at
\(\alpha _{\mathrm{T}} = \alpha _{\mathrm{M}} = 0\),
\(c_{\mathrm{s}}^{2} = 1\) --- GW170817-safe corner; \textbf{less}
constrained, not less fundamental \\
Verlinde (2016) & volume-law de Sitter entropy; elastic/memory response
(§6, §8.1) & independent motivation for the postulate; a₀ ∼ cH₀
consistent at O(1) (§7) \\
interacting DE & the dark-sector energy-exchange ledger (CHR, §6) &
coupling fixed to growth, \(c_{\mathrm{s}}^{2} = 1\), off at high z ⟹
\textbf{instability-free} \\
\end{longtable}

The genuinely irreducible disadvantages, after these borrowings, reduce
to two, and we state them plainly: SEDE is \textbf{empirically younger}
than ΛCDM (the decisive tests are DR3/Euclid, §6), and it rests on the
\textbf{one volume-law postulate} (§8.4) that motivation --- Barrow, the
GSL, CKN, Verlinde --- supports but does not derive. Everything else in
the table is either shared, borrowed, or an honest matter of scope.

\textbf{Symmetry origin of the null predictions.} SEDE's many ``null''
results are not independent assumptions: they follow from a
\emph{single} structural condition --- \emph{the dark sector adds no
symmetry-breaking operator to the GR + Standard-Model effective action;
it is a parity-even scalar
\((\rho _{\mathrm{DE}} = T_{\mathrm{AH}}\cdot s_{\mathrm{grav}}\cdot f_{\mathrm{sat}})\)
minimally coupled through \(T_\mu \nu\).} Each precision null is then
forced by a specific symmetry:

\begin{longtable}[]{@{}
  >{\raggedright\arraybackslash}p{(\linewidth - 2\tabcolsep) * \real{0.5000}}
  >{\raggedright\arraybackslash}p{(\linewidth - 2\tabcolsep) * \real{0.5000}}@{}}
\toprule\noalign{}
\begin{minipage}[b]{\linewidth}\raggedright
null prediction
\end{minipage} & \begin{minipage}[b]{\linewidth}\raggedright
symmetry that forces it
\end{minipage} \\
\midrule\noalign{}
\endhead
\bottomrule\noalign{}
\endlastfoot
β = 0 (no cosmic birefringence) & parity in electromagnetism → no φ F F̃
coupling \\
\(\alpha _{\mathrm{T}} = 0\), \(c_{\mathrm{T}} = c\) (GW170817-safe) &
2-derivative diffeomorphism-invariant gravity → no Chern--Simons /
Galileon term \\
\(\alpha _{\mathrm{M}} = 0\), Ġ/G = 0, no fifth force, equivalence
principle & minimal (metric-only) coupling → no non-minimal ξφ²R \\
no Lorentz violation (Fermi GRB, Holometer) & bulk Lorentz invariance \\
\end{longtable}

So β = 0 and \(c_{\mathrm{T}} = c\) are \emph{the same fact} (§7
birefringence; absence of a parity-violating / modified-propagation
sector), and the whole cluster --- the reason SEDE passes every GW,
fifth-force, birefringence, and Lorentz null \emph{at once} --- is one
principle, not a list. The de Sitter endpoints add a fifth:
\(f_{\mathrm{sat}}\) → const (with H → const) ⟹ a maximally-symmetric de
Sitter phase ⟹ w = −1 by symmetry (the dark-energy bracket of Result 12;
App. A.7), with the crossing in between dynamical. We note the limit
sharply: this symmetry argument protects SEDE's \emph{safety}, not its
\emph{content}. The distinctive, load-bearing inputs --- the volume law
Δ = 1 (a geometric/thermodynamic postulate, §8.4), the coupling γ ≈ 1.5
(halo binding physics, §3.1), and the crossing redshift (a dynamical
balance, Result 13) --- are \textbf{not} symmetry-derivable; symmetry
buys SEDE nothing there, and the one open postulate is untouched. The
symmetry-protected cluster is precisely the part SEDE shares with any
minimally-coupled smooth dark energy; what makes it SEDE remains
thermodynamic and postulate-based.

\subsection{7.6 Relation to information dark energy and other
structure-sourced
models}\label{relation-to-information-dark-energy-and-other-structure-sourced-models}

The idea that dark energy is \emph{sourced by the entropy/information
produced as structure forms} is \textbf{not original to SEDE}, and we
state the priority plainly. It is due to \textbf{M. P. Gough}, in a
sustained but under-cited programme that we date --- from a NASA-ADS
author search --- to \textbf{2006--2008} (``On the Similarity of
Information Energy to Dark Energy,'' 2006; ``Information Equation of
State,'' 2008), developed through \emph{Holographic Dark Information
Energy} {[}\emph{Entropy} \textbf{13}, 924 (2011), arXiv:1105.4461{]} to
a recent observational analysis {[}\emph{Entropy} 2025{]} that fits the
stellar-mass-density history to the dark-energy equation of state from
Planck/Pantheon+/DES/DESI, finding
\(\rho _{\mathrm{DE}} \propto  SMD(a)^{p}\) with p ≈ 0.5 (R² = 0.93). In
Gough's model dark energy is the Landauer energy of the
information/entropy of star-heated gas and plasma. SEDE shares his
central physics --- structure's entropy \emph{is} dark energy, with a
self-regulating feedback (acceleration suppresses growth) that addresses
the coincidence problem and predicts time-limited acceleration --- and
we credit it as the priority for the structure-sourced-dark-energy idea.

What distinguishes SEDE is the \textbf{mechanism and a testable
consequence}, not the founding idea. (i) \emph{Mechanism:} Gough uses
the thermal/information entropy of hot \textbf{baryons} (∝ T, residing
in the bulk); SEDE uses the \textbf{gravitational} binding entropy of
bound structure deposited on the \textbf{apparent horizon}
\((\rho _{\mathrm{DE}} = T_{\mathrm{AH}}\cdot s_{\mathrm{grav}}\),
Barrow volume-law Δ=1). These are physically different entropies in
different places. (ii) \emph{Prediction:} the two are \textbf{not}
degenerate. SEDE's driver is the gravitational growth factor D(z) (dark
matter), Gough's is the baryonic stellar-mass-density SMD(z) (star
formation); computing both (\texttt{run} comparison,
\texttt{output/fig\_gough\_vs\_sede.png}) gives dark-energy histories
that agree to ≲ 6\% for z ≲ 0.6 --- where supernova data are strongest,
explaining why both fit current data --- but \textbf{diverge by up to
\textasciitilde50\% by z = 3}, SEDE staying flatter (more high-z dark
energy) because gravitational growth persists where star formation
declines. The coincidence that both land on a ``0.5'' exponent is
superficial: Gough's is a power of stellar mass density, SEDE's is the
H-exponent λ = 1−Δ/2 from Barrow Δ=1. (iii) SEDE adds a
\textbf{parameter-free} dark sector and a w-crossing rather than a
fitted CPL. So a clean way to separate them observationally is to ask
whether dark energy tracks \textbf{D(z)} (gravity) or \textbf{SMD(z)}
(stars), with the high-z dark-energy fraction (DESI Lyα, future surveys)
as the lever.

Other structure-/entropy-linked programmes differ more sharply and are
catalogued in \texttt{NOVELTY\_CHECK.md}: entropic-force and
Barrow/Tsallis holographic dark energy source the term from the
\textbf{horizon} (structure downstream --- and standard entropic-force
DE in fact \emph{struggles} with structure growth, which SEDE's design
inverts); backreaction (Buchert, Räsänen) and timescape (Wiltshire) make
acceleration an \textbf{apparent} inhomogeneity effect with no real dark
energy; matter-creation/CCDM (Lima) replaces dark energy with a
particle-creation pressure; and cosmologically-coupled black holes
(Farrah et al.) use one special structure via vacuum interiors. A
systematic priority pass over \textbf{both INSPIRE-HEP and NASA ADS}
(citation-ranked, multiple phrasings; documented in
\texttt{NOVELTY\_CHECK.md}) returns \textbf{no precedent} for SEDE's
specific construction --- ``dark energy = gravitational binding entropy
of bound structure on the apparent horizon''
(\texttt{abs:("dark\ energy"\ "binding\ energy"\ horizon)} yields zero
relevant hits in either database) --- nor for its terminology. So while
the \emph{general} structure-sourced-dark-energy idea belongs to Gough
(since \textasciitilde2006--2008), SEDE's \emph{gravitational-horizon
mechanism, its gravity-driven (D(z), not SMD) signature, and its
parameter-free first-principles dark sector} appear to be new. We also
ran Gough's model through our identical CAMB-marginalised pipeline (its
\(\rho _{\mathrm{DE}} \propto  SMD^{0.5}\) expressed as a w(a) table):
both are viable, with SEDE modestly preferred (ΔDIC ≈ −3.3 vs Gough's ≈
−0.45 relative to ΛCDM), and Gough's fit favouring a higher H₀ ---
consistent with his Hubble-tension claim. In summary, SEDE is a
distinct, more-derived realisation of Gough's idea, decided from ΛCDM
and from Gough alike only by the high-z dark-energy and growth data to
come (§6).

\section{8. The quantum-gravity
underpinning}\label{the-quantum-gravity-underpinning}

The energy-conservation account of §2--§3 took two horizon properties as
given: the \textbf{magnitude}, \(\rho _{\mathrm{DE}}\) \textasciitilde{}
\(\rho _{\mathrm{crit}}\) rather than the QFT vacuum's
\(M_{\mathrm{P}}^{4}\), and the \textbf{H-scaling} λ = 1/2, which
follows from the horizon entropy being \textbf{volume-law}
\((S_{\mathrm{grav}} \propto  V_{\mathrm{AH}}\); §3.3) with \emph{no
deformation parameter}. This section supplies their
quantum-gravitational basis --- it is the support beam under the
construction, and a reader content to take the two scales as inputs may
treat it as optional. We argue why the horizon entropy is volume-law
rather than area-law, anchor the magnitude to the holographic bound and
the flatness normalisation, and isolate the single statement that
remains a postulate. Throughout, the Barrow parameter Δ appears only to
\emph{label} the area-to-volume interpolation
\((S \propto  A^{1+\Delta /2}\): Δ=0 area-law, Δ=1 ⟺ volume-law) and as
the \textbf{falsifiable test} of the volume postulate (measured in §5.6,
forecast in §6); it is not an input to the model.

\subsection{8.1 Why volume-law, not area-law --- consistency motivations
(none a
derivation)}\label{why-volume-law-not-area-law-consistency-motivations-none-a-derivation}

The dark-sector entropy is the coarse-grained entanglement entropy
contained in the apparent-horizon \textbf{volume},
\(S_{\mathrm{grav}} \propto  V_{\mathrm{AH}} \propto  R_{\mathrm{AH}}^{3}\)
--- equivalently the Barrow endpoint Δ=1
\((A^{1+\Delta /2} = A^{3/2} = R^{3})\). The open physical question is
sharp: why the volume law, rather than the Bekenstein--Hawking
\emph{area} law? We give three considerations in the main line and
collect three further, more speculative ones in Appendix G. \textbf{None
is a microscopic proof}; together they motivate the volume endpoint and
make the area-law branch disfavoured in the diagnostics.

\begin{enumerate}
\def\labelenumi{\arabic{enumi}.}
\tightlist
\item
  \textbf{Geometric ceiling caps the deformation} (Result 9). The
  horizon's fractal dimension is \(d_{\mathrm{H}} = 2 + \Delta\), capped
  at the space-filling value \(d_{\mathrm{H}} = 3\) for a two-surface in
  three-space. Since (for a volume law in embedding dimension d) Δ =
  2/(d−1), the case d = 3 gives Δ = 1. We are precise about what this
  delivers: the ceiling bounds the deformation, \textbf{Δ ≤ 1} --- it
  does \emph{not} select Δ = 1. Landing \emph{on} the endpoint requires,
  in addition, that the entropy is monotone in Δ (so the cap is
  saturated); since \(\rho _{\mathrm{DE}} = T\cdot s\) makes the free
  energy F = E − TS vanish identically, entropy is the governing
  potential and the generalised second law is consistent with --- is a
  \textbf{fixed point of} --- the saturated (volume-law) state. We
  stress that this is a \textbf{consistency check, not a selection
  principle}: ``total entropy does not decrease'' is not ``the parameter
  takes its entropy-maximising value,'' and the saturation of the cap
  \emph{is} the volume-law postulate (§8.4) in effective form, not an
  independent derivation of it. Volume-law scaling is the signature of a
  thermalised, maximally-scrambled state (Page curve, eigenstate
  thermalisation).
\item
  \textbf{The holographic bound fixes the scale, not the form.} The CKN
  energy bound (§8.2) pins the \emph{magnitude} but not the exponent ---
  its own state count gives a \emph{sub}-area-law (δ = 3/4), the
  opposite of an enhancement. The volume law is therefore a statement
  about the horizon's quantum \emph{state} (typical / thermalised),
  which the energy bound cannot supply. This isolates the open content
  to one sentence: \emph{the horizon is in a typical,
  volume-law-entangled state.}
\item
  \textbf{The area-law branch is disfavoured directly by current data.}
  \emph{Saturating} the CKN bound --- the natural ``no-deformation''
  alternative --- is algebraically the area-law case Δ = 0, whose
  equation of state (w₀ = −0.49) and large early-dark-energy fraction
  are disfavoured in the conditional diagnostics of §5.6. On current
  data the volume branch is \textbf{favoured over the area-law
  alternative}, not merely an optional choice --- though, as §5.6
  stresses, the robust marginalised measurement is still to come.
\end{enumerate}

Three further consistency motivations --- a \emph{driven non-equilibrium
steady state} (the structure gate maintains the volume-law state despite
a long scrambling time), a \emph{roughening / Edwards--Wilkinson
universality-class} analogy for the horizon surface, and the
\emph{independent volume-law entanglement entropy} of Verlinde's
emergent-gravity programme --- are collected in \textbf{Appendix G}.
Each makes the postulate less \emph{ad hoc}; none derives it.

\subsection{\texorpdfstring{8.2 Why \(\rho _{\mathrm{DE}}\)
\textasciitilde{} \(\rho _{\mathrm{crit}}\) --- the CKN bound and the
vacuum-energy
scale}{8.2 Why \textbackslash rho \_\{\textbackslash mathrm\{DE\}\} \textasciitilde{} \textbackslash rho \_\{\textbackslash mathrm\{crit\}\} --- the CKN bound and the vacuum-energy scale}}\label{why-rho-_mathrmde-rho-_mathrmcrit-the-ckn-bound-and-the-vacuum-energy-scale}

The deformation fixes the \emph{scaling} of dark energy with H but not
its \emph{magnitude}. The magnitude comes from the Cohen--Kaplan--Nelson
(CKN) holographic energy bound (Principle 0): an effective field theory
in a region of size L cannot hold more energy than a black hole of that
size, ρ \(L^{3} \lesssim  M_{\mathrm{P}}^{2}\) L,
i.e.~\(\rho  \lesssim  M_{\mathrm{P}}^{2}/L^{2}\). With the IR cutoff at
the cosmic horizon, L = c/H, this gives
\[\rho_{\rm DE} \lesssim M_P^2 H^2 \sim \rho_{\rm crit} \;\Rightarrow\; \Omega_{\rm DE} \sim \mathcal{O}(1).\]
The naïve QFT vacuum, ρ \textasciitilde{} \(M_{\mathrm{P}}^{4}\),
overshoots this by \(\rho _{\mathrm{P}}/\rho _{\mathrm{crit}}\)
\textasciitilde{} 10¹²² --- the cosmological-constant problem. SEDE
never invokes the QFT vacuum: dark energy is the horizon's conjugate
thermal energy density, and the CKN relation supplies the
\textbf{horizon-scale upper bound}
\(\rho _{\mathrm{DE}} \lesssim  M_{\mathrm{P}}^{2}H^{2}\). This is a
bound, not an equality: the present amplitude is fixed by spatial
flatness \((\Omega _{\mathrm{DE0}} = 1 - \Omega _{\mathrm{m}})\), in the
same bookkeeping sense that ΛCDM fixes \(\Omega _\Lambda 0\) --- once
the volume-law branch is chosen, the bound is saturated to O(1) by that
normalisation \textbf{today}. We note explicitly that the saturation is
a present-day statement: since \(\rho _{\mathrm{DE}} \propto  H\)
(volume law) while the bound ∝ H², the ratio
\(\rho _{\mathrm{DE}}/(M_{\mathrm{P}}^{2}H^{2}) = \Omega _{\mathrm{DE}}(z)\)
falls at higher z (0.70 → 0.22 → 0.03 at z = 0, 1, 3), so dark energy
sits \emph{below} the CKN bound in the past --- consistent, as CKN is an
upper bound. So CKN \emph{addresses} the vacuum-energy-scale problem (it
explains why the natural gravitational horizon scale is
\(M_{\mathrm{P}}^{2}H^{2}\) \textasciitilde{} \(\rho _{\mathrm{crit}}\),
not \(M_{\mathrm{P}}^{4}\)) rather than deriving the magnitude from
first principles. The same bound, read as a UV--IR relation, gives a
gravitating-vacuum cutoff \(\Lambda _{\mathrm{UV}} = (M_{\mathrm{P}}\)
\(H_{0})^{1/2} \approx  4\) meV --- numerically the observed dark-energy
scale (§7).

\subsection{8.3 The seam, BBN-safety, and the holographic
scope}\label{the-seam-bbn-safety-and-the-holographic-scope}

A naïve \emph{microscopic} volume-law entropy density would be
Planckian, so a naïve
\(\rho _{\mathrm{DE}} = T\cdot (S_{\mathrm{vol}}/V)\) would overshoot
\(\rho _{\mathrm{crit}}\) by \((M_{\mathrm{P}}/H)^{\Delta }\)
\textasciitilde{} 10⁶¹ --- the square root of the 10¹²² hierarchy
(Result 11). The \emph{same} factor controls three apparently separate
issues --- this seam, the BBN bound on a Barrow-modified Friedmann
equation, and the modified-gravity-vs-holographic fork --- which are one
object. The resolution is the holographic-DE scope adopted in §2: the
deformation acts \textbf{only on the dark-energy fluid} (magnitude from
CKN, scaling from Barrow), \textbf{not} on the gravitational sector. The
early-universe expansion is then exactly standard
\((\Delta N_{\mathrm{eff}} = \Delta Y_{\mathrm{p}} = 0\), verified;
BBN-safe), and the naïve T·(S/V) is precisely the modified-gravity path
the scope forbids.

This is the \textbf{load-bearing assumption of the model}. In the
\emph{modified-gravity} reading --- Barrow entropy deforming the
Friedmann equations {[}arXiv:2110.00059{]} --- BBN forces \textbf{Δ ≲
1.4×10⁻⁴} {[}Barrow, Basilakos \& Saridakis, arXiv:2010.00986{]} and
SN+BAO+H(z) give Δ \textasciitilde{} 10⁻⁴ (the area law at 2σ
{[}arXiv:2108.10998{]}), which read literally would exclude SEDE's Δ =
1. Those bounds do \emph{not} apply to SEDE: it keeps standard gravity
and lets the volume-law entropy act only on the dark-energy fluid via
\(\rho _{\mathrm{DE}} = T\cdot s\) (§2, App. A.1), so there is no extra
Friedmann term, no BBN deviation, and Δ is unconstrained by them. The
headline Δ = 1 thus rests on this single scope choice --- legitimate and
standard (it \emph{is} the holographic-dark-energy framework), but not
free: a reader who rejects it recovers the tightly-bounded
modified-gravity Barrow cosmology instead.

\subsection{8.4 The one irreducible
postulate}\label{the-one-irreducible-postulate}

\begin{quote}
\textbf{Postulate (volume-law horizon).} The cosmic apparent horizon
carries the coarse-grained \emph{entanglement} entropy of its
\textbf{volume},
\(S_{\mathrm{grav}} \propto  V_{\mathrm{AH}} \propto  R^{3}\) --- a
constant horizon entropy density, i.e.~a thermalised, volume-law
(non-extensive) state --- rather than the Bekenstein--Hawking area law.
\end{quote}

The results in §8.1--§8.3 are bounds, consistency arguments, or
motivations \emph{given} this postulate; the postulate itself is not
derived from a microscopic theory of quantum gravity. §8.1
\emph{reduces} it (the CKN bound fixes the scale, so the only open
content is the horizon's entanglement \emph{state}), \emph{motivates} it
(non-extensivity, the GSL, the driven steady state, the
Edwards--Wilkinson universality class, and the independent volume-law
entanglement entropy of Verlinde's emergent gravity), and shows the
\emph{area-law branch is disfavoured by current diagnostics} --- but a
microscopic state count that delivers the volume law remains the model's
single open foundational item (§9). Two points of emphasis. First, the
model carries \textbf{no fitted deformation parameter}: Δ is not an
input but the falsifiable \emph{test} of this postulate --- current data
already prefer the volume endpoint and disfavour the area-law branch
(§5.6, a conditional diagnostic), with the robust marginalised
measurement to come from DESI DR3 + Euclid (§6), alongside cross-horizon
predictions (a deformed black-hole entropy \(S \propto  A^{3/2}\); §7,
presented as speculative). Second, the open question is now stated
plainly --- \emph{why is the horizon entropy volume-law (coarse-grained
entanglement) rather than area-law?} --- rather than hidden in a value
of a tunable parameter. We do not claim to have derived quantum gravity;
we claim to have reduced its dark-energy content to one testable
statement and removed it as a fitted parameter.

\textbf{On the plausibility of a volume law (genre, not derivation).} A
volume-law horizon entropy is a strong hypothesis, but volume-law
\emph{entanglement} entropy is at least not exotic in cosmological field
theory. The entanglement entropy of a field subregion --- area-law in
the de Sitter adiabatic vacuum --- grows and saturates to a \emph{volume
law} after a de Sitter→flat transition, ending in a partially
thermalised (generalised-Gibbs) state {[}Khlebnikov \& Sheoran,
arXiv:1907.00487{]}; volume-law entanglement also arises from
field--curvature coupling {[}arXiv:2306.08357{]}. We cite these only to
establish \emph{genre and plausibility} --- that volume-law entropy is a
real feature of cosmological QFT --- and we explicitly resist
over-reading them. Two caveats: (i) these results concern the
entanglement entropy of field \emph{subregions} under specific
transitions, not the cosmic-\emph{horizon} entropy SEDE invokes; and
(ii) their assignment (de Sitter = area-law, flat = volume-law) is
\emph{not} SEDE's and in fact runs opposite to it (SEDE places the
volume law in the de Sitter brackets). So the literature makes a
volume-law hypothesis less exotic, but it neither \emph{anticipates} nor
\emph{derives} the specific postulate --- which remains the model's one
open foundational item (§9). (We note for completeness that
Padmanabhan's holographic-equipartition route,
\(dV/dt = N_{\mathrm{sur}} - N_{\mathrm{bulk}}\), does \emph{not} help
here: its surface count is area-law,
\(N_{\mathrm{sur}} = A/\ell _{\mathrm{P}}^{2}\), and it reproduces the
\emph{standard} area-law Friedmann equation --- it is the baseline SEDE
departs from, not independent support for the volume law.)

\subsection{\texorpdfstring{8.5 The postulate is a dof-\emph{counting}
claim, not a thermalisation
claim}{8.5 The postulate is a dof-counting claim, not a thermalisation claim}}\label{the-postulate-is-a-dof-counting-claim-not-a-thermalisation-claim}

The phrase ``volume-law'' conflates two claims of very different status.
\textbf{(i) The state:} are the horizon degrees of freedom maximally
entangled (thermal) or in a low-entanglement ground state? \textbf{(ii)
The counting:} does the \emph{number} of horizon dof grow as the area,
\(N \propto  R^{d-2}\), or the spatial volume, \(N \propto  R^{d-1}\)?
Since S ∼ (state factor) × N, the deformation Δ --- through
\(S \propto  A^{1+\Delta /2}\)∝ \(R^{(d-2)(1+\Delta /2)}\) --- is fixed
by the \textbf{counting} (area→Δ=0, volume→Δ=1 in d=4); the state only
sets whether the prefactor is maximal. SEDE's postulate is, precisely,
the \textbf{counting} claim.

\emph{The state half is settled, and is not peculiar to SEDE.} Maximal
scrambling has a precise meaning in the Sachdev--Ye--Kitaev model
{[}Sachdev \& Ye 1993; Kitaev 2015{]}, the canonical chaotic,
black-hole-dual system. But maximal scrambling \textbf{does not fix Δ}:
a black hole is a \emph{maximal} scrambler yet \textbf{area-law} (S =
A/4, Δ=0, reproduced by string microstate counts). Scrambling is
therefore shared with the Δ=0 black hole; and SYK, being all-to-all
(geometry-free), cannot even pose the counting question. (An explicit
Majorana-SYK diagonalisation confirming the state is maximally entangled
--- Wigner--Dyson statistics, Page-value saturation, complete OTOC
scrambling --- is given in Appendix G; it settles the \emph{state}, not
the counting.)

\emph{The counting half is where SEDE lives, and area is the default.}
In causal-set quantum gravity --- the setting where the counting
question is natural --- both scalings appear, but the \emph{canonical}
horizon-entropy object, the causal-\textbf{link} count that recovers the
Bekenstein--Hawking law {[}Sorkin and collaborators{]}, is the
\textbf{area} one \((R^{d-2})\), in flat space and identically in de
Sitter. So even in the framework most hospitable to volume-counting,
\textbf{area-law is the default}, and SEDE's Δ=1 is the specific,
non-default identification of horizon entropy with the \textbf{bulk}
count. (The causal-set sprinkle measurements, the entropy-bound check,
and a Ryu--Takayanagi holographic-code realisation --- volume-law
requires \emph{nonlocal} horizon connectivity --- are in Appendix G;
their verdict is that the volume count is \textbf{allowed and realisable
but not selected}: no consistency argument available to us derives it,
and none excludes it.)

\textbf{The sharpened postulate, properly located.} The one open item is
therefore \emph{not} about thermalisation or chaos (settled, and shared
with area-law black holes) but is purely a \textbf{dof-counting}
statement:

\begin{quote}
the cosmic horizon's entropy counts its \textbf{bulk} (volume-scaling, N
∝ R³ ⟺ \(d_{\mathrm{H}} = 3\)) degrees of freedom rather than its
\textbf{boundary} (area-scaling, links) ones --- a count that exceeds
the holographic area bound by \(R/\ell _{\mathrm{P}}\) (the seam of
§8.3).
\end{quote}

It reduces to two deciders. \emph{Theoretically} it is precisely
\emph{``is the Hilbert-space dimension of the de Sitter static patch
\(e^{Area/4}\) or \(e^{Volume}\)?''} --- a frontier problem a complete
dS holography would settle. \emph{Empirically} it is the deformation Δ:
Δ=0 (area) is already disfavoured by present data (§5.6), and DESI DR3 +
Euclid will measure Δ to \textasciitilde0.09 (§6). SEDE's postulate is
thus not an arbitrary assumption but a \emph{consistent, realisable bet
on a well-posed open question of quantum gravity}, with a decisive test
imminent. Tellingly, the one horizon we can check today --- the black
hole --- is area-law; we name that asymmetry as a genuine conceptual
cost (§7), not a neutral prediction.

\section{9. Discussion: open problems \& honest
assessment}\label{discussion-open-problems-honest-assessment}

\textbf{The one open item.} Every dark-sector input is fixed by a stated
choice (§3) given the single postulate that the horizon is
volume-law-entangled (§8.4). That postulate is motivated by gravity's
non-extensivity and the second law but is not derived from a microscopic
state count; it is the model's one irreducible input. We have argued it
is not closable by reframing (canonical vs microcanonical accounting
gives either area-law λ = 1 or Planckian-volume λ = 1/2, never both for
free) --- it is a genuine open problem, and a \emph{falsifiable} one
(§6).

\textbf{Scope and tensions.} H₀ is out of SEDE's scope (it is
\(r_{\mathrm{d}}\)-pinned; the model does not address the
distance-ladder tension and we do not claim it does). The
structure-sourcing \emph{mechanism} itself --- the defining SEDE claim
that \(f_{\mathrm{sat}}\) tracks growth --- is only weakly confirmed: a
geometry-vs-structure consistency test (E1; the null test
\(\rho _{\mathrm{DE}} = G[D(z)]\) with \(f_{\mathrm{sat}}\) read
independently off the \emph{expansion}, DESI
\(D_{\mathrm{H}}/r_{\mathrm{d}}\) → H(z), and off \emph{structure
growth}, RSD fσ8 → σ8(z)) tracks at z \textless{} 1 but is RSD-limited
at z \textgreater{} 1 and remains inconclusive (current data: χ²/dof =
0.40 in the RSD-error budget, 1.40 in the precise-geometry budget;
Figure 9). \textbf{The pipeline is implemented and staged}
(\texttt{run\_e1\_mechanism.py}) --- it runs now and is ready to apply
the moment Euclid/LSST weak-lensing tomography sharpens the structure
leg, at which point the forecast of §6 makes it decisive (∼4σ). Its data
vector and CONFIRM/REFUTE thresholds are \textbf{pre-registered}
(PREREGISTRATION.md §F) so the staged test cannot be tuned after the
data arrive. We have separated this mechanism into a derived,
parameter-free part --- the growth--expansion lock P2, which follows
from \(f_{\mathrm{sat}} = f(D^{2})\) with no extra assumption (§6),
testable with exactly that data --- and an optional near-critical layer
(deferred to a companion paper) that supplies a further kill test. But a
\emph{derived consistency relation} is not yet a \emph{confirmation},
and the mechanism stays unconfirmed until the lock P2 is measured. We
stress this is a \emph{near-term} gap, not a permanent one: a Fisher
forecast (§6, \texttt{run\_mechanism\_forecast.py}) shows the
growth-side mechanism test reaches ∼4σ with combined Euclid/LSST weak
lensing + CMB-lensing + clusters + RSD --- so the distinctive claim is
decidable on the same timescale as the amplitude Δ, not deferred
indefinitely. And the cosmic- birefringence signal (§7) is a live
tension. We emphasise the honest framing: SEDE's \emph{phenomenology} (w
crossing −1, the volume-law endpoint) is mildly favoured in the current
compiled-data analysis with an encouraging out-of-sample check, but its
\emph{distinctive mechanism} (the \(f_{\mathrm{sat}}\) gate tracking
structure growth) is the unconfirmed frontier.

\textbf{Threats from the recent literature (a red-team).} We list the
strongest recent results that could damage SEDE, with our honest
assessment of each.

\begin{enumerate}
\def\labelenumi{\arabic{enumi}.}
\tightlist
\item
  \emph{Holographic DE disfavoured by full CMB / the early--late
  tension.} Wu et al.~{[}arXiv:2509.02945{]} find that holographic and
  Ricci DE are disfavoured by full ACT DR6 + DESI through an
  early-vs-late tension that compressed (R, ℓ\_A) priors hide --- a
  direct challenge, since SEDE's headline ΔDIC used compressed CMB. We
  tested it. \textbf{Marginalised full primary likelihood:} a
  CAMB-in-the-loop fit on the full Planck \(plik_{\mathrm{lite}}\)
  TTTEEE + lowℓ TT/EE, with the five cosmological parameters and the
  Planck nuisance \emph{all free} (not θ*-fixed), gives
  \textbf{Δχ²(SEDE−ΛCDM) = −0.43} (\texttt{run\_full\_cmb\_mcmc.py}),
  the preference sitting in the high-ℓ peak structure (−0.37) ---
  \emph{negative}, i.e.~mildly favouring SEDE, \textbf{not} the large
  \emph{positive} Δχ² the HDE early--late tension would produce; the
  best-fits are physical \((\Omega _{\mathrm{m}} = 0.295/0.315\) for
  SEDE/ΛCDM). This is corroborated by two θ*-matched single-point
  proxies --- the same \(plik_{\mathrm{lite}}\)
  \((\Delta \chi ^{2}_{\mathrm{CMB}} = -0.50\),
  \texttt{run\_plik\_check.py}) and the official ACT DR6 CMB-lensing
  likelihood \((\Delta \chi ^{2}_{\mathrm{ACT}}\)-lens = −0.32,
  \texttt{run\_act\_lensing.py}, validated against the package's
  fiducial χ² = 14.06). ACT DR6 lensing re-evaluated \textbf{at the
  marginalised primary best-fit} (not a fixed fiducial) gives
  \(\Delta \chi ^{2}_{\mathrm{ACT}}\)-lens = −0.31 (ΛCDM 14.10 → SEDE
  13.79), and the low-ℓ ISW-sensitive channel is flat (−0.03). All are
  \textbar Δχ²\textbar{} \textless{} 1: SEDE does \textbf{not} inherit
  the standard-HDE tension, because it has \emph{less} early dark energy
  than ΛCDM \((f_{\mathrm{sat}}\) → 0,
  \(\Omega _{\mathrm{DE}}(z=3) \approx  0.035 < 0.040\)), not the
  runaway high-z behaviour that breaks original HDE. \textbf{Full
  non-lite plik (foregrounds sampled):} we have since run the
  \emph{complete} high-ℓ plik TTTEEE likelihood with all 15
  foreground/calibration nuisance parameters sampled at their Planck
  priors (\texttt{run\_full\_cmb\_mcmc.py\ -\/-fullplik}), the
  gold-standard primary CMB --- it gives \textbf{Δχ²(SEDE−ΛCDM) = −0.33}
  (TTTEEE −0.27, lowl.TT −0.23, lowl.EE +0.17), confirming the
  foreground-marginalised \(plik_{\mathrm{lite}}\) result (−0.43) to
  within the minimiser noise. So \textbf{letting the foregrounds float
  does not reveal a hidden tension} --- the standard-HDE early--late
  wall is simply absent for SEDE. \emph{Residual scope:} only the ACT
  DR6 \emph{primary} multifrequency likelihood remains not run (ACT
  \emph{lensing}, the DE-sensitive channel, is; §above) ---
  \textbf{formally deferred} under the pre-registered trigger
  (PREREGISTRATION.md §E; \texttt{act\_dr6\_mflike} + multi-GB products
  not installed), with the falsifiable committed expectation
  \textbar Δχ²\textbar{} \textless{} 1 given three concordant CMB
  datasets already at that level; a full MCMC for ΔDIC is deferred
  (expected ≈ Δχ², \(p_{\mathrm{D}}\) shared). The full primary
  likelihood is also folded \emph{into the joint} (§5.1).
\item
  \emph{Cosmic birefringence --- a contingent threat, with a rival
  waiting if it is confirmed.} SEDE predicts β = 0; a nonzero
  \emph{rotation} is reported (ACT DR6 2.9σ {[}arXiv:2509.13654{]};
  Planck \textasciitilde0.3°), and \emph{if} it is cosmic and
  dark-energy-sourced, a single axion explains β + acceleration
  {[}arXiv:2504.17638{]} and an interacting axion sector reproduces the
  phantom crossing + β + dark matter {[}arXiv:2506.12589{]} --- a
  parsimony threat, since the rival does SEDE's crossing and more with
  one field. But the cosmic, DE- sourced nature is \emph{not
  established}: the measurement is of α + β (miscalibration-degenerate),
  Planck and ACT disagree by \textgreater3σ, SPIDER is incompatible with
  the coupling {[}arXiv:2510.25489{]}, and the β × LSS cross-correlation
  is null (p ≈ 37\%, {[}arXiv:2509.22273{]}) --- at SEDE's predicted
  zero, not the axion (§7). So this is a \emph{contingent} threat, not a
  present falsification, and the cleanest current channels favour β = 0.
  The discriminator is the pair \{w-crossing, β\}; β = 0 is itself
  symmetry-derived (§7), tied to GW-safety, so a confirmed cosmic β
  would also predict GW parity violation.
\item
  \emph{The DDE signal may be a low-z SN systematic.} Huang, Cai \& Wang
  {[}arXiv:2502.04212{]} show the DESI evolving-DE preference is biased
  by an external low-z SN sample in DESY5 and drops to \textless2σ once
  corrected. \emph{Defence (partial):} the bias is DESY5-specific (``in
  contrast to the uniform behaviour of PantheonPlus''), and SEDE's
  headline uses Pantheon+; more decisively, SEDE's leave-one-probe-out
  \textbf{drops supernovae entirely and stays preferred (Δχ² = −2.47)},
  and \(F_{\mathrm{AP}} / tracer\)-LOO are SN-free. SEDE is not
  SN-driven --- but if the field's evolving-DE evidence softens, SEDE's
  \emph{context} erodes, and its ``suggestive, not established'' framing
  must hold.
\item
  \emph{The nearest Barrow-HDE fit prefers Δ ≪ 1.} Luciano,
  Paliathanasis \& Saridakis {[}arXiv:2506.03019{]} fit standard
  Barrow-HDE to DESI DR2 + SN + CC and obtain \textbf{Δ \textless{}
  0.47--0.54} (Δ = 0 within 1σ, mild positive tendency) --- which, read
  literally, disfavours SEDE's Δ = 1. \emph{Defence (shown, not
  asserted; Fig. 10, \texttt{run\_delta\_orthogonality.py}):} both
  models share the scaling
  \(\rho _{\mathrm{DE}} \propto  H^{2-\Delta }\); the \emph{only}
  difference is SEDE's structure gate \(f_{\mathrm{sat}}(z)\), and the
  gate is exactly what decouples Δ from the observable the Barrow-HDE
  bound rests on. At the \textbf{same} Δ = 1 the two models differ
  sharply: SEDE's gate \((f_{\mathrm{sat}}\) → 0 at high z) holds the
  early-dark-energy fraction to \(\Omega _{\mathrm{DE}}(z=3) = 0.029\)
  --- CMB-safe --- whereas un-gated Barrow-HDE has
  \(\Omega _{\mathrm{DE}}(z=3) = 0.147\) (5× larger) and w₀ = −0.77 with
  no w = −1 crossing, versus SEDE's w₀ = −1.0 in DESI's crossing
  quadrant. So the Δ-constraint derived from \emph{Barrow-HDE's}
  early-DE/EOS does not transfer: SEDE's Δ = 1 produces viable
  observables precisely where Barrow-HDE's Δ = 1 does not, and in SEDE's
  own construction Δ = 1 is data-preferred over Δ = 0 (§5.6). The bound
  constrains a sibling model, not SEDE --- and Fig. 10 shows the
  orthogonality rather than asserting it.
\end{enumerate}

\begin{figure}
\centering
\pandocbounded{\includegraphics[keepaspectratio]{../../output/fig_delta_orthogonality}}
\caption{Fig10}
\end{figure}

\textbf{Figure 10.} Why the standard Barrow-HDE bound Δ ≪ 1 does not
constrain SEDE's Δ = 1. Both models share the holographic scaling
\(\rho _{\mathrm{DE}} \propto  H^{2-\Delta }\); the only difference is
SEDE's structure gate \(f_{\mathrm{sat}}(z)\). \emph{Left:} the
early-dark-energy fraction \(\Omega _{\mathrm{DE}}(z=3)\) --- the
quantity the CMB and the standard-HDE EOS fit actually constrain ---
versus Δ. Un-gated Barrow-HDE (red) carries large, strongly Δ-dependent
early DE \((\Omega _{\mathrm{DE}}(z=3) = 0.15\) even at Δ = 1, rising to
0.70 at Δ = 0), so its fits pin Δ \textless{} 0.5; SEDE's gate (blue)
drives \(f_{\mathrm{sat}}\) → 0 at high z, holding
\(\Omega _{\mathrm{DE}}(z=3) \approx  0.03\) --- at or below the
ΛCDM/CMB-safe level (grey dashed) --- for \emph{all} Δ. At Δ = 1 the two
models differ five-fold (0.029 vs 0.147), and SEDE additionally crosses
w = −1 (w₀ = −1.0) where un-gated Barrow-HDE does not (w₀ = −0.77).
\emph{Right:} consequently the data prefer \textbf{orthogonal} Δ in the
two parametrisations --- Δ̂ \textless{} 0.5 for Barrow-HDE (EOS/early-DE;
Luciano+ 2025) and Δ̂ ≈ 1.0 for SEDE (amplitude/early-DE; §5.6) --- with
no contradiction. The gate is the difference.
(\texttt{run\_delta\_orthogonality.py}.)

Net: the two threats that targeted the model itself (full-CMB tension;
the Barrow-Δ bound) are answered by direct tests or a stated model
distinction; the two that target the \emph{context} (birefringence; the
low-z SN systematic) are real and partially outside SEDE's control ---
birefringence especially.

\textbf{Evidential status.} SEDE is a falsifiable, fixed-parameter,
Λ-free alternative to ΛCDM, mildly preferred (suggestive
\textasciitilde2--3σ) with an encouraging out-of-sample check, resting
on one open postulate and a few stated theoretical choices. It is
stronger than ``indistinguishable from ΛCDM'' and weaker than
``established.'' Its value is twofold: a \emph{physical story} (no
fine-tuned Λ, w(z) from \(\Omega _{\mathrm{m}}\), the coincidence
reframed) and a \emph{decidable prediction} (the volume law, Δ = 1). The
data, not the theory, will settle it.

\section{10. Conclusions}\label{conclusions}

We have presented Structural Entropy Dark Energy: a fixed-parameter
dark-energy model in which dark energy is identified with the conjugate
thermal energy density of a volume-law-entangled cosmic horizon and
activated through a structure-growth gate. The model is a single
Friedmann fixed-point equation with \textbf{no fitted dark-sector
parameter} --- λ, w(z), γ, and \(c_{\mathrm{s}}^{2}\) are each fixed by
a stated theoretical choice rather than tuned --- that replaces the
cosmological constant by a structure-gated, horizon-coupled term. It
\textbf{addresses} the vacuum-energy-scale problem (the scale is the CKN
holographic bound, normalised by flatness, not the QFT vacuum) and
reframes the coincidence problem (the activation gate tracks the
late-time maturation of the structure-growth history). In a
marginalised, CAMB-in-the-loop joint analysis it is mildly preferred
over ΛCDM (ΔDIC ≈ −2.9); the preference is suggestive (calibration
\textasciitilde2--3σ, not obviously flexibility) with an encouraging
pre-DESI → DESI out-of-sample check, and its equation of state crosses
−1 in DESI's evolving-DE quadrant. Its quantum-gravitational content
reduces to one falsifiable postulate --- a volume-law horizon entropy.
The sharp future test is near: under the stated Fisher assumptions, DESI
DR3 and Euclid could constrain the effective deformation Δ to σ ≈ 0.09,
separating SEDE's volume-law horizon from a smooth one at
\textasciitilde11σ. SEDE is suggestive, not established --- and,
crucially, it is decidable.

\section{Code and data availability}\label{code-and-data-availability}

All code, data vectors, and analysis pipelines are publicly available at
\texttt{github.com/spsingularity/SEDE}. A single entry point,
\texttt{reproduce\_all.py}, regenerates every number and figure in this
paper from staged drivers; the model and likelihood live in the
\texttt{sede/} package, and the verification suite
(\texttt{sede/verification.py}) checks the quoted results numerically.
The observational inputs are standard public datasets (DESI DR2,
Pantheon+/DES-5YR/Union3, Moresco cosmic chronometers, Gold-2018 fσ8,
compressed Planck, eBOSS DR16), retrieved via the documented loaders. A
tagged release with an archival DOI will accompany journal submission.

\section{Use of AI tools}\label{use-of-ai-tools}

Artificial-intelligence tools (large language models) were used to
assist with drafting and editing the manuscript, with developing and
cross-checking the analysis code, and with literature searches. The
author conceived the model, directed and independently verified all
analyses, and takes full responsibility for the content, including all
derivations, results, and conclusions. No AI tool is an author.

\section{References}\label{references}

\emph{Cited inline by author and year. All 45 arXiv identifiers were
verified to resolve on the live arXiv API, with the returned titles
checked against the cited content (\texttt{verify\_refs.py}); the 19
post-2024 identifiers --- the highest hallucination risk --- all resolve
to real papers with matching titles. Journal volume/page details should
be confirmed against the journals of record at proof.}

\textbf{Foundations --- horizon thermodynamics, holography, entropy.}

\begin{itemize}
\tightlist
\item
  Jacobson, T. (1995). Thermodynamics of spacetime: the Einstein
  equation of state. \emph{Phys. Rev.~Lett.} \textbf{75}, 1260.
  arXiv:gr-qc/9504004.
\item
  Cai, R.-G. \& Kim, S. P. (2005). First law of thermodynamics and
  Friedmann equations of Friedmann--Robertson--Walker universe.
  \emph{JHEP} \textbf{02}, 050. arXiv:hep-th/0501055.
\item
  Gibbons, G. W. \& Hawking, S. W. (1977). Cosmological event horizons,
  thermodynamics, and particle creation. \emph{Phys. Rev.~D}
  \textbf{15}, 2738.
\item
  Bousso, R. (2002). The holographic principle. \emph{Rev.~Mod. Phys.}
  \textbf{74}, 825. arXiv:hep-th/0203101.
\item
  Cohen, A. G., Kaplan, D. B. \& Nelson, A. E. (1999). Effective field
  theory, black holes, and the cosmological constant. \emph{Phys.
  Rev.~Lett.} \textbf{82}, 4971. arXiv:hep-th/9803132.
\item
  Page, D. N. (1993). Information in black hole radiation. \emph{Phys.
  Rev.~Lett.} \textbf{71}, 3743. arXiv:hep-th/9306083.
\item
  Padmanabhan, T. (2012). Emergence and expansion of cosmic space as due
  to the quest for holographic equipartition. arXiv:1206.4916.
\item
  Gough, M. P. (2008). Information equation of state. \emph{Entropy}
  \textbf{10}, 150. \emph{(Earliest of the priority lineage --- dark
  energy from the information/entropy of structure; §7.6.)}
\item
  Gough, M. P. (2011). Holographic dark information energy.
  \emph{Entropy} \textbf{13}, 924. arXiv:1105.4461. \emph{(Priority:
  dark energy as the entropy/information energy of star formation ---
  the structure-sourced-dark-energy idea SEDE shares; §7.6.)}
\item
  Gough, M. P. (2025). Evidence for dark energy driven by star
  formation: information dark energy? \emph{Entropy} (2025). \emph{(Fits
  stellar-mass-density history to the DE equation of state;
  \(\rho _{\mathrm{DE}} \propto  SMD^{0.5}\).)}
\item
  Khlebnikov, S. \& Sheoran, A. (2019). The first heat: production of
  entanglement entropy in the early universe. \emph{Phys. Rev.~D}
  \textbf{100}, 103515. arXiv:1907.00487. \emph{(de Sitter → flat:
  entanglement entropy saturates to a volume law.)}
\item
  Belfiglio, A. \& Luongo, O. (2023). Entanglement area-law violation
  from field--curvature coupling. arXiv:2306.08357.
\end{itemize}

\textbf{Non-extensive / deformed entropy and holographic dark energy.}

\begin{itemize}
\tightlist
\item
  Barrow, J. D. (2020). The area of a rough black hole. \emph{Phys.
  Lett. B} \textbf{808}, 135643. arXiv:2004.09444.
\item
  Tsallis, C. \& Cirto, L. J. L. (2013). Black hole thermodynamical
  entropy. \emph{Eur. Phys. J. C} \textbf{73}, 2487. arXiv:1202.2154.
\item
  Saridakis, E. N. (2020). Barrow holographic dark energy. \emph{Phys.
  Rev.~D} \textbf{102}, 123525. arXiv:2005.04115.
\item
  Barrow, J. D., Basilakos, S. \& Saridakis, E. N. (2021). Big Bang
  Nucleosynthesis constraints on Barrow entropy. \emph{Phys. Lett. B}
  \textbf{815}, 136134. arXiv:2010.00986. \emph{(Modified-Friedmann
  route: Δ ≲ 1.4×10⁻⁴.)}
\item
  Asghari, M. \& Sheykhi, A. (2022). Observational constraints of the
  modified cosmology through Barrow entropy. arXiv:2110.00059.
\item
  Leon, G. \emph{et al.} (2021). Barrow entropy cosmology: an
  observational approach. arXiv:2108.10998.
\end{itemize}

\textbf{Emergent gravity and the dark universe.}

\begin{itemize}
\tightlist
\item
  Verlinde, E. P. (2017). Emergent gravity and the dark universe.
  \emph{SciPost Phys.} \textbf{2}, 016. arXiv:1611.02269 (2016).
\item
  Verlinde, E. P. (2011). On the origin of gravity and the laws of
  Newton. \emph{JHEP} \textbf{04}, 029. arXiv:1001.0785.
\item
  Milgrom, M. (1983). A modification of the Newtonian dynamics as a
  possible alternative to the hidden mass hypothesis. \emph{Astrophys.
  J.} \textbf{270}, 365.
\item
  McGaugh, S. S., Lelli, F. \& Schombert, J. M. (2016). Radial
  acceleration relation in rotationally supported galaxies. \emph{Phys.
  Rev.~Lett.} \textbf{117}, 201101. arXiv:1609.05917.
\end{itemize}

\textbf{Structure formation.}

\begin{itemize}
\tightlist
\item
  Press, W. H. \& Schechter, P. (1974). Formation of galaxies and
  clusters of galaxies by self-similar gravitational condensation.
  \emph{Astrophys. J.} \textbf{187}, 425.
\item
  Sheth, R. K. \& Tormen, G. (1999). Large-scale bias and the
  peak-background split. \emph{Mon. Not. R. Astron. Soc.} \textbf{308},
  119. arXiv:astro-ph/9901122.
\end{itemize}

\textbf{Data.}

\begin{itemize}
\tightlist
\item
  DESI Collaboration (Abdul-Karim, M. \emph{et al.}) (2025). DESI DR2
  results II: measurements of baryon acoustic oscillations and
  cosmological constraints. \emph{Phys. Rev.~D} \textbf{112}, 083515.
  arXiv:2503.14738.
\item
  Planck Collaboration (2020). Planck 2018 results. VI. Cosmological
  parameters. \emph{Astron. Astrophys.} \textbf{641}, A6.
  arXiv:1807.06209.
\item
  Brout, D. \emph{et al.} (2022). The Pantheon+ analysis: cosmological
  constraints. \emph{Astrophys. J.} \textbf{938}, 110. arXiv:2202.04077.
\item
  DES Collaboration (2024). The Dark Energy Survey: cosmology results
  with ∼1500 type Ia supernovae (DES-SN5YR). arXiv:2401.02929.
\item
  Rubin, D. \emph{et al.} (2023). Union3 / UNITY1.5: cosmological
  constraints from a unified type Ia supernova compilation.
  arXiv:2311.12098.
\item
  Alam, S. \emph{et al.} (eBOSS Collaboration) (2021). Completed SDSS-IV
  eBOSS: cosmological implications. \emph{Phys. Rev.~D} \textbf{103},
  083533. arXiv:2007.08991.
\item
  Eskilt, J. R. \& Komatsu, E. (2022). Improved constraints on cosmic
  birefringence from the WMAP and Planck data. \emph{Phys. Rev.~D}
  \textbf{106}, 063503. arXiv:2205.13962.
\end{itemize}

\textbf{Post-DESI DR2 dark-energy literature} (arXiv identifiers
verified 2026-06; confirm journal details before submission).

\begin{itemize}
\tightlist
\item
  Turyshev, S. G. (2026). Dark energy after DESI DR2: observational
  status, reconstructions, and physical models. \emph{Phys. Rev.~D}
  (2026). arXiv:2602.05368.
\item
  Wang, D. \& Mota, D. F. (2025). Did DESI DR2 truly reveal dynamical
  dark energy? \emph{Eur. Phys. J. C} (2025). arXiv:2504.15222.
\item
  Chaudhary, H., Capozziello, S. \emph{et al.} (2025). Evidence for
  evolving dark energy from DESI DR2 BAO and Pantheon+, DES-Dovekie, and
  Union3. arXiv:2508.10514. \emph{(LRG1--2-driven preference.)}
\item
  Liu, G., Wang, Y. \emph{et al.} (2024). Impact of LRG1 and LRG2 in
  DESI 2024 BAO data on dark-energy evolution. arXiv:2407.04385.
\item
  You, C., Wang, D. \emph{et al.} (2025). Dynamical dark energy implies
  a coupled dark sector: insights from DESI DR2 via a data-driven
  approach. arXiv:2504.00985.
\item
  Zhang, X., Xu, Y.-H. \emph{et al.} (2025). Reconstruction of dark
  energy using DESI DR2. arXiv:2511.02220. \emph{(GP w(z); crossing
  \(z_{\mathrm{wt}} \approx  0.46\).)}
\item
  Cai, Y., Ren, X., Qiu, T., Li, M. \& Zhang, X. (2025). The quintom
  theory of dark energy after DESI DR2. arXiv:2505.24732.
\item
  Thanankullaphong, P., Sahoo, P. \emph{et al.} (2026). Quintom dark
  energy: future attractor and phantom crossing in light of DESI DR2.
  arXiv:2601.02284.
\end{itemize}

\textbf{Critical analyses and tensions (red-team).}

\begin{itemize}
\tightlist
\item
  Wu, P.-J., Li, T.-N., Du, G.-H. \& Zhang, X. (2025). Observational
  challenges to holographic and Ricci dark energy paradigms: insights
  from ACT DR6 and DESI DR2. arXiv:2509.02945. \emph{(Full-CMB HDE
  tension; addressed in §9 via the ACT-lensing and plik-lite tests.)}
\item
  Huang, L., Cai, R.-G. \& Wang, S.-J. (2025). The DESI DR1/DR2 evidence
  for dynamical dark energy is biased by low-redshift supernovae.
  arXiv:2502.04212.
\item
  Luciano, G. G., Paliathanasis, A. \& Saridakis, E. N. (2025).
  Constraints on Barrow and Tsallis holographic dark energy from DESI
  DR2 BAO data. arXiv:2506.03019. \emph{(Standard Barrow-HDE: Δ
  \textless{} 0.47--0.54.)}
\item
  Diego-Palazuelos, P. \& Komatsu, E. (2026). Cosmic birefringence from
  the Atacama Cosmology Telescope Data Release 6. arXiv:2509.13654.
  \emph{(β = 0.215°±0.074°, 2.9σ.)}
\item
  Barman, B. \& Girmohanta, S. (2026). Implications of DESI for dark
  matter \& cosmic birefringence. arXiv:2506.12589. \emph{(One axion
  sector: phantom crossing + β + DM --- the unified rival to SEDE.)}
\item
  Lin, W., Visinelli, L. \emph{et al.} (2025). Testing quintessence
  axion dark energy with recent cosmological results. arXiv:2504.17638.
  \emph{(Axion: β + acceleration at 1σ.)}
\item
  Nakagawa, S., Nakai, Y. \emph{et al.} (2025). Interpreting cosmic
  birefringence and DESI data with an evolving axion in ΛCDM.
  arXiv:2503.18924.
\item
  Nakagawa, S., Naokawa, F., Namikawa, T. \& Komatsu, E. (2023).
  Redshift evolution of cosmic birefringence in CMB anisotropies.
  \emph{Phys. Rev.~D} \textbf{107}, 083529. arXiv:2301.07971.
  \emph{(Universal late-time β(z) profile.)}
\item
  Arcari, S., Bartolo, N. \emph{et al.} (2025). Stairway to axions:
  cross-correlation of birefringence and galaxies from NPIPE and Quaia.
  arXiv:2509.22273. \emph{(Birefringence tomography.)}
\item
  BICEP/Keck Collaboration (2026). BICEP/Keck XXI: constraints on
  early-universe parity violation from multipole-dependent
  birefringence. arXiv:2603.06812.
\item
  Yin, L., Xiong, S., Kochappan, J., Lee, B.-H. \& Ghosh, T. (2025).
  Constraints on cosmic birefringence from SPIDER, Planck, and ACT
  observations. arXiv:2510.25489. \emph{(α+β degeneracy; Planck--ACT
  differ \textgreater3σ; SPIDER incompatible with the coupling.)}
\end{itemize}

\textbf{Growth--geometry and response probes.}

\begin{itemize}
\tightlist
\item
  Rauhut, S. J., Blake, C. \emph{et al.} (DESI Collaboration) (2025).
  Testing gravitational physics by combining DESI DR1 and weak lensing
  datasets using the \(E_{\mathrm{G}}\) estimator. arXiv:2507.16098.
\item
  Barreira, A. \& Schmidt, F. (2017). Complete super-sample lensing
  covariance in the response approach. \emph{JCAP} \textbf{06}, 053.
  arXiv:1711.07467.
\end{itemize}

\textbf{Codes and inference.}

\begin{itemize}
\tightlist
\item
  Lewis, A., Challinor, A. \& Lasenby, A. (2000). Efficient computation
  of CMB anisotropies in closed FRW models (CAMB). \emph{Astrophys. J.}
  \textbf{538}, 473. arXiv:astro-ph/9911177.
\item
  Blas, D., Lesgourgues, J. \& Tram, T. (2011). The Cosmic Linear
  Anisotropy Solving System (CLASS) II. \emph{JCAP} \textbf{07}, 034.
  arXiv:1104.2933.
\item
  Foreman-Mackey, D., Hogg, D. W., Lang, D. \& Goodman, J. (2013).
  emcee: the MCMC hammer. \emph{Publ. Astron. Soc. Pac.} \textbf{125},
  306. arXiv:1202.3665.
\item
  Torrado, J. \& Lewis, A. (2021). Cobaya: code for Bayesian analysis of
  hierarchical physical models. \emph{JCAP} \textbf{05}, 057.
  arXiv:2005.05290.
\end{itemize}

\section{Appendix A --- Key results, prescriptions, and
assumptions}\label{appendix-a-key-results-prescriptions-and-assumptions}

This appendix documents the six load-bearing items, labelled by their
actual epistemic status --- one defining ansatz (A.1), supporting
derivations (A.2, A.4, A.5), a fixed prescription (A.3), and a motivated
postulate (A.6). The headings give that status; the parenthetical
``Result/Prescription/ Motivation N'' numbering corresponds to the
repository code labels \texttt{THEOREM\_N} in \texttt{sede/theory.py}
(retained there for code stability, \emph{not} a claim that each is a
mathematical theorem). We work in units
\(c = \hbar  = k_{\mathrm{B}} = 1\) and write
\(8\pi G = 1/M_{\mathrm{P}}^{2}\). Source statements
\texttt{THEOREM\_1}, \texttt{THEOREM\_3}, \texttt{THEOREM\_4C},
\texttt{THEOREM\_5}, \texttt{THEOREM\_8}, \texttt{THEOREM\_9}; the
remaining items (2, 4, 4B, 6, 6B, 10--13, 5D) are catalogued at the end
with one-line statements.

\subsection{\texorpdfstring{A.1 The conjugate horizon-fluid ansatz:
\(\rho _{\mathrm{DE}} = T_{\mathrm{AH}}\cdot s_{\mathrm{grav}}\)}{A.1 The conjugate horizon-fluid ansatz: \textbackslash rho \_\{\textbackslash mathrm\{DE\}\} = T\_\{\textbackslash mathrm\{AH\}\}\textbackslash cdot s\_\{\textbackslash mathrm\{grav\}\}}}\label{a.1-the-conjugate-horizon-fluid-ansatz-rho-_mathrmde-t_mathrmahcdot-s_mathrmgrav}

\textbf{Claim.} SEDE \emph{identifies} the dark-energy density with the
conjugate energy of the horizon entropy,
\(\rho _{\mathrm{DE}} = T_{\mathrm{AH}}\cdot s_{\mathrm{grav}}\cdot f_{\mathrm{sat}}\),
with \(s_{\mathrm{grav}}\) the (volume-law) horizon entropy density and
\(f_{\mathrm{sat}} \in  [0\),1{]}. The Clausius relation supplies the
\emph{dynamics}; the absolute density ρ = T·s is the model's defining
\textbf{ansatz}, motivated by horizon thermodynamics --- not a theorem.

\textbf{Setup.} At the flat-FRW apparent horizon
\(R_{\mathrm{AH}} = 1/H\), the dynamical Cai--Kim temperature is
\(T_{\mathrm{AH}} = (H/2\pi )(1 - \epsilon /2)\), ε ≡ −Ḣ/H². We adopt
the Gibbons--Hawking limit \(T_{\mathrm{AH}} = H/2\pi\) (ε → 0), the
leading temperature appropriate for the de Sitter-attractor dark-energy
sector; the dynamical (1 − ε/2) correction is sub-leading where dark
energy matters and is disfavoured in full (§2, §6), so it does not enter
the calibrated E(z). The gravitational entropy is the coarse-grained
\textbf{entanglement entropy in the horizon volume},
\(S_{\mathrm{grav}} \propto  V_{\mathrm{AH}} \propto  R_{\mathrm{AH}}^{3}\)
(§3.3, §8.1), i.e.~a \emph{constant} entropy density
\(s_{\mathrm{grav}} = s_{0}\) (in contrast to the Bekenstein--Hawking
area-law density
\(s_{\mathrm{AH}}\)=¾\(M_{\mathrm{P}}^{2}H \propto  H\), which would
instead give the area-law model \(\rho _{\mathrm{DE}} \propto  H^{2}\)
--- the branch disfavoured in the conditional diagnostics of §5.6).

\textbf{Construction.} \emph{Step 1 (Clausius --- the rigorous,
dynamical content).} Hayward's first law for a dynamical horizon,
\(\delta E = T_{\mathrm{AH}}\) \(dS_{\mathrm{AH}} + W\) δV with work
density W = (ρ − p)/2, applied to a shell of matter crossing the horizon
in time dt with flux −dE = (ρ + p)H V dt, gives the Clausius relation (ρ
+ p)H \(V = T_{\mathrm{AH}}\) Ṡ\_AH. This relates \emph{changes} and is
standard horizon thermodynamics.

\emph{Step 2 (Bousso bound as motivation for a bounded
\(f_{\mathrm{sat}}\)).} The covariant entropy bound gives
\(S_{\mathrm{matter}} \leq  A/(4l_{\mathrm{P}}^{2}) = S_{\mathrm{AH}}\)
on any light-sheet bounded by the horizon. This \emph{motivates}
interpreting \(f_{\mathrm{sat}}\) as a bounded effective saturation
fraction; in SEDE we impose \(0 \leq  f_{\mathrm{sat}} \leq  1\) and
\(f_{\mathrm{sat}}(0) = 1\) with the redshift dependence supplied by the
halo prescription (A.2/A.3), and we do \textbf{not} identify
\(f_{\mathrm{sat}}\) literally with
\(S_{\mathrm{matter}}/S_{\mathrm{AH}}\) (consistent with the small
binding-entropy budget, App. F).

\emph{Step 3 (the SEDE identification --- the ansatz).} SEDE posits the
dark-energy density to be the conjugate energy of the volume-law horizon
entropy,
\[\rho_{\rm DE} \equiv T_{AH}\,s_{\rm grav}\,f_{\rm sat} = \frac{H}{2\pi}\,s_0\,f_{\rm sat} \;\propto\; H\,f_{\rm sat}.\]
Fixing the one constant \(s_{0}\) by spatial flatness
\([\rho _{\mathrm{DE}}(0) = \Omega _{\mathrm{DE0}}\)
\(\rho _{\mathrm{crit}}\),0, \(f_{\mathrm{sat}}(0) = 1\){]} gives
\[\rho_{\rm DE}(z) = \Omega_{\rm DE0}\,\rho_{\rm crit,0}\,\frac{H}{H_0}\,f_{\rm sat}(z)
= \Omega_{\rm DE0}\,\rho_{\rm crit,0}\,E(z)\,f_{\rm sat}(z),\] i.e.~the
fixed-point term \(\Omega _{\mathrm{DE0}}\cdot E\cdot f_{\mathrm{sat}}\)
of eq. (1) --- the \textbf{volume-law (λ = 1/2)} scaling, \emph{not} the
area-law H². This is an identification motivated by horizon
thermodynamics, not a derivation: the Clausius relation (Step 1) fixes
the dynamics, while the absolute density is the postulate (§8.4).

\emph{Step 4 (scale / coincidence).} \(\rho _{\mathrm{DE}}\)
\textasciitilde{} \(M_{\mathrm{P}}^{2}H_{0}^{2}\) \textasciitilde{}
(10⁻³ eV)⁴ is small because \(H_{0}/M_{\mathrm{P}}\) \textasciitilde{}
10⁻⁶¹. The CKN bound (§8.2) \emph{explains} why this horizon scale is
\textasciitilde{}\(\rho _{\mathrm{crit}}\) --- as an upper bound, taken
here to be saturated to O(1) by the flatness normalisation above ---
replacing the QFT-vacuum conflation behind the 10¹²⁰ problem rather than
deriving the magnitude from scratch. ∎

\subsection{\texorpdfstring{A.2 Result 3 --- the structure gate
\(f_{\mathrm{sat}}(z)\)}{A.2 Result 3 --- the structure gate f\_\{\textbackslash mathrm\{sat\}\}(z)}}\label{a.2-result-3-the-structure-gate-f_mathrmsatz}

\textbf{Claim.}
\(f_{\mathrm{sat}}(z) = [1 - e^{-\gamma  x(z)}]/[1 - e^{-\gamma }]\),
where x ≡ σ8²(z)/σ8²(0) = D²(z) (D the linear growth factor, D(0) = 1).

\textbf{Setup.} We model the cumulative activation gate directly
(dimensionless, avoiding any mix of area- and volume-law densities)
through a normalised growth-weighted deposition kernel K(t) ∝ (dσ8²/dt)
\(e^{-\gamma \sigma 8^{2}/\sigma 8^{2}(0)}\) --- the structure-growth
rate dσ8²/dt \textgreater{} 0 times a Boltzmann factor \(e^{-\gamma x}\)
for the exponential rarity of massive halos (γ from A.3) --- taking
\(f_{\mathrm{sat}}\) to be its normalised cumulative from early times to
epoch z (cosmic time t),
\[f_{\rm sat}(z) = \frac{\int_0^{t(z)} K\,dt}{\int_0^{t_0} K\,dt},\] so
that \(f_{\mathrm{sat}}(\infty ) = 0\) (early, t→0) and
\(f_{\mathrm{sat}}(0) = 1\) (today, \(t = t_{0}\)) by construction (the
closed-form gate below is this cumulative-kernel form, not a relaxation
ODE).

\textbf{Proof.} The kernel integrated over cosmic time is K dt =
(dσ8²/dt) \(e^{-\gamma \sigma 8^{2}/\sigma 8^{2}(0)}\) dt =
\(e^{-\gamma u}\) dσ8² with u ≡ σ8²/σ8²(0); the Hubble factors cancel
and the integral reduces to one over u. With u running from 0 (z′→∞, no
early structure) to x(z) ≡ σ8²(z)/σ8²(0) = D²(z),
\[f_{\rm sat}(z) = \frac{\int_0^{x(z)} e^{-\gamma u}\,du}{\int_0^{1} e^{-\gamma u}\,du}
= \frac{1 - e^{-\gamma x(z)}}{1 - e^{-\gamma}}, \qquad x(z) = D^2(z),\]
where the denominator (x = 1, i.e.~z = 0) enforces
\(f_{\mathrm{sat}}(0) = 1\). Checks: \(f_{\mathrm{sat}}(0) = 1\);
\(f_{\mathrm{sat}}(z \gg  1)\) → 0 as x → 0;
\(f_{\mathrm{sat}} \in  [0\),1{]} for x ∈ {[}0,1{]}. The form is the CDF
of a truncated exponential of rate γ on x ∈ {[}0,1{]} --- the normalised
activation fraction reached by epoch z. ∎

\subsection{A.3 Prescription 4C --- the fixed structure coupling γ ≈
1.496}\label{a.3-prescription-4c-the-fixed-structure-coupling-ux3b3-1.496}

\textbf{Claim.} The entropy weight exponent is p = 5/3, giving γ = d ln
\(\Sigma _{\mathrm{S}}^{(5/3)}/d\) ln σ8² = 1.496 (the derivative is
w.r.t. the variance σ8², the gate argument x = D² = σ8²/σ8²(0) of A.2;
equivalently ½ d ln \(\Sigma _{\mathrm{S}}/d\) ln σ8 = 1.496, i.e.~d ln
\(\Sigma _{\mathrm{S}}/d\) ln σ8 = 2.993). A self-similar reduction
gives the closed form γ = (p−1)⟨1/α⟩\_Σ with α ≡ −d ln σ/d ln M the
effective spectral slope over the entropy-weighted mass range; it
reproduces the exact mass-function integral to 0.1\%
(\texttt{run\_gamma\_systematics.py},
\texttt{BINDING\_ENERGY\_MECHANISM.md} §7). Entropy is deposited by the
\(\nu _{\mathrm{S}} \approx  2.4\) (\textasciitilde2σ) clusters at
\(M_{\mathrm{S}}\) \textasciitilde{} 10¹⁴·⁵ M⊙/h, giving an independent
cluster-abundance handle; p = 1 gives ≈ 0.27 (∼5× below p = 5/3), a weak
gate versus the real one.

\textbf{Proof.} In the SEDE horizon-fluid ansatz (A.1),
\(f_{\mathrm{sat}}\) is the effective fraction of the
\textbf{cosmic-horizon} entropy capacity that is gravitationally
activated, and
\(\rho _{\mathrm{DE}} = T_{\mathrm{AH}}\cdot s_{\mathrm{grav}}\cdot f_{\mathrm{sat}}\)
--- so the entropy that defines the \(f_{\mathrm{sat}}\) weighting is
accounted in the \emph{horizon} reservoir at the \emph{horizon}
temperature \(T_{\mathrm{AH}} = H/2\pi\). By Clausius, heat Q deposited
into a reservoir at temperature T adds entropy Q/T. When a halo of mass
M virialises it releases gravitational binding energy
\(E_{\mathrm{bind}} \propto  M^{5/3}\) (NFW;
\(R_{\mathrm{vir}} \propto  M^{1/3}\); verified slope 1.64). Under the
SEDE prescription this released binding energy is accounted in the
horizon entropy ledger at the horizon temperature \(T_{\mathrm{AH}}\)
(we do not claim a demonstrated transport mechanism), so the entropy
assigned in this ledger is
\[\Delta S_{AH} = \frac{E_{\rm bind}}{T_{AH}} \propto \frac{M^{5/3}}{H}.\]
Because \(T_{\mathrm{AH}} \propto  H\) is identical for every halo at a
given epoch (mass-independent), the per-halo mass weight is \(M^{5/3}\):
\textbf{p = 5/3}. With \(\Sigma _{\mathrm{S}}\)=∫ \(M^{5/3}(dn/d\) ln
M)d ln M and a Sheth--Tormen mass function (COLOSSUS, EH98 transfer, M ∈
{[}10¹⁰, 10¹⁶{]} \(M_\odot\)),
\[\gamma = \frac{d\ln\Sigma_S^{(5/3)}}{d\ln\sigma_8} = 1.4964 \approx 1.50.\]
The earlier p = 1 (Result 4B) divided \(E_{\mathrm{bind}}\) by the
\emph{halo's} virial temperature \(T_{\mathrm{vir}} \propto  M^{2/3}\),
which is the entropy of the halo's own internal state --- correct for
the halo, wrong for the entropy added to the cosmic horizon (where
\(f_{\mathrm{sat}}\) lives). The heat is the same; the reservoir, and
hence the temperature, is the horizon's. ∎

\subsection{A.4 Result 5 --- the structural-EOS diagnostic (the w₀
side-note, not canonical
w(z))}\label{a.4-result-5-the-structural-eos-diagnostic-the-wux2080-side-note-not-canonical-wz}

\textbf{Claim (structural EOS).}
\(w_{\mathrm{struct}}(0) = -1 + \gamma /[3(e^{\gamma } - 1)]\),
well-approximated by the closed-form algebraic expression
\(w_{\mathrm{alg}} = (4\Omega _{\mathrm{m}}/3 - 1)/(1 - \Omega _{\mathrm{m}})\).
(This is the side-note structural estimate of §3.2, not the model's
prediction, which is the canonical crossing w(z).)

\textbf{Proof.} From
\(\rho _{\mathrm{DE}}(z) = \rho _{\mathrm{crit}}\),0
\(\Omega _{\mathrm{DE0}}\) \(f_{\mathrm{sat}}(z)\) (A.1) with
\(f_{\mathrm{sat}}(x)\), x = D², the structural EOS in the variance
variable is \(w_{\mathrm{struct}}(0) = -1 + (1/3)\)
\(df_{\mathrm{sat}}/dx\)\textbar\_\{x=1\}. Differentiating the Theorem-3
form,
\[\left.\frac{df_{\rm sat}}{dx}\right|_{x=1} = \frac{\gamma e^{-\gamma}}{1 - e^{-\gamma}} = \frac{\gamma}{e^\gamma - 1},\]
so \(w_{\mathrm{struct}}(0) = -1 + \gamma /[3(e^{\gamma } - 1)]\). At γ
= 1.4964 this is 0.432, giving
\textbf{\(w_{\mathrm{struct}}(0) = -0.856\)} (0.33σ from DESI DR2).
Numerically
\(\gamma /(e^{\gamma } - 1) = 0.432 \approx  \Omega _{\mathrm{m}}/(1 - \Omega _{\mathrm{m}}) = 0.451\)
at the observed parameters, so the structural value is approximated to
\textasciitilde4\% by the closed-form expression
\(w_{\mathrm{alg}} = (4\Omega _{\mathrm{m}}/3 - 1)/(1 - \Omega _{\mathrm{m}}) = -0.849\)
(0.21σ from DESI). This is a numerical approximation in the x = D²
convention, \textbf{not} an exact identity. (The interacting-fluid
effective EOS, from Raychaudhuri + Friedmann at z = 0, is phantom,
\(w_{\mathrm{fluid}}(0) \approx  -1.15\), because \(f_{\mathrm{sat}}\)
is still rising today; the canonical λ = 1/2 background realises this as
a crossing of −1, §5.1.) ∎

\subsection{A.5 Result 8 --- the H-coupling λ = 1 −
Δ/2}\label{a.5-result-8-the-h-coupling-ux3bb-1-ux3b42}

\textbf{Claim.} With the Barrow entropy \(S = (A/A_{0})^{1+\Delta /2}\),
the conjugate horizon-fluid ansatz gives
\(\rho _{\mathrm{DE}} \propto  H^{2\lambda }\) \(f_{\mathrm{sat}}\) with
λ = 1 − Δ/2.

\textbf{Proof.} A quantum-gravitationally rough horizon carries the
Barrow entropy \(S = (A/A_{0})^{1+\Delta /2}\), Δ ∈ {[}0,1{]}, in place
of the area law. With A = 4π/H² and \(V_{\mathrm{AH}} \propto  H^{-3}\),
the \textbf{Barrow-deformed} entropy density (written \(s_\Delta\) to
distinguish it from the Bekenstein--Hawking area-law density
\(s_{\mathrm{AH}} \propto  H\)) is
\[s_{\Delta} = \frac{S}{V_{AH}} \propto \frac{(H^{-2})^{1+\Delta/2}}{H^{-3}} = H^{1-\Delta},\]
\[\rho_{\rm DE} = T_{AH}\,s_{\Delta}\,f_{\rm sat} \propto H\cdot H^{1-\Delta}\cdot f_{\rm sat}
= H^{2-\Delta}f_{\rm sat} \equiv H^{2\lambda}f_{\rm sat}, \qquad \boxed{\lambda = 1 - \tfrac{\Delta}{2}.}\]
This is Barrow holographic dark energy with the Hubble IR cutoff
\((\rho _{\mathrm{DE}} \propto  L^{\Delta -2}\), L = 1/H). The endpoints
are Δ = 0 ⟹ λ = 1 (the area-law branch, disfavoured) and Δ = 1 ⟹ λ = 1/2
(the volume-law postulate, A.6). λ is therefore not a free parameter; it
is set by the Barrow deformation. \textbf{Scope:} the Barrow exponent
multiplies only the \(f_{\mathrm{sat}}\)-gated \(\rho _{\mathrm{DE}}\),
\emph{not} the gravitational sector, so H(z) at BBN equals the standard
matter+radiation value to \textasciitilde8 figures
\((\Omega _{\mathrm{DE}}/\rho _{\mathrm{tot}}\) \textasciitilde{} 10⁻¹⁶
at z \textasciitilde{} 4×10⁹); the modified-gravity BBN bound Δ ≲ 10⁻⁴
does not apply, and constant Δ = 1 is BBN-safe. ∎

\subsection{A.6 Motivation 9 --- Δ = 1 from a geometric ceiling and the
GSL}\label{a.6-motivation-9-ux3b4-1-from-a-geometric-ceiling-and-the-gsl}

\textbf{Claim.} Δ = 1 (the volume-law endpoint) is \emph{strongly
motivated} --- not fitted to cosmological data --- by two independent
principles; it is the central postulate (§8.4), not a microscopic
derivation.

\textbf{Pillar 1 --- geometric ceiling.} The Barrow deformation is the
fractal (Hausdorff) dimension of the horizon surface,
\(d_{\mathrm{H}} = 2 + \Delta\). A 2-surface embedded in 3-space cannot
exceed the space-filling dimension \(d_{\mathrm{H}} = 3\), so \textbf{Δ
≤ 1}, with Δ = 1 the extremum where the area-measure becomes a
volume-measure \((S \propto  A^{3/2} \propto  R^{3})\). This is Barrow's
defining range.

\textbf{Pillar 2 --- thermodynamic attractor.} Three premises: (P1) the
DE sector is purely entropic --- SEDE posits
\(\rho _{\mathrm{DE}} = T_{\mathrm{AH}}\) \(s_{\mathrm{grav}}\), so the
free-energy density \(F = \rho _{\mathrm{DE}} - T_{\mathrm{AH}}\)
\(s_{\mathrm{grav}} = 0\) identically; the governing potential is
therefore the \emph{entropy}, not F = E − TS (this is the load-bearing
step: with a nonzero Δ-dependent free energy, equilibrium would
extremise F and Δ = 1 would not follow). (P2) the generalised second
law: \(S_{\mathrm{total}} = S_\Delta  + S_{\mathrm{matter}}\) is
non-decreasing, \(S_{\mathrm{matter}}\) independent of Δ. (P3) Δ is a
thermalising horizon degree of freedom (the standard
Jacobson--Padmanabhan assumption). By (P1)--(P2), maximising
\(S_{\mathrm{total}}\) over Δ reduces to maximising
\(S_\Delta  = (A/A_{0})^{1+\Delta /2}\), whose derivative is
\[\frac{dS_\Delta}{d\Delta} = \frac{S_\Delta}{2}\ln\!\frac{A}{A_0} + [\text{back-reaction through }A(\Delta)].\]
The explicit term carries ln(A/A₀) ≈ ln(10¹²²) ≈ 280 \textgreater{} 0 (a
vastly super-Planckian horizon). The back-reaction (Δ shifts
\(\rho _{\mathrm{DE}} \propto  H^{2-\Delta }\), hence H, hence A) cannot
flip the sign: it \emph{vanishes} at z = 0 (H = H₀ ⟹ ln H = 0 ⟹ ∂\_Δ
\(H^{2-\Delta } = 0\)) and is \textasciitilde1/280-suppressed elsewhere.
Hence \(dS_\Delta /d\Delta  > 0\) on all of {[}0,1{]}; a strictly
increasing function on a closed interval attains its maximum at the
right endpoint. With the ceiling Δ ≤ 1 (Pillar 1) the constrained
maximum is uniquely \textbf{Δ = 1}, and by (P3) a thermalising horizon
relaxes to it.

\textbf{Conclusion.} Under the additional assumptions that Δ is a
thermalising horizon degree of freedom (P3) and that entropy
maximisation over the Barrow family is the correct variational principle
(P1), the GSL argument is a \textbf{consistency check} --- the endpoint
Δ = 1 is a GSL \emph{fixed point} --- rather than a selection principle:
the geometric ceiling \(d_{\mathrm{H}} = 3\) supplies the bound Δ ≤ 1,
and saturation of that cap \emph{is} the volume-law postulate in
effective form (via λ = 1 − Δ/2 this gives λ = 1/2). These assumptions
\textbf{motivate but do not microscopically derive} the endpoint (§8.4).
Because \(S_\Delta\) is monotone in Δ, the same argument favours Δ = 1
at all epochs (the adopted constant Δ). The current conditional
diagnostics prefer values near Δ(0) ≈ 1.0--1.1 (§5.6), consistent with
the endpoint. This is a strong motivation, not a derivation: it reduces
Δ = 1 from a fitted coupling to a GSL-equilibrium extremum under the
stated premises, leaving the microscopic origin open. ∎

\subsection{A.7 The remaining repository items
(catalogue)}\label{a.7-the-remaining-repository-items-catalogue}

\textbf{Result 2} (reheating master equation): the same logistic
\(df_{\mathrm{sat}}/dt = H\)
\(f_{\mathrm{sat}}(1 - f_{\mathrm{sat}}) - \Gamma\) \(f_{\mathrm{sat}}\)
governs the \emph{early} \(f_{\mathrm{sat}}\) (inflaton→radiation),
giving the U-shape of Result 12. \textbf{Result 4 / 4B} {[}superseded by
4C{]}: the \(M^{5/3}\) energy weight, and the p = 1 virial-temperature
mis-step. \textbf{Result 5D}: the (1 − ε/2) ``λ = 1 cousin'' background
(w₀ = −0.85, no crossing) vs the canonical λ = 1/2 crossing background.
\textbf{Result 6 / 6B}: the smooth-DE growth validation and the
\(c_{\mathrm{s}}^{2} = 1\) closure \((f_{\mathrm{sat}}\) a background
functional). \textbf{Result 10}: the QG chain and the volume-law
postulate. \textbf{Result 11}: the seam = BBN bound =
modgrav/holographic fork. \textbf{Result 12}: the \(f_{\mathrm{sat}}\)
U-shape (inflation and dark energy as the two de Sitter brackets).
\textbf{Result 13}: 1 + w = (1/3){[}2λε − d ln \(f_{\mathrm{sat}}/d\) ln
a{]}, the EOS as horizon-entropy production. Full statements in
\texttt{sede/theory.py}; each is exercised by a numbered check in App.
B.

\section{Appendix B --- The verification
suite}\label{appendix-b-the-verification-suite}

The repository ships a \texttt{sede/verification.py} suite of 63
numerical checks (V1--V63) covering: the fixed-point solution
(closed-loop H(z) to 10⁻¹³ over z ∈ {[}0, 10⁶{]}), BBN-safety
\((\Delta N_{\mathrm{eff}} = \Delta Y_{\mathrm{p}} = 0)\), the
CAMB≡CLASS \(r_{\mathrm{drag}}\) agreement (5×10⁻⁵), the CLASS
perturbation validation (0.1\%), the GSL along cosmic history, the fixed
γ prescription and the smooth-DE \(c_{\mathrm{s}}^{2} = 1\) closure, the
horizon-dimension scaling
\(\lambda (d_{\mathrm{H}}) = 2 - d_{\mathrm{H}}/2\) (V59), the §8.5
state/counting diagnostics (SYK scrambling V60; causal-set dof counting
in Minkowski + de Sitter V61; the entropy-bound audit V62; the
tensor-network realisability test V63), and the cross-horizon and
quantum-sector consistency checks. \texttt{reproduce\_all.py}
regenerates every number and figure in this paper from the staged
drivers.

\section{Appendix C --- Data vectors, covariances, and the fσ8
audit}\label{appendix-c-data-vectors-covariances-and-the-fux3c38-audit}

The full data vectors and covariance matrices, the positive-definiteness
checks, the benign Pantheon+ duplicate-CID treatment, and the fσ8 data
audit (the small set of mis-entered growth points corrected to their
published Gold-2018 values, and the SH0ES value correction) --- all
reproduced from \texttt{data/} and documented in
\texttt{DATA\_AUDIT.md}.

\section{Appendix D --- False-preference calibration
methodology}\label{appendix-d-false-preference-calibration-methodology}

The parametric-bootstrap procedure of §5.2: generation of ΛCDM-truth
mocks with \textbf{all} probes (including the compressed CMB
\(R/l_{\mathrm{A}}\) and SH0ES) re-drawn self-consistently; the refit of
both models to each mock; the null distribution of Δχ²(SEDE − ΛCDM)
(mean +0.37, 2000 mocks); the p ≈ 0.004 (95\% UL 0.006); the documented
bug (CMB held fixed → biased null −2.87, p ≈ 0.5) and its fix; and the
sibling-pipeline cross-check (\textasciitilde2σ,
CMB-distance-engine-dependent). The pre-registration sha256 and contents
are listed. The standalone \texttt{PREREGISTRATION.md} additionally
records the triggered CMB follow-ups and their status (§E: full non-lite
plik \textbf{run}, Δχ²=−0.33; ACT DR6 primary formally deferred) and the
locked E1 structure-sourcing decision rule (§F: data vector +
CONFIRM/REFUTE thresholds). The no-SH0ES robustness (§5.4 ii′;
Δχ²=−1.75, ln B=+0.95, p=0.030) is reproduced by
\texttt{run\_no\_shoes\_robustness.py}.

\section{Appendix E --- The QG derivation
ledger}\label{appendix-e-the-qg-derivation-ledger}

The full chain from §8 with each link labelled by status: -
\textbf{Jacobson--Cai--Kim (MOTIVATED / ANSATZ):} Clausius horizon
thermodynamics motivates the horizon-fluid construction and supplies the
conservation/dynamical logic (δQ = T dS ⟹ Friedmann); the
\emph{absolute} identification \(\rho _{\mathrm{DE}} = T\cdot s\) is the
SEDE ansatz (A.1), not a theorem. - \textbf{Non-extensivity ⟹
Tsallis--Cirto/Barrow power law (MOTIVATED):} the weakest link. -
\textbf{Second-law / GSL ⟹ volume-law endpoint Δ = 1 (MOTIVATED, given
the volume-law thermalisation premise):} the geometric ceiling caps Δ ≤
1 and the GSL is \emph{consistent with} (a fixed point at) the saturated
endpoint --- a consistency check, not a selection principle; a
microscopic derivation remains \textbf{OPEN}. - \textbf{CKN scale
(RIGOROUS BOUND, normalised by flatness):} the holographic energy bound
gives the horizon-scale upper bound; the amplitude is normalised by
flatness (§8.2). - \textbf{The one irreducible postulate (OPEN):} a
microscopic state count delivering the volume law.

The perturbative-QG (log-correction) vs thermalised (volume-law) regimes
are reconciled as different \emph{states}, not a contradiction. This
ledger matches §8.4 and Appendix F.

\section{Appendix F --- The energy and entropy
budget}\label{appendix-f-the-energy-and-entropy-budget}

This appendix records the back-of-envelope that the energy-conservation
framing of §1--§2 turns on, and the resulting ledger of what is derived,
normalised, and open.

\textbf{Energy budget.} Dark energy today is \(\rho _{\mathrm{DE}}\) c²
≈ 0.7 \(\rho _{\mathrm{crit}}\) c² ≈ 6×10⁻¹⁰ J/m³. The gravitational
binding energy released by structure is bounded by
\textbar{}\(E_{\mathrm{bind}}\)\textbar/Mc² \textasciitilde{} σ²/c² per
virialised system, with σ \textasciitilde{} 10³ km/s for clusters (σ²/c²
\textasciitilde{} 10⁻⁵) and \textasciitilde200 km/s for galaxies
(\textasciitilde10⁻⁷); even taking all of \(\Omega _{\mathrm{m}}\)
collapsed, \(\rho _{\mathrm{bind}}\) c² ≲ 10⁻⁵ · 0.3 · 8.5×10⁻¹⁰
\textasciitilde{} 10⁻¹⁵ J/m³. Hence
\(\rho _{\mathrm{bind}}/\rho _{\mathrm{DE}}\) \textasciitilde{} 10⁻⁶:
\textbf{binding energy is six orders of magnitude too small to \emph{be}
the dark energy.} A literal matter→dark-energy transfer is excluded
independently --- moving \(O(\rho _{\mathrm{crit}})\) out of matter over
a Hubble time would spoil \(\rho _{\mathrm{m}} \propto  a^{-3}\) and
with it BBN and the CMB.

\textbf{Entropy budget.} Deposited at the cold horizon temperature, the
binding entropy is \(\Delta S_{\mathrm{AH}}\)=
\(E_{\mathrm{bind}}/T_{\mathrm{AH}}\); with
\(T_{\mathrm{AH}} = \hbar H/2\pi\) this comes to
\(\Delta S_{\mathrm{AH}}\) \textasciitilde{} 10⁻⁷ \(S_{\mathrm{AH}}\),
i.e.~\textasciitilde10⁻⁷ of the horizon entropy. So binding-entropy
deposition does not, by itself, fill \(f_{\mathrm{sat}}\) to O(1); it
sets the \emph{relative} z-dependence (the weight p = 5/3 → γ), while
the normalisation \(f_{\mathrm{sat}}(0) = 1\) is set by the horizon
capacity. This is the same family as the
\((M_{\mathrm{P}}/H)^{\Delta }\) \textasciitilde{} 10⁶¹ ``seam'' (Result
11), resolved by the holographic-DE scope (Barrow acts only on the
\(f_{\mathrm{sat}}\)-gated dark-energy fluid).

\textbf{The ledger.}

\begin{longtable}[]{@{}
  >{\raggedright\arraybackslash}p{(\linewidth - 4\tabcolsep) * \real{0.3333}}
  >{\raggedright\arraybackslash}p{(\linewidth - 4\tabcolsep) * \real{0.3333}}
  >{\raggedright\arraybackslash}p{(\linewidth - 4\tabcolsep) * \real{0.3333}}@{}}
\toprule\noalign{}
\begin{minipage}[b]{\linewidth}\raggedright
quantity
\end{minipage} & \begin{minipage}[b]{\linewidth}\raggedright
status
\end{minipage} & \begin{minipage}[b]{\linewidth}\raggedright
set by
\end{minipage} \\
\midrule\noalign{}
\endhead
\bottomrule\noalign{}
\endlastfoot
magnitude \(\rho _{\mathrm{DE}}\) \textasciitilde{}
\(\rho _{\mathrm{crit}}\) & \textbf{bounded + normalised} & CKN
holographic upper bound (§8.2), amplitude fixed by flatness (O(1)
saturation of the bound assumed) \\
H-scaling λ = 1/2 & \textbf{postulate} & volume-law horizon entropy
\(S_{\mathrm{grav}} \propto  V_{\mathrm{AH}}\) (§8.1) \\
shape γ ≈ 1.5 & \textbf{fixed prescription} & halo-binding-energy
weighting \(E_{\mathrm{bind}} \propto  M^{5/3}\) → p = 5/3 (§3.1) \\
sound speed \(c_{\mathrm{s}}^{2} = 1\) & \textbf{closure} & minimal
smooth-DE treatment, \(f_{\mathrm{sat}}\) a background functional
(§3.4) \\
normalisation \(f_{\mathrm{sat}}(0) = 1\) & \textbf{imposed} &
present-day horizon-entropy normalisation \\
volume-law horizon (Δ = 1 microscopically) & \textbf{open} & the one
irreducible postulate (§8.4) \\
\end{longtable}

The honest reading: none of these is a parameter \emph{fitted to
cosmological data}, but neither are they all ``derived from first
principles'' --- the magnitude is a holographic bound normalised by
flatness (as ΛCDM normalises \(\Omega _\Lambda 0\)), the H-scaling is
the central volume-law postulate, γ is a fixed halo prescription, and
\(c_{\mathrm{s}}^{2} = 1\) is a modelling closure. Binding energy enters
SEDE entropically, fixing the \emph{shape} of the gate; the dark
energy's \emph{magnitude} is the horizon's own thermal scale, not energy
borrowed from structure. The single irreducible foundational input is
the microscopic volume-law postulate.

\section{Appendix G --- Supporting motivations and numerical probes for
the volume-law
postulate}\label{appendix-g-supporting-motivations-and-numerical-probes-for-the-volume-law-postulate}

This appendix collects the material demoted from §8 to keep the main
line focused: three further consistency \emph{motivations} for the
volume endpoint (§8.1) and the \emph{numerical probes} of the state and
counting halves of the postulate (§8.5). \textbf{None is a derivation};
each is a consistency check that makes the postulate less \emph{ad hoc}
without selecting it. (The CHR near-critical layer, the cross-horizon
black-hole reading, and the Verlinde dark-matter connection are
developed in a separate companion paper.)

\subsection{G.1 Three further consistency motivations for the volume
endpoint}\label{g.1-three-further-consistency-motivations-for-the-volume-endpoint}

\emph{(supplementing the geometric ceiling, the CKN scale/form split,
and the direct area-law disfavouring of §8.1)}

\begin{itemize}
\tightlist
\item
  \textbf{A driven steady state.} The horizon need not self-thermalise
  --- its scrambling time (∼ ln S / H ≈ 280/H) far exceeds its age.
  Instead the structure-growth gate (the rise of \(f_{\mathrm{sat}}\))
  \emph{drives} it off the ground state; with a scrambling-limited
  relaxation rate the volume-law state is \emph{maintained} as a
  non-equilibrium steady state throughout structure formation (the long
  scrambling time becomes an asset). The volume law is then a fixed
  point of SEDE's own dynamics, not an external input.
\item
  \textbf{A roughening universality class} (developed with the sibling
  \(SEDE_{\mathrm{V2}}\) collaboration). Modelling the horizon as a
  stochastic growing surface,
  ∂\(h/\partial t = \nu \nabla ^{2}h + (\lambda _{\mathrm{K}}/2)(\nabla h)^{2} + \eta\),
  the deformation is fixed by universality class: the thermal/linear
  Edwards--Wilkinson class \((\lambda _{\mathrm{K}} = 0)\) flows to the
  volume endpoint and corresponds to SEDE's purely-entropic F = 0
  condition, whereas a generic nonlinear (KPZ) surface gives a
  sub-volume value. This is the most speculative of the motivations and
  is included only as such.
\item
  \textbf{An independent emergent-gravity proposal} (Verlinde 2016). In
  Verlinde's programme the de Sitter vacuum carries a
  \textbf{volume-law} contribution to its entanglement entropy,
  \(S \propto  V_{\mathrm{H}}\), in addition to the horizon area law ---
  the same super-area scaling SEDE adopts at Δ = 1. A second,
  independently-motivated route arriving at a volume-law entropy makes
  the postulate less \emph{ad hoc}; it remains a motivation, not a
  proof, since Verlinde's volume-law derivation is itself heuristic.
\end{itemize}

\subsection{G.2 The state is maximally entangled (Majorana-SYK
diagonalisation)}\label{g.2-the-state-is-maximally-entangled-majorana-syk-diagonalisation}

The \emph{state} half of the postulate (§8.5) --- that the horizon dof
are maximally entangled (thermal), not in a low-entanglement ground
state --- is settled and not peculiar to SEDE. Exact diagonalisation of
the Majorana Sachdev--Ye--Kitaev model {[}Sachdev \& Ye 1993; Kitaev
2015{]}, the canonical chaotic, black-hole-dual system saturating the
Maldacena--Shenker--Stanford bound
\(\lambda _{\mathrm{L}} \leq  2\pi T\) {[}arXiv:1503.01409{]}
(\texttt{run\_syk\_scrambling.py}, V60), confirms: the gap ratio
reproduces Wigner--Dyson statistics with the correct N-mod-8 symmetry
class {[}García-García \& Verbaarschot 2016; Cotler et al 2017; Atas et
al 2013{]}, eigenstates saturate the Page value {[}Page 1993; Vidmar \&
Rigol 2017{]}, and the OTOC scrambles completely --- while an integrable
control does none of these. But this \textbf{does not fix Δ}: a black
hole is a \emph{maximal} scrambler yet \textbf{area-law} (S = A/4,
reproduced by string microstate counts), and SYK, being all-to-all
(geometry-free), cannot even pose the counting question.

\pandocbounded{\includegraphics[keepaspectratio]{../../output/fig_syk_scrambling}}
\textbf{Figure G1.} SYK as the horizon scrambler (q=4 chaotic vs q=2
integrable control). \emph{Left:} gap ratio climbs through the N-mod-8
RMT classes (GOE→GUE→GSE) --- the certificate of maximal chaos.
\emph{Centre:} eigenstate entanglement saturates the Page value.
\emph{Right:} the OTOC scrambles completely (C→2). This fixes the
\emph{state}, not the counting.

\subsection{G.3 The counting question --- numerical probes (allowed,
realisable, not
selected)}\label{g.3-the-counting-question-numerical-probes-allowed-realisable-not-selected}

The \emph{counting} half (§8.5) --- does the horizon dof number grow as
area \((N \propto  R^{d-2}\), Δ=0) or bulk volume
\((N \propto  R^{d-1}\), Δ=1)? --- is where SEDE's postulate lives.
Three numerical probes (\texttt{run\_entropy\_bounds.py},
\texttt{run\_causal\_set.py}, \texttt{run\_tensor\_network.py}; V62,
V61, V63):

\begin{longtable}[]{@{}
  >{\raggedright\arraybackslash}p{(\linewidth - 4\tabcolsep) * \real{0.3333}}
  >{\raggedright\arraybackslash}p{(\linewidth - 4\tabcolsep) * \real{0.3333}}
  >{\raggedright\arraybackslash}p{(\linewidth - 4\tabcolsep) * \real{0.3333}}@{}}
\toprule\noalign{}
\begin{minipage}[b]{\linewidth}\raggedright
question
\end{minipage} & \begin{minipage}[b]{\linewidth}\raggedright
probe
\end{minipage} & \begin{minipage}[b]{\linewidth}\raggedright
result
\end{minipage} \\
\midrule\noalign{}
\endhead
\bottomrule\noalign{}
\endlastfoot
\textbf{Allowed?} by the established entropy bounds & volume entropy vs
the smooth-area bounds & S ∝ R³ overshoots the
holographic/Bousso/Bekenstein bound A/4 by ∝ R (\textasciitilde10⁶¹ at
the cosmic horizon); consistent \textbf{only} as a genuine
\(d_{\mathrm{H}}=3\) fractal horizon (true area ∝ R³ ⟹
\(S = A_{\mathrm{fractal}}/4\) \emph{exactly} saturates the bound) \\
\textbf{Selected?} by quantum gravity & causal-set sprinkle, Minkowski
\textbf{and} de Sitter & both scalings present; the canonical horizon
object --- causal \emph{links} --- is \textbf{area} \((R^{d-2}\), the
Bekenstein--Hawking recovery), in flat space and, because dS is
conformally flat, identically in de Sitter. \textbf{Volume is not the
default.} (measured: bulk count 1.92≈2 ⟹ Δ=1; link count 0.98≈1 ⟹ Δ=0,
in 2+1D) \\
\textbf{Realisable?} in a holographic code & Ryu--Takayanagi min-cut,
local vs nonlocal & local (lattice) connectivity ⟹ area law (exponent
1.0); \textbf{nonlocal (expander) connectivity ⟹ volume law} (exponent
1.9). Δ=1 is realisable; its required ingredient is \emph{nonlocal}
horizon connectivity (the all-to-all structure of SYK). \\
\end{longtable}

The pattern: volume-counting is \textbf{allowed} (not by evading the
bounds but by the horizon being genuinely space-filling --- the
postulate itself), \textbf{realisable} (a nonlocal holographic code
produces \(S \propto  R^{d-1}\) exactly), but \textbf{not selected} (the
frameworks we can compute in default to area for the canonical horizon
object, and a black hole is area-law). No consistency argument available
to us \emph{derives} the volume count; equally, none \emph{excludes} it.
The two deciders are stated in §8.5: the dS static-patch Hilbert-space
dimension (theoretical) and the deformation Δ (empirical, §5.6/§6).

\end{document}
